\documentclass[lang=cn,10pt]{elegantbook}

\title{给孩子的语文书之何以中国}
\subtitle{LaTeX 经典之作}

\author{齐齐爸}
\institute{共同进步组织}
\date{June 9, 2026}
\version{0.1}
\bioinfo{自定义}{语文寄语}

\extrainfo{不要以为抹消过去，重新来过，即可发生什么改变。—— 比企谷八幡}

\setcounter{tocdepth}{3}

\logo{logo-blue.png}
\cover{cover.jpg}

% 本文档命令
\usepackage{array}
\newcommand{\ccr}[1]{\makecell{{\color{#1}\rule{1cm}{1cm}}}}

% 修改标题页的橙色带
% \definecolor{customcolor}{RGB}{32,178,170}
% \colorlet{coverlinecolor}{customcolor}

%新增
\usepackage{tabularx}            % 自动换列表格
\usepackage{booktabs}            % 三线表
\newcolumntype{L}{>{\raggedright\arraybackslash}X}  % 左对齐自动换列


\begin{document}

\maketitle
\frontmatter

\tableofcontents

\mainmatter

\chapter{前言}

何以中国时间线。

\begin{table}[htbp]
   \centering
   \caption{\textbf{《何以中国》纪录片各集信息汇总}}
   \begin{tabular}{llll}
    \toprule
    \multicolumn{1}{c}{\textbf{分集}} & \multicolumn{1}{c}{\textbf{时期}} & \multicolumn{1}{c}{\textbf{简介}} & \multicolumn{1}{c}{\textbf{相关遗址\&文化}} \\
     \midrule
     \multirow{5}{*}{秦汉} & \multirow{5}{*}{秦朝-汉朝} & \multirow{5}{5cm}{秦汉一统，四方攸同 \\ 发出何以中国之问} & 湖北云梦睡虎地秦墓 \\
     &&& 湖南里耶古城遗址 \\
     &&& 秦直道遗址（陕西-内蒙古） \\
     &&& 广东广州汉代南越国宫苑遗址 \\
     &&& 甘肃敦煌悬泉置遗址 \\
     &&& 江西南昌海昏侯墓 \\
     &&& 湖南长沙马王堆汉墓 \\
     \midrule
     \multirow{8}{*}{摇篮} & \multirow{8}{2cm}{万年前-5600年前} & \multirow{8}{5cm}{创生初始，文明摇篮 \\ 从考古现场还原先民的 \\ 饮食起居、耕作生产、 \\ 手工业器具等文明演进} & 江西仙人洞遗址 \\
     &&& 内蒙古敖汉旗兴隆洼遗址（玉） \\
     &&& 浙江余姚河姆渡遗址 \\
     &&& 湖南洪江高庙遗址（白陶器） \\
     &&& 浙江浦江上山遗址（稻） \\
     &&& 北京门头沟东胡林遗址（粟） \\
     &&& 陕西西安半坡遗址（仰韶文化-彩陶） \\
     \midrule
     \multirow{6}{*}{星斗} & \multirow{6}{2cm}{距今5500年前后} & \multirow{6}{5cm}{满天星斗，灿烂曙光 \\ 庙底沟彩陶、玉器在区域间传播} & 河南陕县庙底沟遗址（仰韶文化-花卉纹） \\
     &&& 河南灵宝西坡遗址（仰韶文化庙底沟类型） \\
     &&& 山东泰安大汶口遗址（海岱地区） \\
     &&& 安徽含山凌家滩遗址（玉） \\
     &&& 辽宁朝阳牛河梁遗址（红山文化-玉） \\
     &&& 浙江余杭良渚遗址（良渚文化-玉、水利） \\
     \midrule
     \multirow{3}{*}{古国} & \multirow{3}{2cm}{距今5000年前后} & \multirow{3}{5cm}{王国肇建，文明实证 \\ 区域性王权国家} & 湖北天门石家河遗址（长江中游） \\
     &&& 山东济南焦家遗址（黄河下游） \\
     &&& 山东龙山文化（蛋壳陶） \\
     \midrule
     \multirow{3}{*}{择中} & \multirow{3}{2cm}{距今4000年前后} & \multirow{3}{5cm}{九州共望，国之在中 \\ 一体化王朝文明} & 山西临汾陶寺遗址 \\
     &&& 陕西神木石峁遗址 \\
     &&& 河南偃师二里头（青铜） \\
     \midrule
     \multirow{2}{*}{殷商} & \multirow{2}{2cm}{距今3000年前后} & \multirow{2}{5cm}{商邑翼翼，四方之极} & 河南安阳殷墟（青铜、甲骨文） \\
     &&& 四川广汉三星堆遗址（青铜） \\
     \midrule
     \multirow{4}{*}{家国} & \multirow{4}{2cm}{距今2000年前后} & \multirow{4}{5cm}{矢其文德，洽此四方 \\ 宗法制+分封制} & 陕西宝鸡周原遗址 \\
     &&& 山西曲沃北赵晋侯墓地 \\
     &&& 湖北随州曾侯乙墓 \\
     \midrule
     \multirow{2}{*}{天下} & \multirow{2}{2cm}{东周-秦朝} & \multirow{2}{5cm}{天下远近，小大若一} & 甘肃张家川马家塬战国墓地 \\
     &&& 河北平山战国中山王墓 \\
     \bottomrule
   \end{tabular}%
   \label{tab:heyizhongguo}%
 \end{table}%








\chapter{第一集 \quad 秦汉}


\section{正文}


公元前223年，秦国和楚国的战事正陷入胶着，虽然楚王负刍已于一年前被秦军俘获，项燕和昌平军却仍在率军顽抗，即使骁勇如秦人，对上了楚人的坚韧，一时间也久攻难下。黑夫是秦军中的一名普通士兵，自参军以来，他和弟弟经常常并肩作战。然而这次两人却被分配到不同的战场，每一次分开，都可能是来不及告辞的生离死别。

安陆城原为楚国疆土，数十年前便已归秦国所有，如今这里居住着大量秦人，淮河边的战场距此不远，每一天都有家庭收到从军的亲人在战场上死去的消息，其余的家庭则在忧惧与思念中煎熬。因着两个弟弟都在从军，衷的一家也在等待中度日。这天终于有人带回了弟弟们的信，衷几乎是颤抖着展开，匆匆一眼便放下心来，赶紧向家人们转述信中内容。

黑夫和惊暂时无恙，前阵子他们分开作战，现在又见面了。黑夫和惊此时正在淮阳攻打反城已久，前路未卜，伤未可知。他们希望母亲寄来的钱一定不要少，还反复叮嘱衷，务必留意官府的军功授爵文书。

信末一行行罗列着对熟识亲友们的絮絮问候，最后一行是惊对妻女的关怀，他并嘱托妻子勉励照顾家中老人。身处生死瞬息的战场，这些琐碎的牵挂才是战士们心里最坚强的铠甲。安陆城外的连绵丘垄，多是城中居民们的死后安息之所。黑夫和惊曾经牵挂的亲人就葬在西郊一片，后来被称为“睡虎地”的坡地上。

20多个世纪后，考古工作者在其中的四号墓里发现了当年黑夫和惊寄出的信，信写在木牍上，木牍被珍重放置在木椁头箱中，收件人衷或许就是这座墓的墓主。

彼时秦国像黑夫和惊这样的平民，只有履险阻，冒矢石，以性命相系，方能换得军功之爵，以此改善个人与家庭的境况，或许他们最终还是生死暌违，家书便成为了生随死葬的珍藏。从辽东至陇西，自长城至南海，数百年间干戈不止，纷争不休。秦始皇陵兵马俑那每一张都迥异的面目，是曾披肩之锐的赳赳雄狮，也是曾被生离死别磋磨的血肉躯体。他们以及六国故地那更多未曾留名的战士，艰险岁月里为己即是为家，为国亦是为家。与家人间的牵挂，就是暗夜里明亮的光。

黑夫与惊写下家书的两年后，即公元前221年，世人期盼已久的和平局面崭然显露，横扫六合的秦军终能暂时放下手中的武器，看看那曾为之浴血的秀丽山川。

一座座曾经倾颓的城邑内外重新升起了连绵不断的炊烟，数百万平方公里的山河被纳入以中央政府为枢纽的郡县网络中，前所未有的辽阔王朝沐浴在初升的朝晖之下，巍巍伫立。

在连年的征战中，秦朝君臣早已认识到，东周列国之间在制度上无法兼容的差异，乃是攻扦不息的根源之一。度量衡、车轨、钱币、法律和文字，因为涉及到行政、手工业、商业和社会治理等重要领域，被迅速选定成为秦朝建立后需要统一的关键内容。

曾经的六国故地和秦都咸阳所在的关中，都曾经发现大量镶嵌秦诏版的铜权和铁权。权是称量重量的秤砣。权上的诏版写道，皇帝并天下后，百姓安定，令丞相隗状、王绾规范度量衡。凡不一致者，皆需统一。这道统一度量衡的诏书也被模印或铭刻在官定的度器和量器上，流布天下。

公元前219年，秦始皇在仲春之月开始了他的冬巡。此后八年间，他一度封禅，五次冬巡。他登峰山、攀泰山，上琅琊，经之罘，过东观，临碣石，至会稽，巡行之处皆刊石勒铭。

这七处石刻，由于风雨剥蚀，人事相侵，仅有琅琊、泰山二石尚存残玦，峰山刻石尚于摹刻本。刻石文辞以李斯新改定的小篆向关东宣示：乃今皇帝壹家天下，兵不复起，灾害灭除，黔守康定，利泽长久。秦要统治的是文字语言存在极大差异的四海八方，秦朝的文字改革使得汉字形体简化，部首和笔数固定，文字使用的标准得到统一。

湖南里耶古城遗址出土的这枚秦代木方，由六个残片拼接而成，以较为拙朴的篆书记录了其时文字的变更。“酉”如故更“酒”，将酒从酉分出，凡酒之意均用“酒”字，不再写作“酉”；“卿”如故更“鄉”，意思是公卿之“卿”仍然照旧，而记入鄉里之“鄉”，统一更用“鄉”字。

经过规范的文字书写和传播起来更有效率，使中央的政令能通过公文畅通地下达至各个郡县，也使更多人能够通过文字交流获得知识和文化的传承。书同文成为了中华文明赓续几千年而始终文化一体的奥秘。

渤海湾畔至今仍存多处秦代行宫遗址。在辽宁绥中万家镇南边，就发现了六组相互关联的秦汉建筑遗址，其中三组皆在海岸，朝向海中的“姜女石”呈合抱之势。这或许就是当年秦始皇、汉武帝曾登临过的“碣石”。石碑地建筑群可能就是秦汉时的“碣石宫”。

曲尺形宫墙之内，建筑高低错落，疏密相间，依稀可见昔时廊道鳗回，院落镶嵌的格局。洪波高涌，碣石为阙，此为秦之国门。临碣石而东望，日升月落，天旋海沸，终至陆地之极，身后是万里河山。

这是秦始皇东巡时，臣属们镌刻在琅琊台石碑上的文辞，残石饱经沧桑，铭文也大多漫浊不清，但内容依然震撼人心。

六合之内，皇帝之土，西涉流沙，南尽北户，东有东海，北过大夏，人迹所至，无不臣者。但天下并不太平，北方边境地平线上争端频仍，匈奴的铁蹄声一旦响起，劫掠与杀戮就会到来。

30万秦军正等待主帅号令，蒙恬已记不清自己经历了多少场战役，只知道在争夺天下时，他是秦王的长策，四海平定时，他就是皇帝的敲扑。面对匈奴滋扰，蒙恬率军却敌，把秦朝的边境线向北推进700余里，直抵阴山脚下。为了在外敌再度来犯时能够迅速地预警并组织反击，蒙恬受命将原先秦、赵、燕三国的长城与由洮水而来的长城相连，沿途起城邑修亭障，随后领兵驻守上郡十余年。

秦始皇35年（公元前212年），蒙恬率军修建了直道。这条长约743公里、宽达50米的道路，以几近直线的走向连接了关中与北方草原地区。经行于山地和丘陵时，削平山脊，填平山谷；经行于沙地和草原时，以夯打或掺石的方式加固路面；经行山侧河畔时，设置护坡和排水沟。这不仅是皇帝巡行北方的“高速公路”，更是支援长城军事防务的运输线。

长城如盾，直道就是持盾的手臂，它们共同铸就了延袤万余里的坚固藩篱。百余年后，史官司马迁行经直道时，忍不住为这项暂山埋谷的工程惊叹。然而，他也指出长城和直道费人力财力无数，虽与王朝的防守有利，却施之过急，未留给刚刚从战乱中安定下来的众生以喘息之机。尽管赢政在“始皇帝”的名号里寄寓着他千秋万世的构想，却仅二世而亡。

直到汉才带来了较长时间的大部分疆域的和平。作为汉朝的开启者，刘邦最初选择的都城是洛阳，他倚中的大臣都是关东人，无不希望永久地定都于此，只有齐人娄静和留侯张良力主迁回关中。虽然战事以息，局势却仍不稳定。关中坐拥天险，不易受敌人威胁，且土地膏腴，交通便利，未来可容纳更多的人口。此外，还可以割断关东功臣集团们与乡土间的利益关系。公元前202年，刘邦带着他的朝廷西迁后暂居于栎阳，这里是秦人旧都，商鞅在此变法，秦的壮大也正是从这里加速。

两年后，在50多公里以外的渭河南岸，由秦朝渭南离宫之一兴乐宫改建而成的长乐宫竣工，意味着新王朝都城的诞生，长治久安的愿望被寄托在这座新都长安的名称里。

未央宫紧闭的宫门一年后轰然洞开。刘邦为这座新宫室的华丽所震惊，即便他见过秦朝巍峨的咸阳宫，其宏伟也远不及眼前的未央宫。他想起了秦王的教训，想起了刚刚剿灭的叛乱，作势怒斥主持工程的萧何，萧何俯首答道“夫天子四海为家，非壮丽无以重威，且无令后世有以加也。”刘邦已领略了威家四海的滋味，这话令他大悦，但清醒之后仍心有惕惕。此后至死也将主要的政治活动放在长乐宫中。

西汉都城长安是当时世界上规模最大的城市之一。龙首原上，长乐、未央两宫东西并峙，于高处制掣全城。桂宫、北宫、明光宫横列两宫之北。长安有东西二市，八街九陌，12城门160闾里。

渭水之北，帝陵有九，五陵邑拱卫京师，东郊多丘墓，东南郊有文、宣二帝之陵，南郊为礼制建筑，西郊则是离宫苑囿。

200年间多少政治风云由此升起，多少人事兴亡在此兴起，个个后席卷着汉朝疆域的每个角落。使臣陆贾从长安出发，携带着汉朝预备封授给南越王的印绶前往番禺。

南越王赵佗本是秦朝派往岭南的一名地方官员，秦末趁乱据五岭以南自立为王。刘邦平定岭北后，不想再冒着瘴疫的危险翻山越岭，尤其越地丛林深密，水道纵横，并不是汉军熟悉的作战环境。他希望用一种和平的方式，令赵佗主动归顺。曾随高祖征战的陆贾，对一路的旧关防再熟悉不过。如今关防已撤，商贾周流，天下繁华气象初现。

过了五岭便是南越，陆贾一行弃车马，换舟楫顺江而下，鹧鸪声里，两岸的木棉树穿古榕而出，绿水明波荡漾。

番禺城中央的宫苑，自然不及汉家都城宏丽，却处处可见效仿之意。赵佗免冠露髻在大殿之上交脚而坐，倨傲地望向汉家使臣。这在礼仪严明的中原人看来，他分明以走成人自居，摆出了割裂的姿态。陆贾款款陈词，既点出赵佗本为中原人的身份，又在对时事的条分缕剖中暗藏了威慑与怀柔。

听毕，南越王肃然起身，重新庄重地按中原礼仪跪坐相待。赵佗先征询得到陆贾认为他比汉家丞相们更贤能的答案后，又在追问中将比诏的对象换成汉家皇帝。陆贾看出南越王心思，并未直接评价二人才干，却转而分析汉越之间的实力差异。身在岭南，无法与中原王朝抗衡，赵佗如何不明白其中因由，如今借陆贾之口说明白，倒也让多年心事了结。虽然股肱大臣们多是中原人，但赵佗的心里是寂寞的，无人能了解他欲振翅于寰宇的鸿鹄之心，越中无人可共语。陆贾这一介使臣的才干胆识，却令他引以为知交。赵佗已经可以想见汉家朝廷的英才济济。

此后赵佗常常邀陆贾饮酒谈天，有时候宴会设在华音宫畔的曲流石渠旁，有时在波光粼粼的蕃池旁。宫苑中，梅花鹿撷食着灌木的嫩芽，龟鳖从渠畔的斜坡爬上岸享受阳光，宫人兢兢业业地清点每颗从北边移植来的枣树上结了多少颗果实。如此数月，赵佗在听取陆贾对天下大势的剖析后，表示从此臣服汉室，并馈赠陆贾宝物，送其北归。

昔日的南越宫苑虽毁，周遭的繁华却2000年未息，其遗址便沉睡在广州的闹市之下。曾经流水蜿蜒的石渠；曾经天光云影的蕃池，至今还回放着历史每一次的柳暗花明，峰回路转。

一件陶提简盖残片上的戳印暗示了赵佗的心境——华音宫。汉初的君臣常常讨论，为何秦统一了疆域开创了前所未有的辽阔版图，却短祚覆亡？在反思中他们逐渐达成共识，柔化统一之策略，厚积统一之基础，化育民意于无形，静水自可流深。

在数代君臣的治理下因长久战乱而凋敝的民生逐渐恢复，同时恢复的还有那曾因为秦而失去的对于“大一统”的信心。武帝接过的是一个初具气象的王朝，此时董仲舒挟数世儒生积累的智慧而来，提出依托儒家思想构建大一统下的国家认同。这显然触动了汉武帝的所思与所愿。他看到了让六合同风、四方共俗的可能性；看到了不再依靠血缘关系治理社会的可能性。南昌市新建区的郭墩山下，沉睡着刘贺夫妇及其子孙，他的祖母是生前宠冠后宫，死后与武帝合葬的李夫人。刘贺继承了父亲的昌邑王位，昭帝崩逝后，迅速被霍光等定为新帝人选。然而在帝位上仅仅待了27天，刘贺便成为关中群臣攻讦的目标，背负无数罪名被褫夺王位，只得以列侯之身终老海昏。

海昏侯墓内建有400平方米的木构椁室，主墓室的设计模仿刘贺生前的真实生活场景，事死如事生。墓中随葬的大量黄金和精美器物可以窥见当时国家的强盛富足。在主椁室的西室内，摆放着莲枝灯、博山炉和漆案。这面漆木衣镜屏陈列在床榻旁最显眼的位置，镜屏是刘贺生前所用陪伴他走过人生大半，镜背表面绘着孔子和他的五个弟子，这是目前所见年代最早的孔子图像。人像两侧以墨书抄录着人物生平及言行：被任用时就施展抱负，不被任用就藏身自好，只有我和你才能这样吧。

当年汉宣帝将刘贺封至南昌做海昏侯，实质上是要他离开富裕的山东到偏远的“南藩”豫章就国还令他不得回长安祭祀祖先。

关山难越，谁悲失路之人，萍水相逢尽是他乡之客。情难自己时，海昏侯乘流东望，辄愤慨而还。

刘贺墓内随葬大量文书，除签牌、奏牍等外，以规整抄录的儒家经典居多，唯有一版摘抄了《论语》断篇的木犊，字体率性随意与别不同。在废黜之后被时时监视的日子里，刘贺无以纾解，唯有寄情于经典之间。他亲笔抄写下了这样的字句“子曰：吾者知乎哉？毋知，有鄙夫问乎吾，空空如也。叩其两端而竭。”人生多疑惑而少智慧，只得穷极叩问，竭尽求索，或许刘贺就在圣贤的教诲间寻找着安身立命的心境。

武帝推崇儒家像刘贺这样的王侯往往能以大儒、博士为师，但对于普天下的学子来说，负笈求学于都城，才是他们的夙愿和殊荣。自武帝立太学后办学规模日渐攀升，到质帝时，从各地前往洛阳求学的学子已达到3万人之多。太学中黄舍上千，鳞次栉比。灵帝召集当时鸿儒正订六经，刻石立于太学。在没有印刷术的年代这“熹平石经”就是学子们的教科书，前来观视与临摹石经的车骑络绎不绝，甚至一度堵塞了太学门口的街道，秦完成了形式上的统一。

汉则使统一的理念嵌入到思想和文化的基因里，以秦朝的文字统一为根系，两汉以太学为参天巨木之干，将治国的理念与秩序寄寓在教育之中，快速抽生出荣茂的枝叶，伸向万里河山的各个角落。即使在当时的偏远之地，都曾有同样熟悉的诵读声，普天下的莘莘学子，从高门绣户的列侯到边陲关隘的无名小卒，西域精绝城里的当地人也都曾捧着同样的《仓颉篇》识文习字，念着同样的冬寒夏暑，玄气阴阳。汉兼天下，海内并廁。他们都将成长为治理这片辽阔疆土的力量。

在注重文治的同时，国家安全仍然是汉朝最为关注的议题。武帝在位时国内人口充盈、物资丰饶，诸侯国也再不能掣肘，海内为一，兵马整肃，但匈奴在这期间也发展壮大，对北境的安全造成了威胁。

公元前138年，汉武帝公开招募西行联络月氏、大渊和乌孙以抗击匈奴的使者。27岁的汉中青年张骞毅然走上“凿空”西域之途，这一去就是整整13年。

从长安出发时，车马粼粼百余人，回程路上只有两人，13年间有利益相诱，有武力相胁，温柔乡、囹圄境，却从未丢失身为汉使的旌节，从未磨灭此行的初心。

公元前130年，年轻的车骑将军卫青直捣龙城，获得了汉朝对战匈奴的初次胜利。公元前123年，嫖姚校尉霍去病率轻骑直入大漠。两年后，又在祁连山奇袭匈奴的浑邪王、休屠王。

公元前119年卫青、霍去病远征漠北，令单于遁逃，封狼居胥。八年后，汉王朝先后设酒泉、张掖、敦煌、武威四郡，牢牢掌控河西走廊。此后匈奴虽然还有零星骚扰却已元气大伤，不再成为气候。尽管如此，庞大的戍边队伍仍在长城沿线屯驻，兵马不懈卫护着汉朝北境的绝对安全。

转也是其中的一员。秋风渐凉，这日正值休假，同袍们有的凑在一起玩六博戏，有的去与随军的家人团聚，有的约了朋友聚会饮酒……转听闻此前受托去家中探视的友人幼卿，此时已来到相距不远的肩水都仓，很想赶去见上一面，却因公事无法脱身，只得修书一封托人转交。他请幼卿往家中捎话“老人平安便是离乡远行者之幸。天气凉了，千万要记得及时添衣，善进饮食。”正要搁笔时想起秋天是匈奴人经常进犯的季节，赶紧又添上几句嘱咐“幼卿你独自行走于关外，此身只一并无二，千万当心。”

转和幼卿如此平凡，他们漫长的人生我们只能窥见这只字片语的时光，再不知后事如何。转写给幼卿的那封信也一裂为二，散落坞墙之下，长埋黄沙之中。

1973年的夏天，大西北的戈壁上仍旧延续着年复一年如期而至的高温和风沙。昔年的居延曾有“弱水流沙”之称，此时“弱水”已更名为额济纳河，流沙依旧无休止。

从7月到9月，甘肃居延考古队都在肩水发掘，与秦时无两的明月照在汉时的关址上，连流沙和相思也并无两样。发掘常常被风沙中断，停了之后又再重来，如此历经将近三个月，汉代河西的这一座金关终于再度重现。

37个探方中11577枚汉简记录着2000多年前的文书与人事往来。考古工作者耐心地将简牍一枚枚清洗、分类、缀合、释读。在这个过程中，30号探方和26号探方中的两片断简被合而为一。碎了的信重新还原，我们才成为时隔二十个世纪的读信人。

为汉朝卫护北部边境国家安全的不仅有使者、将士也有和亲的公主们。乌孙是西域诸部族中较为强大的一个，随畜逐水草而居。张骞向武帝献策，以财物相诱送公主和亲，说服乌孙与大宛、月氏形成夹击匈奴之势。

解忧公主是楚王刘戊的孙女，她的爷爷在景帝时参与了一场史称“七国之乱”的反叛行动，最终兵败自杀。或许因为皇室宗亲和罪臣后嗣的双重身份，她成为汉朝派往乌孙的第二位公主。此前不久，江都王女刘细君刚刚死在异乡，细君公主生前所作的悲歌早已传遍长安“居常土思兮心内伤，愿为黄鹄兮归故乡。”

与草长莺飞的故乡迥异的风物与习俗，是武帝诏书中“欲与乌孙共灭胡”的重托也无法冲淡的忧愁。在政治漩涡中成长的解忧有着与细君公主不同的坚韧心性。她带着王朝的托付，一意西行陪她同往的还有聪慧的侍者冯燎。

从长安向西他们第一次看见天地间如此空荡的景象，四面八方的风无所阻滞地在莽莽平野间来往。第一次看见尘世间如此壮阔的奇观，雪山、沙漠、戈壁、龙堆以及镶嵌其间的湖泊与绿洲。

到乌孙后，解忧嫁与国王军须靡，冯燎则为右大将之妻。他们居毡帐，饮酪浆，白天听着群马嘶鸣，夜晚枕着猎猎风声。对故乡的思念此刻被化作对故土的使命，开始了他们在西域长达数十年的斡旋。

君须靡死后，解忧遵照乌孙的习俗嫁给了他的堂弟翁归靡。公主在内润物无声的影响着乌孙的政局。她的新丈夫不再游移于汉匈之间，而选择了一心归附于汉。冯夫人在外，持汉节为公主行赏赐，邀迎西域诸部族之心。

解忧在乌孙40余年的经营眼看即将换来长久的安宁，翁归靡却猝然薨逝。

60岁的解忧不得不嫁给了她的第三任丈夫泥靡。这是她第一任丈夫与匈奴人的儿子，对汉朝并无好感，她处境艰难，一度生死悬于一线。事情的转机发生在泥靡被杀之后，在冯燎的帮助下，在汉庭的支持下，他的儿子元贵靡最终成为了乌孙的大昆弥。一度屯积于敦煌待命的汉军得以不战而还。

这一年是甘露二年（公元前52年）解忧已经年近七旬，她尚记得自己来时明眸皓齿，风华正茂，如今遥望家乡时，她的眼神却越来越模糊，她怕再不回去便忘记故土的模样。她一生里大部分的时光都在这伊犁河畔度过，为了故国所赋予的使命，变了语言，改了习惯，默默地适应着这里的一切。

她的三任丈夫，她的儿女都在这里，但是她不能忘记自己的来处，那灼灼的桃花，那苒苒的桑梓，那方正的文字，那绮丽的丝绸，这也是千百年后仍然嵌在中国文化里的基因——对故乡的眷念始终深深镌刻在骨血里。

从玉门到敦煌一路皆是茫茫戈壁，汉代的马甲骑联翩，也曾由此经行。在距敦煌市区还有60公里的柳格高速南边，悬泉置伫立在连绵而荒凉的山脉之下，这是敦煌郡效谷县属下的一处邮驿机构。日常在此工作的吏卒有30余名，负责传递官方文书、军情急报，接待往来的各级官员和各方使者。

啬夫弘是悬泉置存在的前后数百年时间里任职最长的一个。从宣帝元康三年到元帝初元四年的十八年间，他几乎未曾离开过这里。他为都护、使者、公主、将军、列侯们传过无数封信，准备过无数次行程所需的马匹与粮草。他仍清楚记得12年前他刚到悬泉置不久，长罗侯常惠护送少主相夫来到敦煌，不久后又因和亲不成而折返，之后便听说了解忧公主和冯夫人那惊心动魄的经历。

公元前58年，他看见丞相史李尊护送复原的戍卒返回原籍，在归乡的队伍中总有招车上载着的一具具棺木，那是不幸死在异乡的戍卒们。曾经，新征的年轻人从悬泉置经过，去往敦煌等地服役。如今他们也经这里返回各自的家乡。河东、南阳、颍川、东郡、魏郡、淮阳国，目送那些曾卫护过一方安宁的人们离开。啬夫弘仔细录下公文中戍卒们的原籍。

甘露二年（公元前52年）二月十二日傍晚，平望驿骑带来了长罗侯和解忧公主的书信，他们的书信一向与西域局势有关。啬夫弘未敢怠慢，将之交托给驿骑朱定，让他连夜赶往万年驿。

一年后啬夫弘收到了下行的公文，让他为从乌孙归汉的解忧公主一行准备车驾，他这才知道曾送出的那封信，正是公主归家的恳求。那年冬天他看见年老的解忧公主和她的三个儿孙匆忙上前迎接。

从赤古城到悬泉置是1273公里，从悬泉置到长安城是1372公里，这位离家50年的大汉公主已在遥远的西域看过了将近2万轮的日升月落，古稀之年的她终于踏上了归途。岁月轮转，烽火已熄，古道上音尘渺绝，戍卒的书信埋在黄沙下，王侯的宫苑隐在高楼间，而那些城邑、道路、边关、堰渠曾见证着秦汉时期每一个鲜活的人，如何在这片山河湖海间努力经营着自己的一生，慢慢垒就矗立历史长河中的巍巍王朝。

他们曾眼望星辰；他们曾背负长天；他们曾跋山涉水；他们曾生死离别；他们曾在不知觉间成为两千年后今日中国的奠基。2000多年前长安城中的汉武帝望向九州华夏的舆图，他决意为汉王朝打开连接外部世界的通衢，沿着士兵、公主、使臣们曾经往来的河西走廊，越过葱岭向西而行；沿着商贾、船队曾经航行的线路，穿过大海向西而行。海陆并行的通途寄寓着强大的国家对外文化经贸交流的渴望，也孕育着中国大地此后对外交通的格局。

位于山东临沂的北寨汉画像石墓，刻画着一幅充满日常生活气息的乐舞百戏图。跳丸、飞剑、走索、寻橦……文献中记载的民间杂技鲜活的展示在人们面前，更特别的当属经丝绸之路传来的“鱼龙曼衍”之戏。据说这种戏使用西域幻术能够呈现出鱼龙变化的神异场景，不知不觉间外来文化沿着河西走廊流入生机勃勃、兼容并蓄的东方。走入寻常百姓家，而曾经横亘在东西方之间的山海重重已成通途。

统一而不断强大的中国，从秦汉时代启程，通过强有力的中央政府建构起多元一体的文明持续发展至今。

独特的中华文明如何形成？秦汉之前是否有更早的中国雏形？三皇五帝的古史系统是否可靠？中国文明5000年是否可信？让我们追随考古学家的脚步，作万年的时光旅程。

看我们的文明在摇篮里的孩提模样，被她稚嫩的笑容打动，看她在辽阔的山川大地上蓬勃成长，在5000多年前形成一个可以称作中国的共同体，长成为真正的文明；看她经历龙山时代500年激荡整合，将共同体凝聚为政体，为夏商王朝的诞生欢欣鼓舞；看她究天人之际，由尊帝到尚德，由编户至齐民，完成政体、社会乃至华夏文明的整合。

我们自秦汉溯源，探寻何以中国，贯通五千年文明血脉。我们经秦汉传续，虽历五千年而不老，中国依旧如旭日朝阳。


\section{典故}

\begin{table}[htbp]
\centering
\small
\renewcommand{\arraystretch}{1.2}
\caption{\textbf{第一集“秦汉”典故与文化知识}}
\begin{tabularx}{\textwidth}{@{}c p{3cm} L@{}}
\toprule
%每列的表头内容居中
\multicolumn{1}{c}{\textbf{序号}} & \multicolumn{1}{c}{\textbf{典故/文化名词}} & \multicolumn{1}{c}{\textbf{文中相关语句（或出处）}} \\
\midrule
1 & 黑夫与惊 & “黑夫是秦军中的一名普通士兵……黑夫和惊暂时无恙” \\
2 & 睡虎地秦墓 & “后来被称为‘睡虎地’的坡地上……四号墓里发现了当年黑夫和惊寄出的信” \\
3 & 秦始皇封禅（shàn） & “他一度封禅，五次冬巡” \\
4 & 李斯小篆（zhuàn） & “刻石文辞以李斯新改定的小篆” \\
5 & 书同文 & “书同文成为了中华文明赓（gēng）续几千年而始终文化一体的奥秘” \\
6 & 碣（jié）石宫 & “石碑地建筑群可能就是秦汉时的‘碣石宫’” \\
7 & 蒙恬（tián）修长城与直道 & “蒙恬受命将原先秦、赵、燕三国的长城相连”“蒙恬率军修建了直道” \\
8 & 司马迁叹直道 & “百余年后，史官司马迁行经直道时，忍不住为这项堑（qiàn）山堙（yīn）谷的工程惊叹” \\
9 & 萧何建未央宫 & “夫天子四海为家，非壮丽无以重威” \\
10 & 陆贾出使南越 & “陆贾款款陈词……赵佗肃然起身，重新庄重地按中原礼仪跪坐相待” \\
11 & 赵佗（tuó）与南越国 & “南越王赵佗本是秦朝派往岭南的一名地方官员，秦末趁乱据五岭以南自立为王” \\
12 & 华音宫 & “一件陶提筒盖残片上的戳（chuō）印暗示了赵佗的心境——华音宫” \\
13 & 董仲舒 & “董仲舒挟（xié）数世儒生积累的智慧而来，提出依托儒家思想构建大一统下的国家认同” \\
14 & 海昏侯刘贺 & “南昌市新建区的墎（guō）墩山下，沉睡着刘贺夫妇及其子孙” \\
15 & 关山难越，萍水相逢 & “关山难越，谁悲失路之人？萍水相逢，尽是他乡之客” \\
16 & 熹（xī）平石经 & “灵帝召集当时鸿儒正订六经，刻石立于太学” \\
17 & 太学 & “自武帝立太学后办学规模日渐攀升……到质帝时已达3万人” \\
18 & 张骞（qiān）凿（záo）空 & “27岁的汉中青年张骞毅然走上‘凿空’西域之途” \\
19 & 卫青、霍去病 & “卫青直捣（dǎo）龙城……霍去病封狼居胥（xū）” \\
20 & 封狼居胥（xū） & “令单于遁逃，封狼居胥” \\
21 & 解忧公主 & “解忧公主是楚王刘戊（wù）的孙女……一意西行” \\
22 & 冯嫽（liáo） & “聪慧的侍者冯嫽” \\
23 & 悬泉置 & “悬泉置伫（zhù）立在连绵而荒凉的山脉之下……啬（sè）夫弘是任职最长的一个” \\
24 & 居延汉简 & “37个探方中11577枚汉简记录着2000多年前的文书与人事往来” \\
25 & 鱼龙曼衍（yǎn） & “经丝绸之路传来的‘鱼龙曼衍’之戏” \\
26 & 龙山时代 & “经历龙山时代500年激荡整合” \\
27 & 究天人之际 & “看她究天人之际，由尊帝到尚德，由编户至齐民” \\
\bottomrule
\end{tabularx}
\label{tab:heyizhongguo_fixed}
\end{table}


\section{经典词语}

\begin{table}[htbp]
\centering
\small
\renewcommand{\arraystretch}{1.2}
\caption{\textbf{第一集“秦汉”精彩成语与词语（注音+解释+例句）}}
\begin{tabularx}{\textwidth}{@{}c p{2.6cm} p{2.6cm} X@{}}
\toprule
\multicolumn{1}{c}{\textbf{序号}} & \multicolumn{1}{c}{\textbf{成语/词语}} & \multicolumn{1}{c}{\textbf{名词解释}} & \multicolumn{1}{c}{\textbf{示例句}} \\
\midrule
1 & 胶着（jiāo zhuó） & 比喻相持不下或工作不能进展，如胶黏着。 & “敌人西侵的部队在长江的南北两岸都呈着胶着状态了。”——郭沫若《羽书集·致华南友人们》 \\
2 & 生离死别（shēng lí sǐ bié） & 很难再见的离别或永久的离别。 & “独有王夫人和宝钗娘儿两个倒像生离死别的一般，那眼泪也不知从那里来的，直流下来，几乎失声哭出。”——曹雪芹《红楼梦》第一百十九回 \\
3 & 忧心忡忡（yōu xīn chōng chōng） & 形容忧虑不安的样子。 & “忧心忡忡，待旦而入。”——韩愈《送李愿归盘谷序》 \\
4 & 絮絮叨叨（xù xù dāo dāo） & 形容说话啰嗦，连续不断。 & “那婆子絮絮叨叨，只管说些闲话。”——吴敬梓《儒林外史》第二十回 \\
5 & 披坚执锐（pī jiān zhí ruì） & 穿着坚固的铠甲，拿着锋利的武器。指奔赴战场作战。 & “披坚执锐，临难不顾，身先士卒。”——陈寿《三国志·魏书·曹彰传》 \\
6 & 生死暌违（shēng sǐ kuí wéi） & 生与死分离，不得相见。暌违：分离，隔离。 & “新知虽已乐，旧爱尽暌违。”——何逊《赠诸游旧》 \\
7 & 干戈不息（gān gē bù xī） & 战争不停息。干戈：古代兵器，借指战争。 & “干戈不息，征徭未休。”——白居易《与元九书》 \\
8 & 赳赳武夫（jiū jiū wǔ fū） & 威武雄壮的武士。 & “赳赳武夫，公侯干城。”——《诗经·周南·兔罝》 \\
9 & 李代桃僵（lǐ dài táo jiāng） & 以此代彼，代人受过。原比喻兄弟相爱相助。 & “桃生露井上，李树生桃旁。虫来啮桃根，李树代桃僵。”——《乐府诗集·鸡鸣篇》 \\
10 & 鳞次栉比（lín cì zhì bǐ） & 像鱼鳞和梳子齿那样有次序地排列着。多用来形容房屋或船只等排列得很密很整齐。 & “市井街坊鳞次比，鲛绡泉贝积如山。”——北宋·李邴《咏宋代泉州海外交通》 \\
11 & 势如破竹（shì rú pò zhú） & 比喻节节胜利，毫无阻碍。 & “今兵威已振，譬如破竹，数节之后，皆迎刃而解。”——房玄龄等《晋书·杜预传》 \\
12 & 背水一战（bèi shuǐ yī zhàn） & 比喻决一死战。 & “信乃使万人先行，出，背水陈。赵军望见而大笑。”——司马迁《史记·淮阴侯列传》 \\
13 & 直捣黄龙（zhí dǎo huáng lóng） & 比喻直捣敌人巢穴。 & “直捣黄龙府，与诸君痛饮耳。”——《宋史·岳飞传》 \\
14 & 封狼居胥（fēng láng jū xū） & 比喻建立显赫功勋。 & “骠骑将军去病率师躬将所获荤允之士，登临瀚海，封狼居胥山。”——班固《汉书·霍去病传》 \\
15 & 筚路蓝缕（bì lù lán lǚ） & 形容创业艰苦。 & “筚路蓝缕，以启山林。”——左丘明《左传·宣公十二年》 \\
16 & 家喻户晓（jiā yù hù xiǎo） & 家家户户都知道。形容人所共知。 & “使其家喻户晓，则人人自危。”——欧阳修《乞罢刽子配隶人户札子》 \\
17 & 前所未有（qián suǒ wèi yǒu） & 历史上从来没有过。 & “而中国政体，实为前所未有之奇变。”——梁启超《新民说》 \\
18 & 浴血奋战（yù xuè fèn zhàn） & 形容顽强地拼死战斗。 & 现代常用成语，多见于抗战文学。例：“八路军在阵地上浴血奋战，终于守住了防线。” \\
19 & 惊心动魄（jīng xīn dòng pò） & 形容使人感受很深，震动很大。 & “其文则惊心动魄，一字千金。”——钟嵘《诗品》卷上 \\
20 & 欣欣向荣（xīn xīn xiàng róng） & 比喻事业蓬勃发展，兴旺昌盛。 & “木欣欣以向荣，泉涓涓而始流。”——陶渊明《归去来兮辞》 \\
\bottomrule
\end{tabularx}
\label{tab:excellent_words}
\end{table}


\chapter{第二集 \quad 摇篮}

\section{正文}

天地玄黄，宇宙洪荒，日月盈仄，辰宿列张。我们的中国正从远古混沌中渐渐清晰。6,500万年前青藏高原隆起，黄土高原、云贵高原、内蒙古高原，抬升形成了中国三级阶梯式的基本地理格局。自此，大江东去，日月西沉，疾风地北，烟雨江南，是我们所熟悉的中国。距今约 200万年前，直立人出现在东亚大陆东方古人类繁衍生息持续演化的故事拉开帷幕。

云南元谋人、陕西蓝田人、北京人、广东马坝人、山顶洞人，我们只能依据发现地点给这些远古的祖先命名。他们离开我们太过久远，久远的骨骼已经石化，还执拗地保留着中国古人类独有的特征。头骨正中的矢状脊，突出的面部，高而前凸的颧骨、阔鼻、铲形上门齿、下颌的圆枕．.

北京房山周口店遗址的20多个地点都发现了古人类生活过的遗迹。在位于第27 地点的田园洞地层里，出土了距今约4万年的人类遗骸化石。其中包括下颌骨、脊椎骨、股骨跟骨等体骨的大部分骨骼，属于同一位中年男性个体。

古人类遗骸或遗迹中残存着极其微量的 DNA片段，让我们可以直接研究过去人群的遗传信息。正是从田园洞人的腿骨上，古遗传学家成功捕获到测序较完整的基因组序列，与现代人基因特征完全符合。这是中国第一例人类古基因组，也是目前为止获得的东亚最古老的人类基因组。遗传意义上的亚洲人的祖先，在距今4万年前已经出现，他们与世界各地现代人的出现基本同时。

然而，挑战在3万年前后降临，严酷的冰期席卷全球，随着气温不断下降，巨大的冰盖覆盖了北美、欧洲和亚洲北部的大部分地区。

曾经水草丰美，猛犸成群的草场突然变成了一望无垠的冰原和荒漠；曾经湿热瘟瘴的丛林，却可能变成了气候宜人的河谷。适宜人类生存的区域不断偏移和缩小，风雪交加的世界里，田园洞人的身影渐渐模糊．...

约两万年前，江西万年县大源盆地的山间，一群人正在黄昏的光线里寻找着入夜后的栖息之所。眼尖的少年看见小河山的脚下，因为河床的缩小和后退，新露出了一处山洞。山洞内外地势平常，离水源不远也不近，是个十分理想的住所。于是人们对洞内的地面稍作平整后，生起一堆火，在这里度过了大雨滂沱的寒冷夜晚，这是被称为末次冰期最盛期的时代。环境发生着急剧的变化，熟悉的食物逐渐消失，现代人面临着自出现以来最严峻的生存挑战。

初升的日光冷清而凛冽，仙人洞族人已早早开始了一天的忙碌。老人敲击燧石砸出石片，锋利的刃部足以割开兽皮；女人用骨针缝补兽皮，制成过冬之衣；族人们收获野生禾草，用骨鱼杈捕鱼，也捕捞河蚌和田螺；孩子们则在原野和林间操练着捕猎的技能。

河岸边的少年就地取用夹杂着石英砂砾的泥土，摻水做成泥条一圈圈盘筑成園底的深腹罐。他在罐内外不停的拍打，使得泥条间的粘合更加规整、紧密。天色越来越明亮，少年将做好的罐子晾在一边，再将另一些早已阴干的陶坯放进露天的火堆中，渐渐的泥土变得干燥而坚硬。

陶器是数百万年的演化以来，人类创造出的第一种人工材料。这是目前已知世界上最早的陶容器﹣﹣仙人洞出土的陶器残片。外表面仍留存有烟食与火烧的痕迹，这显示他们应该是炊具。从此在生食和烧烤之外，人类还可以通过粒食和蒸煮的方式加工食物，从富含淀粉的食材和肉类中获取更多的能量。

外出的人们带着一天的收获归来，山洞里再度热闹起来，有人用锋利的刮削器肢解了一头斑鹿，放在火堆上炙烤。孩子们围陶罐而坐，煮熟的食物香气扑鼻，有人敲开鹿骨将骨髓和另一些软烂的食物留给老人。

天气越来越冷，食物的获取并不稳定，不是每天都能像这样饱餐。因此所有的食物都要物尽其用。在蒸腾的热气中，山洞里的火光明明灭灭仿佛悬挂在远古文明之路上的一盏夜灯，远在数百万年前当古猿人开始直立行走那一刻，他们看到了更远的世界，却也暴露出柔软的腹部。此后，漫长的狩猎采集生涯里，更需要彼此守护。

广东英德青塘 13,500年前的拂晓，一位豆蔻年华的少女刚刚离世，同伴们将她抬出日常活动的洞穴，穿过一片疏林来到另一个狭长幽深的洞中。他们选了一处凹坑，将少女摆出蹲踞的姿势放进去，在她背后垫了六块石灰岩角砾，把她生前常常使用的一枚骨针也放了进去。

这一次仍然生存的人们心头涌起了悲哀和思念，他们不能再像其他的动物一样任由死去的同伴在自然中慢慢分解，他们想用这样一种"仪式"来告别曾经一同生活过的人，让她用一种特别的方式沉睡，就像婴儿回到了母体，就像活着的人和死去的人之间永远有一种联结。

一万年前随着全新世的到来，气候逐渐变暖，人类的生存条件好转，上山文化的先民自洞穴走向旷野，在浙江的金衢盆地停留下来，沿着钱塘江的支流营建房屋和村落，过上了定居的生活。

村庄里有经验的前辈们，春天在这里试着播下种子，随着雨季后水位下降，裸露的河岸边如今结出了饱满的稻穗。少年收割稻穗，更年幼的孩子也来帮忙，用手轻拍还没完全驯化的谷粒就纷纷落下了。母亲们用磨盘、磨棒，碾磨橡子、稗子、菱角、块根放到陶制的大口盆中，加热搅动后就能煨成糜烂的粥。

制陶匠人把散落的稻壳、稻叶和稻杆掺在陶土中，可以增加陶土的延展性，提高耐热性能，避免陶器在烧制过程中发生破裂。这些掺杂着稻壳和稻叶的陶片，成为植物考古学家研究的对象。

通过对小穗轴基盘形态的观察，可以确认有些稻穗不是自然脱落的，已有人工驯化迹象。这是世界上最早的碳化稻种子，距今9,000 多年发现于浙江上山遗址 461号灰坑的填土中，它可能是古人在炊煮过程中，偶然掉落在火塘边的。虽然有些残缺，但它仍然保留了早期稻作农业的珍贵信息。

早在更久远的时代，人们就已经认识到野生稻种子能用来果腹，或许他们无意间发现生长于河岸边的这些稻属植物在第二年还能发芽抽穗，于是尝试着最初的种植。尽管作物的完全驯化需要经历数千年的演变过程，又或许稻米在当时人们的食谱中仍非主导，但农耕的出现却永久的改变了历史的走向。

在上山文化延续的近两千年间，炊烟袅袅，人口蕃息……在浦江上山，人们修建了 10 余米长的干栏式房屋，在义乌桥头聚落三面环壕，一面临河，其间发现了房屋、墓葬和窖穴遗迹，反映出丰富的定居生活内容。

安定祥和的生活里，长出了从容有度的生命，比之一万年前就连陶器的面貌也焕然一新，多种形态的盆、罐、钵、杯、盘是为了应对更精细的生活需求。

以白彩施于红底之上的太阳纹，对顶三角纹形似八卦卦象的短线组合纹等，则投射了更深邃的精神世界和更丰富的审美情趣。

这件义乌桥头遗址出土的陶壶，虽历经8,500年岁月斑驳，仍难掩动人魅力，造型简约而优雅，线条流畅而柔和，陶色明丽但不炫目，静静凝聚成历史那头的一抹红。

与此同时，定居社会与农业革命的强音也在北方奏响。河北西北部的张家口尚义四台遗址发现成排半地穴式房址，测年数据达到距今1万年左右，是北方地区最早的明确的定居聚落。遗落在地面上的打制刮削器，工艺成熟的细石器以及大量的动物骨骼，表明食物种类变得更加丰富，但狩猎采集仍然是获取食物的重要手段。磨盘、磨棒和陶器的出现则显示食材的精细加工和烹制已开始普及，定居生活方式的出现成为旱作农业在北方地区起源发展的重要条件。

从目前的考古证据来看，北方旱作农业最早的发生区并非中原，而是更靠北的燕山南北地带。粟和黍都在中国传统的"五谷"之列。北京西郊的东胡林遗址中，发现了迄今所见年代最早的碳化粟粒。粟就是小米，田间地头随处可见的狗尾草，是粟的野生祖本。狗尾草的籽粒狭长，经过人类的驯化，粟粒的形状已经接近球形。

阴山山脉以北科尔沁沙地的边缘地区，兴隆洼的先民们则已经开启了以黍为主要农作物的村落生活。由于它耐寒、耐旱的特点，黍在8000年前的整个北方地区，成为了最重要的旱作农业品种。

今天每一种出现在我们餐桌上的食物，都曾见证了远古先民们的饥寒与艰辛，他们播下稻种、粟种，从此不再被动等待自然的赐予，而是主动参与到万物生长的进程中。中华文明以农为本的基因早已渗透在我们的思维和血液里。

8000年前兴隆洼猎人的队伍，穿行在西辽河流域的平原丘陵间。尽管这里的人们已经过上了以种植粟黍为主的农业生活，但面对将到来的漫长冬季，通过狩猎来储备食物依然不可或缺。

英武的首领是聚落中最优秀的猎手，正带领着族人向丛林深处行进，人们身背长弓，腰悬兽皮箭囊。首领示意猎人们在草丛中隐蔽下来，他侧耳倾听又仰面观察树叶的摆动，前方忽然闪现一只野猪的身影。首领摘下长弓，抽出羽箭。

这令人惊叹的考古现场正是兴隆洼人曾经的家，它的面积达3万平方米，围沟环绕中 180余座半地穴式房屋成排分布，从小型流动的狩猎采集人群，发展为规划有度的大型村落社会。兴隆洼人经历了巨大的环境、资源和社会的挑战。

当年在聚落中心最大的房屋里，首领正在接待远方来的贵客，众人围坐在火塘边相谈正欢。首领拿出一柄刃部装嵌有细石叶的骨埂石刃刀，这是他爱用的狩猎工具。

8,000年后考古学家在这间房屋的东北部地面下，发现了一座墓葬。墓主人是一位50多岁的男性猎人首领。把少数生前或许具有特殊地位的人埋在居室内是兴隆洼文化的特殊葬俗。墓主右侧葬有两只整猪，一雌一雄均呈仰卧状，似乎是祈求猎物丰盛之意。在他的耳部发现了一对玉玦，根据使用痕迹来看应该是墓主人非常珍爱并日常佩戴的耳饰。

玉是美石，经过琢磨、抛光后温润而莹泽，后来的中国发展出了世界上独一无二的玉文化，此为重要先声。

在农业起源的初期，中华大地已经呈现出多元的发展景象，南北方形成了各具特色的农业模式。淮河流域则表现出南北交汇，错落相间的文化格局。

距今 9,000 年到 7800 年前生活在河南舞阳县贾湖遗址的先民，因附近有湖泊湿地，鸟兽成群，便以稻作、渔猎为生。少年蹲在莎草丛边的石上，用湖水仔细的清洗着一段丹顶鹤翅膀上的尺骨，身后的父亲接过尺骨，用手指认真的比度后，刻下几道细线交给少年，少年根据标记用燧石尖刃器小心地在尺骨上钻下了7个孔。

骨笛的制作涉及到复杂的音律和数学知识，在氏族里只有少数人能传承古笛制作和吹奏的技艺。一旦制成，骨笛也往往永远伴随着主人生死不离。

贾湖村落中不同家族的房屋各自成组，错落环绕在中心广场周围，房屋之旁还有制作陶器的窑场，储藏食物的窖穴，死去的人们则安睡在离房屋不远的地下。

墓葬是我们解读过去的锁钥。在贾湖二期的墓葬里344号墓的墓主为男性，无头，代之以叉形器和一组龟甲，其中一件龟甲上还有意义不明的刻符，左臂旁放着一件骨饰﹣﹣两只骨笛，下肢处则放置一些渔猎工具。推测墓主人可能是身具多种才能的智者，担任着氏族的族长兼巫师。

尽管比起数万年前生活已经好过了一些，但人们种植水稻、饲养家猪的规模相当有限，很大程度上还需要依靠渔猎采集获取食物，冬春季节分外难捱。

于是这夜有一场隆重的仪式，氏族成员们集中在村落中央的广场，族长和女祭司肃立其间，高亢的骨笛之声划破夜空。年轻人将酒倒入钵中，用大米、蜂蜜和山楂酿成的酒，甜香沁脾，族长一饮而尽，手持刻符石柄的女祭司高举起双手。成年族人的小腿上绑着装有石子的龟甲响铃，随着舞蹈的节奏怦然作响，少女们的颈下和腰间挂着一串串鸟肢骨管横截而成的小骨环，舞姿轻灵。通过这样的仪式，族人们祈求神灵庇佑；祈求年岁丰穰；祈求氏族人口繁衍；祈求春来水患不扰．...

聚落之外，少年静坐河岸，月光如水洒满河滩，少年抽出那支崭新的七孔笛放在唇边。祖先常在，神灵常在，四时循环，万物化生，万籁俱寂中少年感受到大音希声之妙境，风声、水声、芦苇声、鹤鸣声共作，涌到唇间，送入笛身。他越奏越入神，手指按向骨笛的第7孔。广场上忙碌的人群都为此声惊喜，静止不动，静听余音在星空中盘旋．....

时光荏苒，昔日的少年成为了新的族长。他已经拥有了自己独立制作的骨笛，但仍细心保留着已经故去的父亲当年教他制作的第一支骨笛。那只骨笛后来被失手摔成三段，他便在断裂处钻上几对缀合孔，用线细细缠裹，音色仍旧优美如初。

杭州湾畔的跨湖桥遗址，湖岸边的独木舟被固定在桩架设施中，奖置于船边，这样一停就是8,000年，似乎仍在等待起航。独木舟取材自一整棵的马尾松，制作时先用火烘烤木材表面，然后再用坚硬的石器一点点地木为舟，至今残长犹有5.6米。船旁发现了一些编织物，或许曾是一张船帆，在等待扬帆起航的日子里，它被海浪吞没。那是来自远古的一个信号，面对风浪，面朝大海，面向远方已经站立起我们的少年中国。

8,300 年前古宁波湾的海平面上升，海水灌入四明山北麓，余姚井头山遗址距离现在的海岸约20多公里，当年却是依山傍海。晨雾散去，岸边高地上露出一座村落，用粗木搭建起框架，地板架空在地面上，以竹篾为墙，茅草为顶，这样的干栏式房屋正适合温暖潮湿的海边环境。

井头山人享用了海鲜大餐后，就把鱼骨和蛎、蚝、螺的壳丢弃在村落旁。日复一日，年复一年，堆积如丘，这就是考古学中所称的"贝丘遗址"。

晨光里有人在海岸边踏浪，拾起还来不及随潮水褪去的贝类；有人双手翻飞用芦苇编织着背篓和渔罩；有人忙着收拾晾晒好的肉脯和鱼干，将橡子、麻栋果和核桃一起放进储藏坑，而拥有高超技艺的木匠们，自如的制作着各类木器。

我们所以为的历史往往与权力、战争和王朝更替相关，其实真正支撑人类数千年文明与温暖绵延的，是那些昼出夜伏，炊烟袅袅的日常生活。

石刀刮削木材的声音打破宁静，父亲在修整刚制作好的木桨，他校正杆部是否笔直，桨部的角度是否合适。父亲深知桨在航海中的重要性，在陆地上沿着嶙峋的海岸寻找合适的滩涂，费时且艰难，人们想依靠舟船去更远的地方。少女抬头向东望去，大海是一片灰白而忧郁的原野。

井头山遗址的库房中，这只木桨保存的完好程度令人吃惊，它好像刚刚完成了一次出海。发掘者在10米深的地下找到了这处遗址，这里后来曾被海淹没，覆盖着厚达8米的淤泥，淤泥容易流动坍方，因此在发掘前需要预建钢结构围护的发掘基坑。如果不是考古学家的远见和坚持，我们不会知道这海相沉积之下，曾也是一群人的家。

人与海进退之间，相伴相生。千年之后海平面趋于稳定，陆地重新露出，淡水溪河再次流淌于平原之上，成为人类宜居的家园。一叶竹筏在落日余晖中，摇进了青山绿水环抱的村庄。河姆渡同样的地理坐标，经历海岸线变化下的古环境演变，已然沧海桑田。成排的干栏式木屋错落于菱荷丛生的河湖溪塘间，低洼的平地上稻穗金黄，芦荻作雪飞，小桥流水，饭稻羹渔，颇有几分后世的江南模样。

河姆渡先民制作着翻土的骨耜，精耕细作成为这片土地的传统。与此同时人们对生活更多了分精致的追求，炊煮食物用釜，蒸食物用甑，煮水用盘，盛放食物用豆、盘、钵和盆。简单蒸煮食物以求果腹的时代已经过去，这些优雅美观的饮食器具的制作者似乎正于日常生活中感知着烹饪的愉悦。

这只圈足木碗外表还晕着朱红色的漆。时光已逝，考古学使它复活，那消失的生命的温度重返人间。蚕纹象牙杖端饰、黑陶朱纹钵，这些堪称原始艺术的物件都是先民智慧与信仰的结晶。

这件雕刻"双鸟朝阳纹"的象牙蝶形器，器身正面中央阴刻五环重圈纹，周围绕以火焰状光芒，两侧各有一只鸟昂首相望，又似在引吭啼鸣。

远古先民在漫长的采集渔猎和农业活动中，观察季节变化与万物生长，体验寒暑交替和昼夜节律，产生了对太阳的敬畏与向往。当他们把这种情感寄托在如鸟一般能通天达地的动物身上，神灵崇拜就应运而生。

一轮红日正冉冉升起，驱散了山间的溟漾晨雾，沅水在湘西的千丘万壑间蜿蜒穿行。远古的匠人以深红色的着色剂，用软笔在白陶簋形器的外底部绘出一个太阳纹的图像，如从远古幽冥中穿云破雾而来，瞬间光芒万丈。

这件亚腰白陶罐，肩部饰篦点凤鸟，雄健的勾喙，敏锐的环眼，张扬的双翅，展示着非凡的神采。它的身侧是蝶形兽面图像，似乎蕴含庄严的神力。这类刻画着神圣图像的白陶器是高庙文化的典型器物。在这件白陶罐的领部由戳印细篦点构成的双羽翅獠牙兽面纹，构图复杂，如同一张长满利齿的阔口，出獠牙，还伴有羽翼和羽饰。兽面两侧各有一座云梯萦纤的"高阙"，它是不是远古祭祀中的天梯形象呢？

7,000多年前，这些构思诡谲的通灵祭器犹如一道天光照亮了幽暗的史前丛林。它产生于远离江湖平原的河谷山川，那里充满了神秘的色彩和浪漫的情调。在洪水频发和瘴疫横行的年代里，疾病与死亡的威胁无处不在。为了获得神灵的恩赐与庇佑，先民奉献出最虔诚的信念和最精湛的技艺，以获得心灵的慰藉。

高庙文化并不是一个人丁兴旺、持续长久的考古学文化，却凭借一群狂热的天才"艺术家"和他们创造的神秘作品在历史上留下了独特的影响力。

随后以洞庭湖为中心的汤家岗文化继承了它的传统，在他们的白陶器上大量涌现各种复杂的几何纹样。有学者猜测，八角星纹代表一种宇宙观，它常见于圈足盘的底部，把这些器物倒扣放置正如天体的模型，中间的圆形是天，外侧的八角则象征大地的八方。无论如何，器具本身只是一种实物形式的载体，它真正要表达的是人们寄予这种载体之上的精神信仰。原始信仰驱散了许多萦绕在人们心中对未知世界的恐惧，也成为不同人群、不同聚落之间连接的纽带。

从汤家岗文化所在的长江流域向北越过绵延的秦岭，就到了地理意义上的北方。秦岭北麓发源的渭河汇入黄河之处，正是辽阔的"八百里秦川"也是丰沃的关中平原。

沪河东岸的半坡村落，一座半地穴房屋刚刚建造完成。屋顶新铺的茅草散发着清香，男人抬入盛满谷物的大陶瓮放在后墙边，又将陶盆、陶钵、陶罐、石斧、石铲等一一摆放好。进门的左右两侧各有一个矮矮的土台子，那便是床铺。女人给它铺上新编的草席和松软的兽皮，火塘中灶火正燃，褐色的夹砂陶罐架在灶上，外表布满烟苔，里面的小米粥已经煮熟，颜色金黄诱人。

距今 7,000年至6,000 年左右，新石器时代的社会得以显著发展。仰韶文化在黄河流域开启了属于自己的时代，其早期被命名为半坡类型，一个个围沟环绕的农业村落像点点繁星一般出现在黄土大地的川谷盆地间。

半坡村落的中心广场上燃起一堆篝火，一只整猪被对半剖开放在木架上烧烤，浓香四溢。新居收拾停当，庆贺乔迁的宴饮拉开序幕。聚落中的男子们在长长的草席上落座；氏族的长辈们被安排在首座；妇女们端出烹制好的食物，有小米饭、羹汤还有烤熟的猪肉。男人们抱起小口尖底瓶浑浊的米酒倒满一盆，他们用小杯舀酒向老者致敬，然后一饮而尽。女人和孩子们也开始享用丰盛的食物，整个聚落一片欢声笑语。

考古学家在半坡聚落的壕沟内发掘出45座房子，200 多个窖穴以及牲畜围栏遗迹。整个居住区以一条小沟分为南北两片，中间有道路相通，壕沟以北有公共墓地，以东则有公共窑场。每个房屋内居住着一个家庭，每个单元是一个扩展家庭，而整个聚落就是一个氏族。

临潼姜寨遗址更清楚地表现出聚落规划中对亲属关系的关注，圆形围沟内房屋有100多座，明显分成5个单元，围绕着近4,000平方米的中心广场分布。每单元都有大中小型房址，所有房门都朝向中心广场，在这样的空间格局中每一个居住者都时刻感受着个人与家族部落的血脉联系。

逝去的亲人按照亲属关系分组埋葬在聚落周围。考古学家发现一种特殊的埋葬方式对血缘的强调近乎极致。河南灵宝的城烟遗址发现多座多人合葬墓，均属仰韶文化早期。人骨成层安放，其中第94号墓内共有人骨19具。这些人骨是从之前不同时代的墓坑中被迁移而来，重新埋葬在一起，似乎更多在彰显群体的凝聚力和亲密关系。生则共居，死则同葬，血缘不会被死亡割断。

皓月当空，半坡聚落的仪式进入高潮，酒酣之际有人吹起陶员，悠悠咽咽在夜空回荡。少年被埙声打动，悄悄起身独自回到居室。那置于高处的彩陶盆是少年的牵挂，它是为聚落里逝去的婴儿准备的，盆的内壁用黑彩绘出人面鱼纹。

仰韶文化半坡类型的彩陶中最有代表性的象生图案便是鱼纹，有些是单独出现而较为写实的鱼纹；有些是以直线与弧线描绘，圆点、弧线和弧边三角穿插而显得活泼灵动的简化鱼纹。纹样格式除平展式外，还出现了回旋、跳跃等姿态。

半坡类型晚期，彩陶上的单独鱼纹采取了夸张变形的艺术处理，变成上下对称的式样，趋于几何化。鱼与人面相结合的形象是半坡彩陶独有的图案，如著名的人面鱼纹，人和寄寓又互相转借，意味着人和鱼是交融的共同体。

被人格化了的鱼类图像和各式鱼纹，可能具有半坡氏族保护神的性质，将早夭的婴儿装殓在陶瓮中，以瓮为棺，以盆为盖，埋入土中，这是一种叫瓮棺葬的习俗。人面鱼纹绘制在陶盆内壁，当它倒扣在瓮口的时候，画案只有其内的亡灵才能看得到。瓮棺上还特地留出一个圆孔，让孩子的灵魂可以自由出入。

百万年来无边的夜空曾令我们的先祖们畏惧，猛兽、寒冷、黑暗，总是一再侵袭的饥饿，随时突然降临的死亡使每一个夜晚都曾如此艰难和漫长。直到第一双拾取火种的手开始让光亮在寒夜驻守；第一双制作工具的手开始让万物可被改造利用；第一双撒下谷种的手开始为来年种下希望。

为了生存我们的祖先一次次走向远方，平原、山地或是丘陵、海岸，他们披荆斩棘，筚路蓝缕，聚族而居，生生不息。从此这片苍茫大地上夜空深邃，群星闪耀，不再惧怕黑夜的人终于能抬头仰望星河。


\section{思辨句式:我们所以为的***，然而，真正***，恰恰是***}

\subsection{原文}

我们所以为的历史往往与权力、战争和王朝更替相关，其实真正支撑人类数千年文明与温暖绵延的，是那些昼出夜伏，炊烟袅袅的日常生活。

\subsection{思辨结构解析}

提出大多数人认同的常识，进一步分析这背后令人吃惊的本质。（甚至相反）

\subsection{写 2024 新高考一卷}

\textbf{作文题：} 随着互联网的普及、人工智能的应用，越来越多的问题能很快得到答案。那么，我们的问题是否会越来越少？

\textbf{仿写：} 我们所以为的技术革新总是与算法、数据和机器进步相关，其实，真正在暗流涌动中引领文明破浪前行，赋予人类以跨越式解放与蜕变的，是那根孜孜不倦、探索无界的好奇心，是那些深邃不息、发现并挑战未知之难越的灵魂。


\subsection{写 2024 全国甲卷}

\textbf{作文题：} 每个人都要学习与他人相处。有时，我们为避免冲突而不愿表达自己的想法。其实，坦诚交流才有可能迎来真正的相遇。

\textbf{仿写：} 我们所以为的缓解冲突往往与保持沉默、息事宁人相关，其实，真正能够化解矛盾、消融分歧、弥合人心裂痕的，是那些坦诚相待、推心置腹的深入交流。在冲突面前选择沉默，只会致使矛盾在岁月的冲刷下愈发根深蒂固；选择坦诚交流，才能共同耕耘出理解的沃土，迎来真正的相遇。


\subsection{写 2024 北京卷}

\textbf{作文题：} 几千年来，古老的经典常读常新，杰出的思想常用常新，中华民族的伟大精神历久弥新……很多事物，在时间的淬炼中，愈显活力和价值。请以“历久弥新”为题目，写一篇议论文。

\textbf{仿写：} 我们往往将历史的磅礴归因于那些璀璨的经典、深邃的思想与崇高的精神，然而，真正赋予“历久弥新”以深刻内涵的，恰是那些经过时光雕琢，在人民群众日常琐碎中不断焕发生机与意蕴的生活细节。


\subsection{写 2024 天津卷}

\textbf{作文题：} 在纷纷的世界中，无论是个人、群体还是国家，都会面对别人对我们的定义。我们要认真对待“被定义”，明辨是非，去芜存真，为自己的提升助力；也要勇于通过“自定义”来塑造自我，彰显风华，用自己的方式前进。

\textbf{仿写：} 我们常常误以为他人的定义即是我们存在的全部意义，然而，真正决定我们本质与未来的，恰恰在于我们如何审慎对待外界的“定义”，明辨是非，去芜存真，以此作为自我提升的助燃剂；更在于我们是否有勇气进行“自我定义”，以独树一帜的姿态塑造并展露内在风华，坚定不移地沿着自己选择的道路砥砺前行。


\subsection{写 2024 上海卷}

\textbf{作文题：} 生活中，人们常用认可度判别事物，区分高下。请写一篇文章，谈谈你对“认可度”的认识和思考。

\textbf{仿写：} 我们惯于以浅薄的认可度作为评判事物高下的准绳，却忽略了其背后更深层的本质与意义。其实，真正剖析“认可度”的内涵，我们会发现，它也可能是盲从与主观臆断的结合体，而事物的真正价值，应超越外界的喧嚣认可，直抵其内核的深刻与独特。


\chapter{第三集 \quad 星斗}


\section{正文}


长江、黄河自雪域发源，宕落高原蜿蜒于群山与莽原之间，向东汇入大海，古老的中国文明在大江大河之间生长。然而北有荒漠雪山、西伯利亚荒原和蒙古高原横亘；门上有青藏高原、横断山脉阻隔；东、南皆是茫茫大海。

早期的先民基本处于一个相对独立的地貌环境，与南亚、西亚、中亚和地中海等区域相比，古中国文明与欧亚其他早期文明之间道阻且长，难以建立起直接的联系。

于是在漫长的时间岁月里，这片大地上的先祖们经风雨，历霜雪，斩斫着沿路荆棘，探索着漫漫前进之途。

终于在距今5,500多年前后，因着长期的稳定与繁荣，人烟蕃息之处越来越多，晋陕豫三省交界处这片横楔在秦岭与黄河之间的狭长平原出现了高陵杨官寨、华县泉护村、灵宝北阳平、灵宝西坡这些面积数十万甚至百万平方米的仰韶文化大型中心聚落。

初夏野蔷薇盛开在华山脚下的泉护村，绘陶青年的面前摆满了兽骨、兽牙、树枝等绘画工具以及各种可化为颜料的矿石。他将所需矿石砸碎，再放置研磨器中，用石棒研磨成粉。这是迄今发现最早的一套完整的仰韶文化绘画用具，由陶杯、石盖、石砚、研磨棒和氧化铁颜料组成。石砚臼窝内壁及砚面上还残留红色颜料的痕迹。

在青年工匠的笔下行云流水地绽开一朵朵花，经过火的淬炼，黄底黑彩的陶器便能再现5,500年前那个夏日人们所见过的绚丽，闻过的香气，陶上怒放的野蔷薇将不再步入花开花谢的轮回，而将永世盛开。这类绘有花瓣纹样的陶器，被今天的考古学家命名为仰韶文化庙底沟期彩陶。

仰韶文化历时近2千年，庙底沟期是它的中期，因河南三门峡庙底沟遗址而得名。花卉纹是庙底沟彩陶最具代表性的纹样，除此之外还有回旋勾连纹、豆荚纹等。

在低饱和度的黄色或红色素地上，以黑彩、白彩或红彩描绘纹饰，通过连续、反复、对称、共用的手法，实现了写实与抽象的意向表达，阴纹与阳纹的相互映衬，平视与俯瞰的视觉效果。

曾经的仰韶文化早期各地虽然都流行彩陶，却各具区域特色，到庙底沟时期从关中平原山西南部到河南西部，彩陶的纹饰模式高度近似，这跨越距离的呼应却是因何而起？

考古学家发现史前人群的分野往往不仅以血统为要素，更以生活方式或文化传统当作彼此辨识与认同的密码。

庙底沟人原以旱作农业为主，也逐渐接纳了稻作农业的升级方式，这就是不同人群接触交流甚至融合的重要迹象。

不仅如此，籍由成熟早作农业人群的扩散和流动，庙底沟彩陶的影响先是到达了中原河洛地区，再以整个黄河中上游为起点，挟风行之势，向着北抵阴山，东到大海，西达甘青，南至长江的更广阔天地辐射。

向东山东海岱地区大汶口彩陶的叶片纹、花瓣纹和旋纹是自庙底沟彩陶掀来的风潮；向西传承陕西关中的甘肃彩陶继续朝着更远的陇西和青海东界传递影响；在稍晚的马家窑文化时期，甘青地区将发展成为新的彩陶文化中心；北溯汾河和南流黄河两岸的谷地、高原而上，庙底沟彩陶为内蒙古中南部及东北的红山文化彩陶注入了新的生命力；南延汉江而下，庙底沟彩陶在长江中游的豫南、鄂西北引发追捧后，继而流播至湖南、浙江等地。

古老的华夏大地由此拉开第一次大范围交流、大规模互动的序幕。在没有广告效应和发达物流的史前时代，器物传播依靠的是人与人之间的传递。

在各地都已发展出成熟陶器传统的那个时代，由庙底沟掀起的彩陶艺术浪潮能够波及这样广阔的区域，应该代表着远距离交流的频繁以及以彩陶为基础的文化共识，这或许就是文化意义上的最早中国的雏形。

在后世的《史记》中保留了有关黄帝这位华夏共祖的传说和记忆。据说他的足迹东至大海，西至陇右，南至长江，北至幽燕，与庙底沟彩陶的影响范围惊人的重合。

在古汉语中"花""华"同音，"华"的本意为"花"。金文中的"华"字就是花朵加上花蒂的样子，因此考古学家苏秉琦将庙底沟彩陶中常见的纹样喻为"华山玫瑰"，而把创造庙底沟之花的人群称为最初的"华人"。美丽的花卉纹似乎可以解读为华夏之"华"的最初由来。

土地肥沃河流纵横的关中平原是仰韶文化绽放的核心区域之一。泾河与渭河交汇处的杨官寨遗址面积达百万平方米，考古发掘揭露出由大型环壕、西门址、中心水利系统、东区墓地等组成的规划有序的村落布局。东区成人墓地已经发掘500多座墓葬，显示了聚落繁荣发展的面貌。

杨官寨遗址出土的特殊器物，如镂空人面漆座、朱砂人面饰陶器、动物纹彩陶盆等，可能与聚落进行的重大祭祀活动有关。表明当时已经拥有相对复杂的政治组织，为社会向更高一级发展奠定了基础。

秦岭余脉绵延向东直指中原，黄河南下穿过龙门转折东流，河南西部灵宝的铸鼎原，南依秦岭，北面大河，是仰韶文化的另一个核心区域。其中西坡遗址是经发掘确认面积最大的聚落之一，它的中心是4,000多平方米的广场，广场四角各有一处大型建筑。

室内的火塘中烈焰正燃，一场庄严的成年礼即将在这里举行。一位青年站在大房子之外难以激动的心情，只等完成这场隆重的仪式后他就能成为真正参与族中事务的一员，这是他自孩提时代起就抱有的梦想。

大房子的建造曾是村中最大的工程。不包括砍伐、加工和运输树木等工序，单单建造房子本身就让百余名壮劳力施工足足3个月，彼时还是稚龄的他充满好奇总想一探究竟。建造这样的大型建筑，先民采用了相当复杂的工艺。先要开挖房基坑并垫土夯打，在基坑周围挖出墙基槽，槽中竖立粗大的木柱，用夯土填满基槽使木柱牢固。通过夯打与夯填形成内外墙体，此后设置室内柱础石，挖门道、火塘，栽立室内柱，再用不同种类的黏土和夯土铺垫坑底，以石灰膏泥抹出平整的地面和墙壁，最后依托木柱搭建起回廊、门道以及屋顶。这类大型建筑的出现体现了聚落上层日益强大的组织能力，是庙底沟社会发展到新阶段的又一有力证据。

当年的西坡少年曾趁着众人休息时偷偷溜进大房子里，眼前所见令他无比吃惊，就像有几回夕阳余晖染红天地河的景象，墙壁和眼前的柱子都被刷成了红色。恍然之间他发现位于中央的一根柱子还没上完色，他知道这是大房子里最神圣的区域，他忍不住蘸了颜料就要去刷柱子，却被门外突然的响动吓个正着。

未成年的族人是不允许进入这样的重要场所的，可出乎意料的是面对这擅自闯入的少年，族长却没有一丝恼怒，他轻轻扶起眼前的孩子，示意他应该从上往下刷。少年迟疑地探出手臂却一不小心将红色的颜料甩落在族长的脸上。

当这回忆沉淀，步伐就愈加坚定，那个当年偷溜进来的孩童已经长大成人。穿过狭长的门道绕过烈焰腾腾的火塘，进入族人们簇拥着的庄严空间。这铺天盖地、动人心魄的红，红色的立柱、红色的墙壁、红色的地面仍一如当年．..

这座大房子曾见证过无数重要的时刻，西坡和附近小聚落的居民们在特定的时刻总会汇集在此，举行各种各样重要的仪式。神圣的宗教典礼、社群的公共仪式、盛大的婚礼、集体的宴饮、狂欢的节庆。今天的主角是这位西坡青年，在成年礼之后他就要为族中承担起更大的责任。

后来考古学家在西坡遗址的中心发掘出这座大型半地穴式建筑遗址。房址的室内面积为204平方米，加上外部回廊，总占地面积超过500平方米，是迄今发现的那个时代最大的单体建筑。

庙底沟时期许多仰韶文化聚落中的大房子规模宏伟、工艺复杂。他们往往位于聚落中心，宽敞明亮，结构考究，地面层层加工铺设，墙面多施彩绘，想必不是普通的住所而是整个聚落的公共聚会场所。

当屏息静气，仿佛已坍塌的房屋，逆转时光再度矗立，还能听见某个首领或慷慨或舒缓的陈辞，还能听见聚在这里的人们，那热闹的喧哗。

随着彩陶的广泛传播，仰韶文化的社会内部和外部均建立起空前紧密的交流网络。中原和位于东部的海岱地区便有着频繁而持续的互动。

父辈们曾经赶赴大海之滨的壮阔旅程，远方来客不断诉说的奇风异俗，频频拨弄了青年的心，踌躇满志的西坡青年决定在成年礼的第二天就向东出发。

西坡青年沿着黄河一路向东，从西坡到大汶口的距离是600多公里，如果按照每小时行进5公里，每天行走6小时计算大约需要20多天。青年朝行夜宿，路上看过几十轮日升月落，穿过辽阔的平原行至泰山脚下，终于到达了大汶口聚落。

当黄河中游经仰韶文化的整合而掀起风起云涌的变革时，下游也正由大汶口文化引领着书写下气势恢宏的新章。大汶口文化自岱宗泰山和渤海黄海之间的古济水与泗水、汶水流域崛起，一统今山东省全境及皖北、苏北、豫东等地大有风雷之势。

西坡青年受到了大汶口人的热情款待，参加了令他生平难忘的一场盛宴。宾主面前陈列着各种食器和酒器，陶豆中装满了肉干和水果，陶鼎内盛满了金黄的黍糜，精致的高柄杯中斟满了清香的美酒．..

大汶口社会上层豪气而富有，热衷于举办盛宴款待宾客，这也是他们展示财富和威望，广泛交友，扩大往来的媒介。

"好客"或是远古的基因传延至今。宴饮的奢华并未让西坡青年过分艳羡，席上尊卑有序的饮酒进食之礼，却让他格外留意。

考古研究者发现大汶口文化的大型墓葬往往随葬大量食品以及与饮食相关的陶器。大汶口遗址2005号墓的棺椁之间有数组陶器，一组是饮器，如高柄杯、觚形杯；一组是食器，如鼎、豆、壶、三足钵、三足碗等。椁外放置了成套食器，器内多盛有肉食。

大汶口文化不同区域的大型墓葬，几乎都采用了类似葬仪。这说明即便地跨山东、苏北和皖北，各聚落的社会上层之间仍保持了频繁的交流。

相似的考古发现背后是一致的文化认同和社会秩序。以饮食宴享和成套器具为媒介构建的礼仪，则是大汶口文化权利与地位的象征。席间那些绘有花瓣的彩陶瓶，让西坡青年觉得格外亲切，这他熟悉的故乡风物，是他的先辈与这东方之族曾远程往来的凭据，也是仰韶彩陶之风被初兴的大汶口文化效慕接纳的明证。

主人一杯杯地劝，客人便一杯杯地饮。交谈时西坡青年注意到，在座不少大汶口人的上颚都缺了一对侧门齿，这种拔牙之风乃是族中世代相传的习俗，或许是要用对于痛苦的忍耐来显示勇气，宣告成年。

那天的盛宴上除了西坡青年，还有另一位来自异乡的年轻人。在一片欢声笑语中，那位异乡人却显得拘谨而沉默，他是从淮河流域北上的凌家滩人。

大汶口人夸耀着丰盛的食物和财富，强调着世俗社会的秩序和礼仪，而他所来自的凌家滩社会却更注重宗教精神的引领。

现实的欢乐和规仪并不能打动凌家滩青年，唯一让他心有戚戚的是主人手边的一件陶豆。红陶之上以白彩描画的八角星纹格外显眼。在他的家乡八角星纹是神圣的标志，这难道是一种巧合吗？

凌家滩在安徽巢湖一带，与同在长江下游的崧泽人、北阴阳营人一样都有着尚玉之风。大部分玉器都是非实用性的装饰品或仪式用品，这意味着那个举全族之力只为生存的艰险岁月已经过去了。人们产生了除温饱之外的更多需求，社会已能释放出部分人力来专门制造这些"无用之用"的器物。

那时的玉器制作工艺已经相当成熟，开坯时广泛采用线切割或片切割的技术；开孔时则采用了管钻技术，坚致的玉能被人们自如地雕琢成想要的样子。

凌家滩青年的父亲是部落里最好的玉匠。父亲的那双手简直有着无穷魔力，他在童年时总是如痴如醉地观看父亲切割、穿凿、琢磨，一块块看似不起眼的石料经他之手便能化为精巧的美玉。乍见这片玉板之上的"八角星"，他蓦地怔住。回过神后不断地追问父亲其中的含义，父亲只是笑而不语，让他长大后自己去寻找答案。

父亲曾为族里离世的祭司和首领制作过数件玉人。玉人或坐或立，双臂紧贴身体，弯曲上举于胸前，臂上穿戴着成排的镯饰，似是在虔诚祈祷。这种姿势颇含仪式意味，是族人们在特定时刻的程式化肢体语言，也是祖先入葬时的造型。

最令青年赞叹的还是父亲制作的一件玉鹰，那鹰展翅飞翔，仰首侧视，展开的双翼各如一只猪首，敞露出胸腹部赫然的八角星纹，说不出的神秘莫测。

凌家滩人在玉器中寄寓着他们的信仰观念，在日益分化的社会里也以玉器作为展现财富和权力的重要载体。凌家滩墓地建在大型祭坛之上，逝者按照身份的区别被葬入不同的区域，拥有多寡不等和种类不一的随葬品。

07M23号墓位于墓地南侧的中央区域，墓主人的遗骨因数千年环境腐蚀已经荡然无存，但随葬的300件玉石器却印证了他的首领身份。

考古学家通过详细的现场记录和反复摆放分析，复原了当时复杂又严格的礼俗。葬具的最底层用穿孔石钺横向摆成数道"横木"的样子，其上就是七排严丝合缝整齐排列的石、石凿，刃部朝向墓主人的脚步，密密麻麻地砌出一座"平台"。在长2.65米，宽0.8米的葬具内，墓主人的遗体就曾安放于这象征尊崇身份的基底之上。

一种特别的"花石头"被专门用来制作大孔弧形刃部的特殊石钺，在胸口部分放置两排六件，然后沿着墓主身形的中线，整齐往脚步首尾相连排成一线，这些石钺显然没有使用过，入葬时也未装柄，而是摆成一个"T"字形，成为墓主置身的第二层基础。

玉器是身上最重要的配饰，墓主脖梗上挂着一套穿缀繁复的玉璜组佩，两臂上分别穿戴10件玉镯，腰腹部则穿挂一套内置玉签的特殊礼器，造型如龟壳互扣。穿戴整齐的墓主人身上还要继续用石钺和玉环、玉玦来铺陈。

石钺从材质和技术上都能看出三件一组摆放的规律，分别压在身体两侧，尺寸从大到小。成套制作的玉环、玉玦密密地挂坠在葬具的两端，墓坑内散落的玉环可能是最后才撒在葬具上的。除了随身佩戴的玉饰外，墓中出现的大量的玉石器都是为了这样的仪式活动而定制的。用特殊的玉石资源，结合专门技术生产的手工业"奢侈品"由此成为了早期中国体现"身份"的重要内容。

在敷土填满整个墓坑的最后，一件重达88千克的玉石猪被特意摆放在墓土中。这件长 72厘米，宽32厘米的玉石雕是目前中国史前发现尺寸最大的作品。整个葬礼奢华又严格，是凌家滩聚落体现个人地位的重要仪式。

显贵者们被埋葬在墓地的南区，青年的父亲去世后跟其他族人一起葬在西北的墓区。编号98M20的墓葬中除了凌家滩常见的玉石器外，还出土了111件加工玉器后所剩下的玉芯和材料，正是这些特殊的随葬品为考古学家推测墓主人的身份和职业提供了线索，他可能是一位专职的玉工。

一个好的匠人不仅要有精湛的技艺，还需有内心的充盈。父亲的这句话他不曾忘怀。多年前父亲曾听人说起遥远的北方有一处圣地，那里的人们也善于琢玉，并通过仪式中复杂精美的玉器来表达信仰与寄托。父亲心向往之，但终未成行。如今对父亲最好的纪念就是走他想走的路，寻找那跨越生死与时间的终极意义。

是时候离开华服盛宴的大汶口了，凌家滩青年与西坡青年道别，又向好客的主人家辞行，坚定地向着圣地红山行进。

凌家滩青年翻过泰山，沿着茫茫的海河平原向东北前行，当他穿过辽西走廊来到巍峨的大青山脚下便进入了红山文化的核心范围。努鲁儿虎山深处，牛儿河畔的山谷间有一道漫长的山梁，当地人依河名，称其为牛河梁。穿过起伏的山峦，圆坛、神庙和一座座积石高家，在林端时隐时现。

红山人对精神文化的需求高于一切。公共权力自宗教层面产生形成了层级严明的体系，而牛河梁拥有着其中规格最高的、完整的坛、庙、冢祭祀体系。于山梁之上修建山台，西南建筑群采用中轴对称布局，轴线南端直指远处的猪首山山峰，如此规划可谓苦心孤诣。然而红山人的空间视野远不止于此，在牛河梁周围近50平方公里的范围内，他们充分利用山形地势的变化，勾勒出宏大的布局。

云雾深处，红山女巫的身影若隐若现。凌家滩青年亦步亦趋，登上主梁。沿着漫长的山脊穿越掩映的密林，青年来至女神庙前。

这是北方史前时期流行的半地穴式建筑，但采用了较少见的多室结构，平面呈"亞"字形，总面积约 75平方米。

走进狭长的门道神秘庄严的气氛扑面而来，神庙之内装修极为考究，地面经火烘烤，木骨泥墙上装饰有仿木建筑构件，墙面还绘有赭红、黄白相间的彩色几何纹样。相当于真人三倍、两倍和等同的数尊女神塑像，伫立于北主室内，女神头像眼嵌玉石、齿贴蛙饰。

在那个时代产妇和婴儿的死亡率很高，人们哀劬母亲诞育子女的不易，也期望人口能顺利繁衍，将之投射到信仰中便是对于女神或母神的崇敬。

走出女神庙时已是黄昏，落日映红牛河梁的峦。山间散步的积石冢是族中先人们的安葬之处，也是后代族人崇奉先祖的场所。

一冢之中往往有多座墓葬，因墓主身份的不同而被分别安置在大小各异的石砌葬具之中。在牛河梁的积石冢内，玉器几乎是唯一的随葬品，但数量并不算多。大型墓葬中一般也不超过10件，组合也不十分固定，与材质单一数量不多的随葬品形成鲜明对比的却是积石冢那异常巨大的规模。

这种不相称恰恰反映出"唯玉而葬"现象的内涵，那便是玉器无关财富的占有，也不刻意显示等级的差别，而是与特定的象征功能相关，或许是作为通神的工具。

5号地点1号家的这座中心大墓，墓主头部两侧对称摆放着玉璧，腹部放置勾云形器，双手各握一件玉龟。

2号地点1号家的南侧墓群中，埋葬着一位富有的成年男性，随葬玉器20件，斜口筒形器、兽面牌饰、勾云形器、玉龟等等。这些明显具有观念意识的器物折射出红山文化积石冢和葬礼仪式中浓厚的原始宗教特色。

双手置于胸前的玉人、蜷曲的玉龙、斜口筒形器、玉龟……不少红山文化玉器与凌家滩文化的玉器颇具相似之处。环绕积石冢的无底彩陶筒形器，器表施有勾连涡纹，又让人联想到庙底沟彩陶的影响。这些发现似乎证实着5,000多年前黄河中下游、长江中下游和辽河流域等地的社会上层之间已经存在远距离的物质和精神文化交流。

青年望向不远处的圆形祭坛，红山女巫已经站立在祭坛中央。祭坛由三个同心圆石圈构成，每一圈的边缘都竖立着多棱形的花岗岩块，最内圈的整个地面铺满白色石灰石。由外至内的三圈或许分别象征太阳在冬至、春分、秋分以及夏至运行的轨道。这样的结构特征与唐长安城南郊圆丘、明清天坛较为相似，体现出古人对天地方圆早有了超越时间的共通之理。红山女巫默诵祷词，双手置于胸前，虔诚而肃穆地展开双臂，如同展翅的雄鹰将飞升而上。凌家滩青年目睹这庄严而神圣的情景，内心如有和鸣，豁然洞明，唯有将心灵与精神托付于那伟大的自然规律，敬畏天地，方能得到持久的安宁。

西坡青年在大汶口停留了一段时间后踏上了回家的路。回顾这段旅程，大汶口的礼仪和物质生活令他印象深刻，而那些与他所在的中原既有联系也有区别的种种是一段更加难以磨灭的记忆。

他迫切想把此行的所见所闻告诉部落里的每个人，见识了更大的世界是一个青年成人的骄傲。当青年滔滔不绝地讲述起这段旅程的种种，是否会引起族人们长久的惊叹？考古论证止步于此，难以描述生动的细节，但这位经过漫长旅行洗礼的年轻人或许终将成为西坡村落新一代强有力的领导者。

生活似水，又流过了几代人的光阴。西坡村落也在悄然发生着变化。自先辈肇始的对外交流，逐渐扩容了本地文化的基因。

考古学家从这一时期的大型墓葬中辨识出了仰韶先民开创独特社会发展道路的足迹。5,000多年前的盛夏，在黄土塬上的西坡聚落中正举行着一场盛大的葬礼。这是显示社会身份，维系社会组织和秩序的重要仪式。

精细的考古发掘可以让我们推想那场葬礼的部分细节。整个葬礼大致是从遗体处理开始的，首先要清洗逝者的身体和头发。西坡墓地中有些墓主头戴发簪，应该就是清洗头发以后插上的，待清洗过程完成后，人们用织物小心地包裹逝者的身体。在下葬的日子，长长的送葬队伍从中心广场出发，缓缓南行，由聚落的南门跨过壕沟来到地势高亢的墓地，应该会有幽咽的陶损吹响吧。

苍翠的秦岭近在眼前，回首北望黄河隐没在远山之间，在众人肃穆的注视下，包裹着麻布的墓主人被轻轻放置在幽深的墓中。

人们仔细地摆设随葬品，然后用木板封盖墓室和脚坑。

经过多学科研究，我们对这位身高约 1.65米年龄在35岁左右的男性墓主有了更深的认识。他的右侧肋骨中部有骨折错位愈合现象，生前可能因参加狩猎、战斗或竞技而受伤。寄生虫专家提取他盆骨内的土样进行分析后，发现了大量因食用猪肉滋生的寄生虫卵，而猪肉是那个时代奢侈的美食。

我们或许永远无法知道这位男子姓甚名谁，但有一点可以明确，他是庙底沟社会等级化的代表，是距今5,000多年前贵族阶层中的重要一员。

值得注意的是他下颌的两颗门齿在生前就被有意拔除，这是大汶口人流行的风俗。他的墓中还随葬着东方风格的大口缸或许随着大汶口文化的发展壮大，连带着当地的一些饮食器具已逆着彩陶东去的路途，反向输回到中原地区。

庙底沟的先民建造了大型聚落、大型公共建筑以及大型的墓葬，但考古发现中既无奢华的随葬品，也无浓厚的宗教气氛，而用来凝聚族群和巩固权力的大墓葬仪却可以看到跨区域远距离文化交流的借鉴和影响。

简朴实用与多元开放成为中原地区社会文化的特质，对于此后中国文明的形成与发展有着深远影响。

距今6,000年以来，随着庙底沟文化彩陶的扩散，大汶口文化的西渐以及东方各玉文化中心的形成，黄河流域、长江流域和西辽河流域间的文化互动交融日趋紧密。

这满天星斗竞相闪耀的大地，塑造了中国的地理核心。从此区域文化各有所长，发展并进，积累起中华文明多元化的深厚土壤。

区域间日趋深化的交流与互动为一体化的将来积蓄着力量。那一群群阔步山川间的年轻人正如我们的青春中国，已昂首伫立在文明盛放的前夜。


\chapter{第四集 \quad 古国}

\section{正文}

蜿蜒清翠的天目山向着衔接于东麓的杭嘉湖平原，伸出两支余脉，如一双臂膀环拥着余杭山水。

5300年前在散落于平原的孤峰残丘之上，村落连绵迭起，人们在山林狩猎采集，在湿地种植水稻。当人口逐渐蕃息，耕作技术发展，村庄和稻田便向平原深处蔓延。为了防备洪涝灾患，先民们堆土成墩，居住在墩上，死后也葬在墩上。取土时形成的低洼地带被顺势整治成河陂，再与自然水网连通便成为来往运输的水道。

时光荏苒，三四百年后这里渐渐成为方圆百里间的中心，良渚古城拔地而起。

4900年前的夏秋之际，一场大雨倾泻而下，袭击了整个天目山区。掌握尖端制玉技术的良渚匠人正在雕刻一件玉琮。神徽的纹理细如发丝，稍有差错便前功尽弃。这时屋外突然电闪雷鸣，眼前变得晦暗不清，他知无法再继续只得搁下手中的刻刀。

此时的莫角山上，王族们隔着雨幕向远处望去，忧心忡忡……连日暴雨威胁着良渚城的安危。在漫长的时光中，水滋润着长江下游这片丰腴的土地，发展出高度发达的稻作农业。洗练出自马家浜始，经岭家滩、松泽至良渚这样数千年间一脉相承的区域文化。然而，水也曾一次次化身为惊涛和洪涝，阻断先民迈向文明的步伐。是以 5000年前良渚先民决定在这世代享溪的水乡平原上修建这座城址时，选址和布局都经过精心的考量。

良渚城三面环山形成天然屏翳，东边的腹地则平畴千里，水网交织，是再好不过的宜居之地。只是夏季多雨将导致北面山洪倾泻；秋天伏旱又会使水稻绝收、河道干涸。

良渚城的营建者们深知要想突破发展的桎梏，必须解决这频发的水患困扰。于是他们开始构画一个宏大的水坝系统的蓝图，经考古发现的11条堤坝遗址构成了具有上下游两级水库的完整水利系统。

西北方向大遮山中的岗公岭、老虎岭、周家畈、秋坞、石坞、蜜蜂垄组成了高水坝群；西南方向狮子山、鲤鱼山、官山、梧桐弄组成了低水坝群；东南方向塘山是平行于大遮山走向的山前长堤。

良渚堤坝的结构类似于现代的黏土心墙坝，坝体外侧用黄色黏土作为坝壳，坝体内则填筑淤泥和草裹泥。"草裹泥"就是用南荻等植物茎秆包裹，再用竹条进行绑扎固定的块状泥土，类似于现代抗洪时使用的编织袋装土。将一块块的"草裹泥"纵横交叠，堆筑成堤，既能提高坝体的强度，也能加快固结过程，在运输和垒筑时更是便于分工协作。

良渚人还巧妙地利用自然环境的特征，对城外山体进行改造，从新发掘的这段石岭头坝体能清楚地看到长堤的迎水面处采用了石砌结构，因此能承受住剧烈水流的冲击。

汛期时山洪被长堤引向西侧，高低二级水坝形成的库区中，可避免良渚城所在的山间平原出现洪灾；伏旱期时库区又可以向城里城外供水，解决农田灌溉和生活用水的需求，使都城的交通命脉保持畅通。这是同时期世界上规模最大的水坝系统，蓄水量达4600余万立方米，相当于3个西湖。

4900年前那场持续多日的大雨终于停息，精密设计的水坝系统化解了暴雨和山洪的危险。云破天青，阳光漫照在平畴绿野、丘陇山林、民居王宫间，城内外又恢复了往日的繁忙。这是一座无与伦比的水上都城，9座城门中有8座均为水门，城内的51条河道和内外护城河绝大多数都是由人工开挖而成。城内外水系相互连通形成了完整的水上交通网络。钟家港古河道也是良渚城这3万多米的人工水道中的一段，在其两岸的部分台地边缘，还发现有保存良好的木构护岸遗迹。河道内的生活废弃垃圾不多，说明当时为保证水上交通进行了较严格的管控。

研究者溯溪而上，探索建城材料的来源。满目青翠间，与良渚人的竹筏隔着5000年的时光错身而过。生长在天目山中海拔四五百米高处的木材，经由水路被运送到良渚城内的钟家港岸边。在这里，麻栎树干被加工成圆木驯，覃树干被加工成方木，它们将被运至莫角山上成为架设神殿的梁柱。

良诸城墙的总长度将近6公里，宽度平均为50米左右。为了增强稳定性，避免地下水对墙体的渗透，人们在纯净黄土夯筑的墙体底部铺垫石块作为基础，这些石料几乎都来自城外西面，东南面的山中。

考古工作者凝视着残余的墙体，分辨着地层，观察着制作技术。城墙之内仿佛站着古时的一群人，他们也曾望向没有竣工的墙体，望向未知的历史那头．.

古城中央那高达10米，面积约30万平方米的莫角山是依托自然体堆铸成的巨型土台。高台之上的三土墩建筑群是大大小小的宫殿，土墩之间7万平方米的夯柱沙土广场上也曾矗立着一些礼仪性的建筑。

良渚的日常居住人口可能是两到三万。莫角山上居住着王和贵族们，宫殿区外围人工堆筑的丘阜上分布着大大小小的作坊，从事玉器、漆器等高端手工业生产的人群居住于此。厚重的城墙形成环抱之势，东北之雉山和西南之凤山各倚一角，状如后世瞭敌警备的角楼。城墙以外自北向南的长条形区域内，数座人工堆筑的土垄上分布着小型聚落。

莫角山南的池中寺遗址出土了相当于 19.7万公斤稻米的碳化稻谷堆积，足够5,000人一季度的消耗。对稻谷的锶同位素分析结果显示，他们来自良渚城以外的不同地方。城内出土的猪骨同样也有多个来源，这表明良渚古城的统治者们可以在大范围内调配基本的食物资源。当原先的城址规模渐趋饱和后，人们扩建了外郭城，历经数百年的开发和建设，良渚城的总面积竟致630万平方米之巨，是那个时代当之无愧的超级"都市"。

良渚古城和水利系统工程的土方量总计约1100万立方米。假使1万人每日风雨无阻的连续施工，也需要11年时间才能完成。

是什么样的力量能聚集这么大范围的人群和资源，建立和维续着如此辉煌的都市？是以信仰为纽带组织起来的权利。

如果我们必须通过空中的上帝视角，才能看清楚良渚人规划的古城和水利之宏大，那么我们需要放大镜，才能观察到良渚信仰的精微与幽深。

玉琮是良渚人至高的通神法器，这件玉丛重6.5公斤，是目前已知的良渚玉琮之首。它出自反山墓地12号墓，墓主应是某代良渚王。玉琮四面的直槽内都以减地微雕而成，一上一下两个神人兽面纹，只有3厘米高，4厘米宽。其间雕琢极为细密，一毫米内就刻画了四五道线条；玉琮的四角则刻有上下分节的神人和兽面；两侧是抽象的侧身飞鸟纹。

这是良渚人心中的神明，他头戴"介"字形羽冠，耸肩、平臂、屈肘叉腰。腰部的位置有一圆目、阔、带獠牙的兽面，双足似爪。这神人兽面像便是良渚的神徽。

神徽中神人兽面的组合表现变化多样，但万变却又不离其宗。它不仅在玉器上被大量表现，也见于其他材质的载体，如象牙器、漆器、陶器等。

早在良渚建成之前，瑶山、北村、官井头等墓地就已出土了大量的玉器，围绕这些玉器的使用，已经建立起了一套严格的规范，体现出明确的尊卑之别，男女之分。早期女性贵族的地位非常突出，墓中可发现一套套完整的玉璜和圆牌饰组合以及偶见但重要的艺术形象。

良渚的信仰和观念也已形成，龙首、人面、神兽面，成为在各类玉器上反复出现的母题。尽管刻划风格时有拙趣，也常常用到镂空工艺，但同良渚古国发展高峰出现的具象的"神徽"显然是同源同意。

刻纹玉器之于良渚人而言，不仅仅是社会权力的表现，更是精神世界的映照。在良渚的社会结构中鲜明的等级差异贯穿生死。

城外的汴家山上葬着平民，文家山上葬着较低等级的贵族，高级贵族先后葬在城内的姜家山和城外的瑶山、汇观山等高台墓地上。莫角山西侧的反山长垄就曾是数代王族的安葬之处。

考古工作者正在绘图记录着反山第20号墓的情景。这是反山墓地中出土玉器种类最齐全、数量最多的一座墓葬。半圆形饰、成组锥形器、三叉形器和冠状器构成了男性头部玉饰的制度化组合。那由300多颗玉管和玉珠所形成的串饰，并非墓主人生前的日常配饰，而是重大仪式中的特殊用器。墓内的4件玉琮、1件玉钺、24件石钺和43件玉璧，均具有象征性的含义。

琮是良渚信仰体系的载体，钺是男性贵族社会权利的物化表现，玉璧则是财富的象征。种种迹象表明这也是一座王墓，年代稍晚于12号墓。

透过专业的目光和笔触，散落的玉管连成珠串，繁复的玉饰各归其位，叠压的玉璧、石铖依着顺序逐层入葬。它让我们越过时间和尘土，同良渚人一起见证那场贵族的葬礼。

与20号墓大致同时的22号墓，其玉器组合却迥然有别。墓主的头颈部放置一组由12件玉管和1件刻纹半圆牌饰组成的串饰。她的胸腹部纵向放着一件玉璜，一组圆形龙纹牌式的组玉佩。这些事物与一同随葬的冠状器、琮式管、环镯、玉璧、玉鱼等玉器，共同彰显着一位女性墓主的权贵身份。

良渚权贵通过一整套玉礼器及其背后的礼仪系统，显示对神权的控制，进而延伸至对世俗王权、军权和财权的占有。

这套权力信仰体系在临平横山、青浦福泉山、吴县草鞋山、武进寺墩和无锡邱承墩等地的高等级墓葬中呈现出惊人的相似性。

在高度一致的信仰系统和文化认同中，各区域中心之间交流分配着象征权力与信仰的高等级玉器和图像。由此建立良渚古国间的政治和宗教联系，构建稳定的社会权利网络，创造了整个良渚文化区域内的共同秩序。毫无疑问，良渚遗址群为代表的社会组织已经进入早期国家的文明阶段。

距今 4,900 年前，深秋黎明的莫角山上，庄严的神殿浸染在霞光中。神殿广场之上已聚集了来自良渚古国各地的首领。一场盛大的祭祀仪式即将开始。

大巫师上前将玉琮交到良渚王手中。良渚王头戴玉锥形器冠，额束缝缀半圆形神人兽面纹玉牌的丝带，发髻上插着三叉形器与长玉管组成的玉饰，脖项间挂着玉管串饰。并肩立于他身旁的王后，髻上插着玉梳背，头上戴着半圆形玉牌与玉管串成的冠饰，颈间挂着玉璜和圆牌饰组成的繁复组佩，双腕戴着玉镯，行动间琳琅珑璁。

良渚王郑重地举起玉琮，那一刻玉琮上的神徽无比清楚而神圣，仿佛神的眼睛正注视着良渚古国的芸芸众生。

通过开创性的政治实践，长江下游的先民缔造了良渚古国。与此同时长江中游的先民们也在探索着区域性整合的道路。

最早铸造城墙的城头山，在距今5,300年前扩建为规整的圆形城池。高耸的城墙，宽深的护城河，这些不再只作为排涝防洪的设施，同时也成为了严密的军事防御屏障。

城头山扩建后不久，在16公里外一座新城出现，这就是鸡叫城。起初居民依照传统修建了环壕聚落，直到距今5000年前随着聚落规模的扩大，人们决定拓展居住空间。他们于是填平环壕，在更外围开挖护城河，挖出的土堆在护城河内侧，便筑起宽阔的城垣。

从普通环壕聚落演变为城池，不仅意味着生活空间的拓展和改造，更意味着通过城壕等大型公共设施的修建，发展了在大规模人群内部进行协作的机制。

随着鸡叫城的发展壮大，周边部落纷纷归附，城外不断出现新的居民点，这些蜂拥而至的人群造成了粮食紧缺的局面。城中的长老们为此争论不休，有人提出应强制驱赶一部分城外的居民点，有人提议抢占青河城和走马岭城的土地，也有人主张扩大耕田面积并将粮食进行集中分配。但旁人均觉得此法见效太慢，于是这声音便被淹没在一片喧扰间。

原野的高阜上，城郭剪影浮现，青年首领走向那位略显落寞的长老。这片肥沃土地是先祖贻留的宝贵遗产，值得为之殚心竭虑，矢志不渝。

不久后，在鸡叫城的周边，人们丈量土地，规划空间，开沟掘土。一条条相互平行的沟渠在大地上铺展开，所经之处便是新造的水田。

考古发掘揭露了城中一片距今4,800年左右的谷糠堆积层，这层谷糠的堆积厚度约15厘米，揭取的面积为80平方米，可大致推算出其代表的稻谷重量约为2.2万公斤。而这仅仅是本地稻作农业规模的冰山一角，长期持续的农田水利建设彻底改变了周围的大地景观。数百年后，鸡叫城及周边区域形成了以城墙为中心的三重环壕。第一重环壕之内是城的主体，护城河和第二重环壕之间为居住和生活区，二重环壕之外是稻田耕作区，一道道平行水渠贯通着环壕之间的水系。

随着社会生产和生活的空前繁荣，鸡叫城得以牢牢控制周边的几十处从属聚落，取代城头山成为整个澧阳平原的中心聚落。

如果说古国是一种文明形态，那鸡叫城就是从自身泥土和稻田里生长出来的农业文明。鸡叫城的社会组织行为开启了新的发展阶段，多组可能用于公共活动的大型建筑遗迹的出现就是明证。

其中63号房址的保存完好度最为令人吃惊，房屋由主体建筑和外围廊道组成。从东到西面阔5间，室内建筑面积420平方米，加上廊道总面积为630平方米左右。先民开挖基槽后，垫长木板以作基础，于木板一侧立柱，有些木板边还可见抬板时留下的绳索。木柱极为考究，以直径约0.5米的半圆形大木柱为主体，间以长方形小木柱，两侧均有斜穿孔。这座距今4,700年左右的史前大宅，间架有序，布局规整，是先民凭着仅有的石制、木制工具以及他们的双手和巧思建造起来的。中国木构建筑传统的特色已然尽显其中。

城头山和鸡叫城虽然是澧阳平原上两个重要的中心聚落，但若从整个长江中游来看，它已远远无法与江汉平原上正迅速崛起的诸城相匹敌。

汉江北岸的大洪山南麓，屈家岭文化的早期城址聚落群在5,300年前便初具规模。直到500年后那始于青萍之末的溪洞终于酿成壮观的风潮。

石家河城扩展成占地达120万平方米的大城，周边的直接控制区域更是足有8平方公里之广。以石家河城为首，江汉之间的聚落进入新一轮的发展高峰。残留于地面上的西城墙曾经在护城河中流淌的水，还能让我们依稀看见四五千年前石家河古城的巍峨。

在石家河城引领下的各个次级城址，如京山屈家岭、应城门板湾、安陆王谷溜、沙洋城河、天门笑城和孝感叶家庙，无不规模可观。

在长江中游一个强大的古国联盟正在形成，要凝聚这样的大型政治实体与社群，必然需要社会意识形态的革新与再造。

石家河城址中央台地上的谭家岭曾发现过大型建筑的遗存，或许是统治者曾经的居处。西北角的邓家湾是一处祭祀场所，在石家河城未起时就存在，城起之后仪式不断。

陶人偶和陶塑动物出于低洼之处和祭祀坑之内，足有上万件，或许象征着献祭之牲。人偶抱着鱼跪坐于地，如虔诚奉献。这样的大型祭祀场所不止一处，城外西侧的印信台上有5个人工堆筑的台基。平整过的地面上密布着祭祀坑，坑内埋藏着套接在一起的大口缸，多的可达25千，延续10米之长。

刻符陶缸起源于遥远的山东地区，大汶口人通常将其用于随葬。石家河在远距离文化交流中并未直接因袭，而是结合本地的实际需求进行了改造，使之成为凝聚本邦各区域政治实体的仪式工具。

站在石家河城内西南隅的三房湾遗址旁，很难不被这密集堆放的红陶杯震惊。它们厚达数米，总数可以百万计。这些陶杯的形态不尽相同，以斜腹杯、斜弧腹杯、喇叭口卷沿杯为主。同一地点分布的黄土坑、烧土面和洗泥池，显示这里曾是专门烧制红陶杯的窑厂，且沿用了数百年之久。这大规模的专业生产线折射出石家河社会的资源掌控能力，工艺制作实力以及举办大型仪式的频繁需求。

4300年前的石家河城，声势浩大的盟会始于一场庄重的祭祀仪式。长江中游诸城的首领们跟随石家河君长步入会场。君长将红陶杯中的酒酶于地面。

巫师将桃人偶和陶塑动物放于预先挖好的祭祀坑前，以敬奉祖先与神灵。首领们郑重捧起刻有本城符号的套缸，与他城的套缸口底相接，象征着参盟者紧密团结，休戚与共。

长江中游的先民们在江汉沅澧间扎根，农作已数千年。由于人口增多，交流频密且河湖交错，洪涝频繁，一个个规模庞大人口集中的古城在湖沼边缘的平原地带拔地而起，呈串联之势。通过仪式性的联盟、公共性的协作，层级城市体系的构建，他们凝聚起共识，认同于彼此。用千百年的时间共同耕耘一片土地，用千万双手共同搭建一座座城，直到缔结一个稳定的区域性政治联盟，造就一方土地的强盛。

盛夏的岱宗之北，古济水之滨，焦家聚落周围，流水潺湲，林木葱郁，座座村舍升起炊烟袅袅。

大概在距今5000年前后的大汶口文化中期，焦家先民开壕筑城，成为目前考古所见海岱地区年代最早的城。有别于良渚古国中贵族与平民异地而葬的制度，焦家遗址目前发现的近400座墓葬，大小墓葬共处于一个墓地，按等级分区，排列相对井然有序。

152号墓葬的墓主是一位老年男性，墓葬占地近12平方米，使用了两椁一棺作为葬具，是目前发现最早的具有三重棺椁的墓葬。内层棺中墓主头部有束发的骨梳和耳坠，手臂佩戴着玉环臂，手指上套有玉指环，下腹部放置象征军权和王权的多孔玉刀、装柄玉钺，身侧和小腿部位放置着作为法器或响器使用的龟甲器。

棺椁之间随葬了27件陶器，有陶鼎、陶豆等炊煮盛器，还有成组的背壶和高柄杯等水具饮器。功能各异的陶器被有序地陈列在独属于死者的空间内。

对于焦家贵族而言，死亡并非终结而是延续，他们希望通过这些沉默的器物展现出繁复讲究的宴飨仪轨，并将那雍雍肃肃的隆重，从生前一直贯穿至地下。

与这些复杂而奢华的大中型墓形成强烈对比的则是那些几近一无所有的小墓。焦家遗址所显露的阶层分化现象意味着大汶口文化中晚期的鲁北地区，已经开启了社会复杂化的进程。由焦家遗址向南约200多公里，鲁中南山地的西南麓是山区至平原的过渡地带。这里背倚青岭，沃野千里，古泗水静静流淌。

距今 5000年前后，岗上城之始建，最终发展成为海岱地区这一时期最大的城，统领周边60余处聚落。

在岗上遗址的大汶口文化晚期墓地中，考古工作者发现了成排排列的墓葬，最为特殊的是一个合葬大墓，面积达10.56平方米，四位墓主均为男性，自北向南年龄递减。葬具是并排的三联棺，带有头箱和边箱。

这座墓随葬的陶器数量达到300多件，以壶、杯、豆等宴饮器为主，而觚和大口尊等礼器则是墓主显要身份的象征。墓中4人均随葬玉钺，多放置于腰胯附近，但这些玉钺在形制、大小、玉料材质等方面判然有别，又体现出同一阶层内不同墓主间也存在明显的等级差异。女性大墓的体量和随葬品与男性大墓相当，但葬具为一棺，头部有骨梳，颈部放着一串绿松石串饰，随葬玉器多为饰品。二层台西南角放置大量猪上颌骨、猪蹄，是墓主生前财富的重要表现形式。

大汶口各区域中心的大墓虽各具特色，但其统一且制度化的一面已经成型，由此开启了中国特有的"礼制"和葬仪。

例如营建大型墓穴和层层相套的多重棺椁，又如以配玉和玉兵等玉器的使用，实现对个人身份的认定与表达。而在棺椁间依次放入成组成套建制规整的陶器组合，则是具有公共性的仪式活动和社会地位的展现。

在后世儒家的礼治中将社会秩序和伦理以物化的形式体现在丧葬仪式活动中，对自天子至庶民各等级死后的棺椁层数，随葬品的种类数量和摆放方式等都有明确的规定。

而 5000年前海岱山岳各处的大汶口高等级墓葬已经拥有这样一套具备严密秩序的丧葬礼仪。在死后长眠处，大汶口文化人群表现出对于统一社会规则和理念的强烈依从性。以死观生，生前的组织严密性更是可想而知。

黄河流域关于古国的探索并未止步于早期礼制的构建。一场更为深刻的社会变革即将掀开新章，被后来的考古学家们称之为"龙山"的崭新时代喷薄欲出。

4500年前一个盛夏的晌午，日照海滨的阳光灼热而刺眼。丹土城的大殿中贵族之间正发生一场激烈的争吵，穿越大汶口文化晚期此时的丹土已有百年建城史。社会分层固化，年长的贵族们依靠血缘关系结成的权力网络大大的限制了酋长的权力，也压制了反对者的声音。与丹土城毗邻的尧王城，无论在选址还是规模上都不可小觑。团结奋进的气势更是让丹土年轻人羡慕，或许正是鲁东南地区这些早期复杂社会的发展与竞争，促进了大汶口文化迈向龙山时代的社会变革。

要获得更为充分的社会掌控权和财产权，自龙山时代开始的分家在所难免。酋长终于决定带领着年轻人去往丹土以东4公里外的两城镇，踌躇满志地要在这黄海之滨开辟一片新的天地。山东龙山文化从大汶口文化脱胎出来的过程中，聚落数量快速增加，在一些地区呈现几何量级的增长，聚落规模也在快速扩张。

鲁北的寿光边线王、临淄桐林、邹平丁公、章丘城子崖和鲁西的景阳冈，鲁东南的五莲丹土，日照两城镇。尧王城和鲁南苏北的藤花落都是所在区域的核心城址。串起这两个城市群的应该是当时的主要交通路线。

龙山先民不再像大汶口时代那般，营建以血缘原则安排的大规模族墓地，这或许又意味着城内的居民们已处于杂居状态，最初的城市生活开始了。

黄海之滨的两城镇在4,000多年前的改革者的推动下，成为现在所能见到面积最大的龙山文化遗址之一。大汶口人在手工业经济领域的优势已被龙山县民继承，并发展出初级专业化和社会化的标准。

两城镇便是一处以制造业见长的中心聚落，作为当时最发达的工业产品，陶器有着巨大的产量和复杂的器类，其中蛋壳陶就代表了当时最顶尖的制陶技术。

色墨如漆，光泽莹润，器体轻薄，质硬如瓷这便是龙山文化的薄胎黑陶杯，因壁薄如蛋壳而得名"蛋壳陶"。一件这样的杯子，最薄处的杯口仅有0.2毫米，重量不过93克，即便以现代工艺也难以完全复制。杯柄中部鼓出部位中空并装饰细密的镂孔，其内还放置一粒陶丸，将杯子拿在手中晃动时，陶丸碰撞笼壁会发出清脆的响声。它的造型挺拔秀丽，脆弱又精致，如用墨色晕染的一株含苞欲放的百合，静静伫立在时光的薄暮中……

核心城址依赖达的制陶业等手工业经济成为区域的中心。周边的小型城市和聚落对其产品存在需求，于是产生了频繁的交换，乃至商贸活动。在对这些经济活动进行管理时，便诞生了制度与权利。

随着城市社会的产生，社会分化进一步加剧。玉器是当时的高端产品，目前发现的大部分玉器，比如多孔刀、钺、璇玑、璧、琮、璋、圭等礼器，基本都出于大型的中心聚落或是祭祀遗迹中，应该属于专供上层或仪式使用的奢侈品。

而位于山东省中部地区的临胸西珠峰遗址发现了目前考古所见，龙山时代规格最高的大墓。墓内所处的两只玉笄精美绝伦，那种处于整个社会最顶端，可能具有"王"的身份的阶层已经产生。

在邹平丁公遗址出土的一片大平底盆底部残片上，俨然可见线条流畅，排列规律的五行 11个符号，刻写有一定章法，为探讨原始文字的产生提供了珍贵的实物资料。

古国时代列城林立的景观正与先秦文献中上古时代天下万国的记忆合节。

东部地区良渚古国以信仰为基石，以权力为利器一统长江下游，已然迈入文明的门槛，成为东亚最早的区域性国家。

在良渚人以神权维系长江下游古国时，黄河下游的海岱地区却走向了更为世俗化的社会。他们自现实生活中发展出社会治理的智慧，较之鬼神却更崇人力，继承并发展更为严密的礼制实践，成为了"礼出东方"的强大力量。

石家河扎根于传统，在长江中游建立起强大的连城网络，形成千里同风的联盟。大约在4,100年前来自北方的力量南下，改变了石家河文化的面貌，长江中游也由此纳入早期中国的版图。后石家河文化中的玉器和图像元素均在日后的跨区域交流中扮演重要的角色。

除此之外所有地区都在不同程度地参与着此后作为整体的华夏文明的构建。正是他们的发展造就了广域中国的多元基础，又为王朝中国的延绵发展提供了生态、文化、人口和资源等各方面腾挪的空间。

黄河之水奔腾而下，中原即将成为新的时代引领者，下一场历史变革就要到来。

\chapter{第五集 \quad 择中}

\section{正文}


黄河在内蒙古高原和黄土高原之间形成的大“几”字形，是造化之力在神州大地书写出的最壮阔一笔。紧邻这个“几”字形大弯的东侧，是吕梁山脉、太行山脉和它们之间的一连串狭长盆地。距今4300年前，在其中的临汾盆地出现了一个曾辉煌一时，却因岁月久远而湮没无闻的古国，它的都邑就在今天山西襄汾陶寺村一带，研究者遂以地为名，称之为陶寺。雾中，考古工作者们忙碌却有序地清理着陶寺遗址的一角。随着覆盖其上的泥土逐渐被揭去，这里赫然显露一场灾难的证据，4000年前的一位年轻贵族，侥幸从这危机重重中逃脱，他宁愿投身那似乎终古不见天日的密林中，也不愿再面对那曾经繁盛的都邑里发生的一切。

对于整个陶寺都邑中的贵族来说，公元前2000年，是动荡与离丧的开端。曾经守卫森严的夯土城墙被破坏、被废弃，最终也未能护住城中的安宁。耗费无数人与物才建就的宫殿已被粗暴的拆毁，昔日宫室中，那戳印精美图案或绘制蓝彩的白灰墙，曾是令造访者炫目的华丽，已成为散落在泥土中的建筑垃圾。政令所出的宫殿，被改造成制造骨器和石器的手工作坊。最令人触目惊心的，是作坊旁的垃圾沟内，那散乱混杂在骨石废料中的人骨。就连早已作古的人们也没能逃过噩运，陶寺的王墓和贵族墓中，大多有直捣中央棺室的扰坑，死者的头颅、碎骨被随意扬弃于墓内各处。种种迹象表明，这非是为求财而来的盗墓，而是一场宣泄式的毁墓虐尸行为。

这场祸及先人和生者的暴乱因何而起？沉默的陶寺遗址并未告诉我们答案。日夜奔流的汾河边上，陶寺之邑曾经强盛一时。宽近八米的城墙维护着280万平方米的土地，守城的卫兵需要快走一小时，才能完成巡城一周的任务。当时的黄河流域和长江流域，一座座城邑拔地而起，风云竞涌，古国林立，俨然正是传说中那万邦并伫，英雄辈出的景象。然而放眼彼时之世间，在规模上能比肩陶寺城址的也并不多。这样庞大的工程规模和城区空间，足能展露出陶寺社会的综合实力，不难想见城内的人口数量之巨，统治者的社会动员力之强。22号墓是陶寺的大墓之一，在经历过发冢之劫后，排场依旧惊人。

墓穴至今仍深7米，陡直的墓壁上装饰着五周平行的手抹草拌泥宽带，或许是对生前居室装修风格的模仿。这壮观的大墓里盛装着的，不只是陶寺贵族曾经的荣光和尊严，还有属于那个黄金时代的迷梦。在曾经的安葬仪式上，人们把一具船形木棺安放在深深的墓底正中。木棺由一根整木挖凿而成，通体红彩，鲜艳夺目，上覆布置棺罩。一个完整的公猪下颌骨，被高高的放置在墓主人头端的墓壁中央。以其为中轴，这面墓壁下方两侧，各倒置着三件彩漆柄的玉石钺和戚，兵器的陈设显现出仪仗的威风。棺木的左侧与墓壁间，排列着4柄青石刀和7块木案板，石刀下的骨骸和朽灰表明，下葬时案板上还盛放着鲜猪肉。

墓主人足端一侧，摆放着10头对半劈开的猪，墓的四壁底部掏有11个壁龛，里面放置精美的玉器、漆器和彩绘陶器。纵然迄今探明的陶寺墓葬已有数千座之多，其中九成都是仅能容身、空无一物的小墓。像22号墓这样规模的大墓实属少数，这反映出陶寺古国的金字塔式等级结构。王族高居塔尖，在其下依次是贵族、平民、赤贫乃至非自由人，社会分化相当严重。权贵们的大墓即便在幽暗的地下，也极尽所能的炫耀着自己对于珍稀资源的掌握，彰显着非比寻常的身份。在大墓中，屡屡可见融汇四方之风的华美器具，玉琮和玉璧宛然可见黄河上游齐家文化的印记。

镂空兽面，散发着长江中游后石家河文化的神秘气息，鼍鼓之上似能听到来自海岱地区礼乐文化的节拍；陶器上的彩绘图案，甚至依稀有东北地区文化之遗风。当然，繁盛如陶寺，又岂会满足于仅以远方殊物表现自己的权威？他们早已掀起新的文化浪潮，建构了自己独特的制度逻辑。在陶寺的几座大墓中，考古工作者发现了一种从未见过的彩绘陶盘，每墓仅有一件，全都统一放置在墓主右侧偏上部。这样的现象，表明陶盘之珍罕，它们或许是只有特定人群才能拥有的，具有特殊含义的仪式用物。这些陶盘通体覆有一层黑色磨光陶衣，盘内以朱红色颜料绘制一条盘曲的动物纹。

此动物，长躯麟身、方首圆目、巨口长舌、无角无爪，似蛇非蛇，似鳄非鳄，显然并非尘世中物，而是陶寺人揉合了多种动物的特点，创造出来的灵兽。它的形态与中国人所熟悉的龙十分相似。在其中一座随葬陶龙盘的大墓内，墓主足端还放置有一组鼍鼓和石磬。鼍，是扬子鳄的古称，鼍鼓即蒙以鳄鱼皮的木鼓，鼓腔呈竖立筒状，高1米，以树干挖制而成，外壁通体施彩绘。石磬是打击乐器，后世的磬大小相递，成套使用，是以又称编磬。陶寺之磬却仅为单件，因此称为“特磬”。另有数座大墓也在墓主足端放置了这样的特磬与鼍鼓或陶鼓之组合，磬音清越，鼓声嘭嘭，应节而起。

陶寺的红铜铃是迄今所知中国历史上的第一件金属乐器，它的出现，预示着“金石之声”时代的即将到来。金声即铜声，铜与玉石和鸣，正是夏商周三代音乐文明的重要内涵。鼍鼓声沉，石磬声清，铜铃声在更高的频率上带来了另一个声部、另一种音色的和鸣。音乐也呼应着世间烂漫生长的万物，众声皆备，这是初萌于陶寺的合乐之声。铜铃虽小，制作工艺却比大件的实心工具和兵器要更复杂，难度也更大。铸造普通的实心铜质工具仅用单范便可以完成，而铜铃的腹腔中空，铸造时必须使用复合范，也就是多于两块以上的范。因此，陶寺的红铜铃是迄今所知年代最早的，以复合陶范铸造的完整铜器，这种铸造技术也正是后来中国青铜时代的标志性工艺。

陶寺文化所处的龙山时代，在时间上正接续在青铜时代之前。陶寺所处的晋西南，在空间上本就是大中原的组成部分，而且恰恰位于中原与北方交流的重要孔道上。虽然陶寺的铜器是使用天然铜铸造的红铜器，有别于后来的青铜合金，但从时空两重因素推测，陶寺一定曾在中原青铜文明的崛起过程中，扮演过重要的角色。四方的资源技术汇集到陶寺，也造就了这座古城规划有度的布局，城内功能分区明确，宫殿、贵族生活区、祭祀场所、手工业作坊、仓储区和不同规模、规格的墓地一应俱全，是这一时期大中原地区，“城市要素”最为完备的巨型都邑。传说与现实的交织引人遐想，莫非尧曾在此驻足？

西倚黄河，东望太行的陶寺，一度曾如此辉煌。这里社会内部等级划分明确，礼乐器群逐渐创备，礼制作为建立新秩序的重要支柱，已经初行于世，耀熠着即将到来的青铜时代。而与兼收并蓄的广阔胸襟和都邑布局的宏大气势，形成鲜明对比的是，考古调查表明，陶寺文化的聚落，基本上只分布于其都邑所在的临汾盆地周边，山河表里，是可凭恃的天险，却也是可屏翳的阻隔。大邑小国的陶寺，数以万计的人口聚集于辐辏之地，至中晚期，不同阶层间又发生着激烈的斗争。距今4000年后，陶寺文化和陶寺都城开始走向衰亡。

石峁，是黄河的“几”字形大湾所包围的黄土高原上生长起来的一个古国。它处在一片被称为“半月形地带”的内弧一侧。半月形地带是中国大陆自东北、西北而西南的一道美丽弧线，形似一枚嵌在东半部的第一、二级阶地外围的新月，月牙朝向孕育古老华夏文明的黄河中下游和长江中下游。这道半月形弧线上，层峦叠嶂，连绵而跌宕。自东北的大兴安岭、内蒙古的阴山、宁夏的贺兰山，直至川西、滇北的横断山。这片半月形地带的存在，最初是由于中原农耕人群向周边的扩散而形成的。因为人口膨胀，许多人需要向外寻找适宜生存的地方。他们通过智慧、双手以及人彼此间的协作，在与平原迥异的生态环境中也发展出了新的生存策略。他们将黄河流域和长江流域的粟黍、稻作农耕技术带到这些山地里春耘秋收，形成星点散落其间的各个定居社会。

而这个地带，也成为了与欧亚大陆其他文化产生碰撞、融合的前沿阵地。当黄河中下游、长江中下游的大部分地区先后进入古国形态时，地处更外围的半月形地带也迎向了时代的风潮。在毛乌素沙漠东南缘，黄土高原北端的黄河西岸，秃尾河畔的黄土梁峁和剥蚀山丘上，海拔千米之处屹立起一座石峁古城。这座始建于4300年前，面积达400万平方米的巨型城址，是那个时代陕北、晋西至蒙中地区的中心。由于在地面以上留存了相当高的一部分，曾一度被后人误认为是战国时期的秦长城附属边塞遗址。万万没想到，一座史前的城址竟能保存的如此完好，秋日的夕阳，只留一抹染红石峁城。

这石城之壮观，为史前东亚地区所罕见，其核心区是高大巍峨的“皇城台”，这是座四周以石块包砌成阶状护坡的台城，顶部十分宽敞，足有8万余平方米。石峁城有内外二重城墙，内城将“皇城台”包围其中，城墙依山势大致呈东北、西南向分布，面积约210万平方米，外城利用内城东南部墙体继续向东南方向扩筑一道弧形石墙，形成封闭空间，城内面积约190万平方米。石峁内外城，各发现四座城门，多数城门地表仍可见墩台、瓮城、角楼等各类附属建筑，甚至还有疑似后世的“马面”结构，这些多是适于布设驻防力量的设施，将石峁城层层卫护，防御森严。

距今4000年前的深秋，夜幕覆于千万道坼裂的黄土梁峁之上，天地间一点点微光徐徐闪烁，火把之间，石峁首领正缓慢地、庄重地登上皇城台，登得越高，就离苍穹越近。明月孤悬，台城高矗，月光、人影与火光摇曳交映，流动过立于地面的神面纹圆形石柱，流动过嵌在石墙中的玉器，以及那些人面、动物和神面的石雕。环眼神面，是源自江汉之间石家河古城的庄严肃杀，人射马，是来自辽阔草原的生机勃发，两条龙盘曲如环，身披鳞甲，充满着莫测的力量。他伸手摩挲着石雕龙身，神龙犹如被唤醒，在暗夜的微光里闪耀着灵动的生命感。这些石雕散落着嵌于石墙间，有的甚至被倒置，但丝毫不影响那粗犷线条凿刻出的原始美感，石墙缝间，还常有故意嵌入的玉制刀钺，扁薄锋利的带刃玉器原是古国时代各区域首领下葬时，显示军权、王权的身份标志。

在石峁，它们却被广泛用于皇城台的营建、贵族宫殿的奠基等公共工程中。藏玉于墙，这是石峁人独特的习俗，玉所具备的透明、坚韧、锋利的特质，仿如一种神力，同石雕图像一起守卫着这座城，给予特别的护佑。首领持玉璋向天祝祷，群山之间，他听见秃尾河汨汨流淌的声音，听见虫鸣和风呼啸的声音，他对这片土地如此熟悉，即便身处光线微弱的黑夜里，也知道月色底下那山谷之间的小型石城——都在何方。往东南，渡过黄河，翻越吕梁山，便是大邑陶寺，再往东往南是更广阔的中原、海岱和江汉。中原，是石峁人的根脉，仰韶时代，祖辈们从太行山以东和关中地区北上移民至此，带来了故土的文化，也在新的家园生长出新的传统。

从石峁往北，是辽阔的欧亚草原，石峁人从那里引进了绵羊、山羊和黄牛，过上了半农半牧的生活。他们还从那里引进了砷铜合金或锡铜合金的青铜刀，获得了坚利的工具，石峁的一些石雕也能在广义阿尔泰地区的石雕形象中找到母本。但4000年前的那个夜晚，让石峁首领陷入沉思的，并不是自何而来的问题，而是往何处去的答案。从对古气候的研究中，我们知道，当时的全球气候正值又一轮周期性变化的阶段，进入了干冷时期。这轮气候变化带来的生存压力，尤其被处于农牧交界地带的石峁集团深刻感知。环境的恶化、人口的增长、资源的紧张，无不是随时可能压垮北方脆弱生态的最后稻草。各聚落之间，更是为了生存而征伐不休。

某个聚落的人家里，孩子因为饥饿而哭泣难以入眠，怀抱孩童的母亲耐心的哄劝，父亲在旁叹了声气，望向窗外的浓稠夜色，这一刻，皇城台上的石峁首领也正望向空荡荡的天地间，忧思重重。在危机时刻，趋利避害，凝聚族群，自是当务之急。但他也决定听从智者们的建议，向外去寻找更多的资源。何去何从，需祈祷上天，敬告祖先！他将口簧含入嘴里，一阵神秘之声响起，这声音震荡着山谷的夜色，仿若以此，达成与祖先神祇的某种神秘共识，南下，是唯一的选择。他相信漫漫长夜之后，石峁城将在黎明中复苏。

龙山文化晚期，征战、对峙、暴力、冲突在中原和北方地区成为常态，活跃于史前中国大地异彩纷呈的区域性古国先后走向衰落，江汉平原和江南地区人烟稀少，山东龙山文化逐渐式微。在陶寺湮灭之后，石峁也开始衰落，中原其他城池纷纷退出历史舞台。而时代的战车，将由新的王者驾驭。距今3700多年前，洛阳盆地北面，邙山的一处高岗之上，夏王少康面南而立，洛河经山下浩浩东去，南面，伊水两岸双峰耸立，如同门阙；西面，晨雾尽头是狭长的函谷，直通关中盆地。根据古文献记载，夏是中国第一个王朝，是王权世袭的开端，开辟九州的禹，为天下苍生栉风沐雨，披荆斩棘。此后的数代王者，有奋发与开拓之主，也有失德和流亡之君。几十年前，先王太康在一次田猎活动中，被东夷的后裔突袭，下落不明。太康失国之后，他的五个兄弟辗转流亡，在洛水边悲吟《五子之歌》。

“明明我祖，万邦之君。有典有则，贻厥子孙。民可近，不可下。民惟邦本，本固邦宁。”如今，东夷之人已经被降服，少康夙夜匪懈，终成复兴大业，他重回洛阳盆地，定都斟鄩。位于中原腹地的河南偃师二里头遗址，历经半个多世纪的田野考古工作，让人们看到了传说中夏都斟鄩的影子。它的现存面积达300万平方米，周边还环绕多个区域中心、次级中心及众多更小聚落，组成了金字塔式的聚落等级结构和众星捧月式的分布格局。这是一个构思宏大、规划缜密、布局严整、等级有序的超级都邑。文献只记录了它的地理方位，考古发现终于展露出它动人的细节。

二里头都城的中心，两条南北向和两条东西向道路纵横交错，呈方正规整的“井”字形，构成中国最早的城市主干道网，在连接交通的同时，又分割出不同的功能区。居中的是大中型夯土建筑基址集中的宫殿区；官营手工业作坊区和祭祀区，分别位于宫殿区的南、北两侧；贵族居住和墓葬区拱卫在宫殿区周围；再向外的部分，则是一般性居住活动区域，择天下之中而立国，择国之中而立宫，择宫之中而立庙。二里头都城所处天下之中的位置，体现出早期国家等级分明的社会结构和秩序井然的统治格局。而九宫格式的都城规划布局，又与《禹贡》以山川划分九州的“天下观”颇为契合，足以显示二里头王者在都邑建设中“辨方正位、体国经野”的政治抱负。

中条山谷中，王朝的矿工们摸索行进，他们的使命是在这群山中找到新的铜矿。那些往来于运城盐湖和二里头王都之间的盐队，常常需要翻越中条山，其中有人曾在连日雨后目睹远处山石上绽开绿花的奇景。矿工们知道，这或许正是铜矿存在的标志，是以扎入深山四处找寻。此前已有入山寻矿的人殒身于此，于是他们格外注意脚下的路。此时山雨突至，只好寻了一处岩厦暂避。连日的疲惫中，众人沉沉睡去，唯有一人在火堆边默默疗伤，一只翠蓝的蝴蝶蓦然飞起，使他心头油然生出一种奇异的异象，便不由自主地起身追逐。一路追寻而来的矿工，看见蝴蝶渐渐地落在对面的山崖上，那处山体因下雨而发生了滑坡，就在他的眼前，那石上赫然绽开一小丛一小丛铜绿之华，这正是他们要为王朝寻找的涉及国家命脉的战略资源。

山西绛县的西吴壁矿冶遗址，仍留存有二里头文化晚期的冶炼遗迹，铜矿原料便来自附近的中条山。这一时期的木炭窑，其功能是用来烧制炼铜使用的木炭，透过炼炉的残迹和炼渣的堆积，似乎就能看到铜矿石在炉中熔化，沉入炉底，化为纯铜。在欧亚草原上，青铜原本仅被用于制造武器、工具、装饰品等小件器物，且多为锻打成型。冶金术传入中原后，迅速发展出以复合范铸造技术为核心的铜容器生产业，新的资源技术与本身的礼制文化系统相结合，创造出独特的青铜文明。二里头遗址发现了迄今所知中国最早的青铜钺，而秉持斧钺之人就是有军事统帅权的首领。

铜爵是一种酒器，也是中国最早出现的青铜容器。青铜是合金，需要将铜与锡等矿石按照一定的比例调配才能冶炼而成。铜矿可在中条山获得，锡矿的主要来源却可能是大兴安岭南麓地区或是江西、湖南等长江以南地区。这意味着如要大规模冶炼青铜，就必须要具备远距离的矿料开采和运输的能力，加之铜爵造型复杂，涉及高精的铸造技术，这一切唯有具备比此前任何一个时代更强大的社会调动和资源控制能力，才有可能做到。二里头都邑的贵族们不仅垄断了青铜礼器的生产，也独占了青铜礼器的使用权。青铜，这种人类创造的物质材料，也被转化成为与社会秩序和等级制度相关联的政治资源。工匠们将日常器具的形体赋予它们。

食器、酒器、乐器、武器，一旦与青铜相关联，这些原本具有实用功能的器型，却也成为了一种超脱于日常的象征符号，具备了礼器、祭器的功能，也与权力身份相挂钩。当以烈火铸就的坚利金属，被塑成庄重敦厚的礼器，文雅莹润的玉石却制成了杀伐征戮的武器，这也是礼制观念中奇妙的二元对立。饮食器皿不再日常，成为青铜重器，干戈坚刃敛去杀气，化为玉石韫藉，这最早王朝的礼制观，自兹历千年而不衰。而那些被寄寓了特殊含义的礼器群，爵、钺、璋等器种，持续使用超过千年，进而成为中国古代社会政治文化的重要符号。除此之外，一种传说中的神兽更是成为流传至今的中国图腾，这便是二里头的龙，它的用工之巨、制作之精，前所罕见。

这条龙全长70厘米，眼、鼻皆嵌白玉，其余部位则由2000多片绿松石拼嵌而成，每片绿松石的厚度大多不足一毫米，长度也不过几毫米，原本应是粘嵌在木、革等材料之上，为了便于粘嵌，每个嵌片都制成楔形。嵌片形状几乎无一重复，虽历三千余年，色泽犹翠绿鲜艳。这是已发掘的二里头贵族墓葬中，位置最接近所在宫室中轴线的一座，墓葬虽只有两平方米，随葬品却有上百件，墓主人是一名30多岁的成年男子，大量海贝置于其颈部，头骨上方发现三件斗笠状白陶器，可能为头饰或冠饰的组件。他怀抱绿松石龙形器，手系铜铃，可能是一名巫师。二里头遗址的另几座墓葬中，墓主胸腹间与铜铃一道放置着镶嵌绿松石的铜牌饰，牌饰正面铸出寿面纹，再以数百枚细小的绿松石片细密镶嵌，以大颗的绿松石珠点睛，赫然勾画出一幅神兽形象，这或许正是龙形器的简化和抽象表现。

碧龙飞升，铜铃和鸣，神灵降临，亡灵升天。一如后来《诗经》中所记，周王祭祀宗庙时，“龙旗阳阳，和铃央央”的场景。神巫一足略跛，随着节奏，且行且舞，足迹画出大大的八字，这样的步伐被称为“禹步”，表现着大禹治水行山之艰难。新砦陶器盖上的龙首、陶寺陶盘上的蛇形蟠龙纹和石峁石雕上的龙蛇形象，在此刻合而为一，成为后世中国人精神世界里永恒的图腾。与之相呼应，东亚文明的复兴地区即黄河和长江流域，开始由多元化的古国文明走向一体化的王朝文明。传说舜帝统治天下时，南方的三苗不愿归顺，禹请求带兵征伐，却遭到了舜的拒绝。舜认为轻易动武乃是执政者德行不够的表现，于是偃武修政，执兵器而舞，三苗乃服。这传承至今的干戚之舞，展示的既是礼节仪式之肃穆，也是武力军容之威慑，不战而胜，这是夏王少康的追求。

在九州团聚的盛会上，自各方远道而来的嘉宾们奉上本地的珍惜物产，表达对华夏的敬意，热带海域的海贝贵重而绮丽，长江下游制作的印纹硬陶盉坚实精致，来自西北高原的是片状有刃的大型玉器，成群的黄牛与绵羊从北方草原被驱赶而来，而华夏用以回馈九州之敬意的，是那些象征王朝威仪的礼制和礼器文化。起源于龙山时代的玉石璋，以二里头都邑作为起点，向长江中上游、珠江三角洲一带传播，甚至到达今天的东南亚地区。盉、爵等二里头风格的陶礼器，北可至燕山以北，南可及长江中上游，西可达甘肃、青海一带。这正是最美好的时节，江山丽日，花草芳香，那田间、桑下，城中，殿前，欣欣向荣，一派蓬勃景象。

而少康也有他的职责，他关于这片国土的理想，“明明我祖，万邦之君。有典有则，贻厥子孙。民可近，不可下。民惟邦本，本固邦宁。”黄河九曲，夭矫如龙，九州共望，国之在中。在何以中国的追问中，二里头文化具有划时代的意义。中国在此时完成了从新石器时代到青铜时代的转变，开启了以铜礼器为核心，政治权力与伦理秩序合一的三代文明。二里头的崛起，是中国这个文化共同体内，第一次出现超强的政治文化核心，其影响所及超越良渚、陶寺和石峁等各个时期的代表，呈现融合集中的“王朝气象”。这一文化核心正处于周人所称的“中国”，影响波及四方，确立了中华文明向心力和凝聚力的基点，从此历经沧海桑田，永不离散。


\chapter{第六集 \quad 殷殇}

\section{正文}

传说中，女子简狄水中行浴，见玄鸟堕其卵，简狄取而吞之，因孕生契，契长大后辅佐大禹治水有功，受封于商，这以后便有了商族。商代甲骨文中，商王祖先亥的名号中冠以鸟形，这正是商人对玄鸟感生神话的崇信。从太行山中，到冀南、豫北平原，历代先祖经历几多艰险、几多征战，才由弱变强。作为玄鸟子孙，成汤得天命，自奋蹄，约公元前1600年，在山西南部的鸣条之战中，夏王帝桀被成汤率领的商军彻底击溃。而洛水之阳的斟鄩之地，辉煌一时的二里头都城已现衰败之象。大型礼仪和政治性工程因遭破坏而废弃，来自海岱地区的岳石文化，以及被认为与商人先祖有关的下七垣文化人群，大举进入城中，商人还在二里头都邑东北约6公里外建起了偃师商城。

城门狭窄，易守难攻，加之城内有大量储粮囷仓和储物府库，厉兵秣马的气息扑面而来，也许商军就屯驻在这座固若金汤的城池里，监视着不远处那二里头城中的夏遗民。与此同时，古济水之南，一座占地约25平方公里的巨大都城在今天的河南郑州拔地而起。它由内城和外郭组成，其间错落分布着宫城、居住区、作坊区和墓葬区。城垣经历了3000余年的风吹雨蚀、兵火燹灾，至今仍保留着高出地面4~5米、宽20~30米的规模。商人，将这座立国之都称为“亳”，成汤正是在这里统治初生的商王朝。商代早期，逢连年大旱，百姓陷于困厄之中，成汤决定亲赴桑林祈雨，他并未遵从杀人献祭的惯例，而是剪下自己的头发与指甲，放于柴堆之上，作为身体的一部分代己为祭。

兵力强盛，人心归附，商很快成为大邑。《诗经·殷武》中追忆道，“昔有成汤，自彼氐羌，莫敢不来享，莫敢不来王，曰商是常。”以郑州商城和偃师商城为核心，商王朝开始构建覆盖黄河、长江流域的庞大城池网络，一边沿用或改造夏人留下的城池，一边按照新的理念建设新城。中条山下，垣曲商城衔山抱水。长江北岸，盘龙城高大崔嵬。盘龙城的布局，与郑州偃师商城类似，主要由宫殿和手工业作坊区等构成，埋葬风俗和出土的青铜礼器也几乎与中原地区保持一致。盘龙城中的宫殿基址，虽然仅有郑州商城的十分之一，但在当时的长江地区，规格却无有出其右者。这座城，是商人在南方疆域的政治和经济中心，或许，也是控制南方珍贵资源的最佳中转地。

几乎与郑州商城建起同时，冶炼青铜的烈火，就已经照亮了郑州南关外的天空。此处的大型铸铜作坊，生产规模惊人，工艺精湛，不仅制造小件的青铜兵器和工具，还能铸造鼎、鬲、爵、觚等青铜礼器。商人先祖曾生活在北方，很早便接触到经草原传来的青铜冶铸术，后来，更是吸收了二里头的影响，最终发展出那个时代最尖端的大型青铜器铸造技术。此时，熔炉里的铜锡逐渐熔化为合金液体，司炉们敲碎浇道口的炉壁，金色的溶液奔流而出，注入浇铸孔，待其冷却后打碎外范，则大鼎可成。青铜，是商王朝至为重要的政治资源，也象征着这个正冉冉升起的新王朝那举世无双的实力。为了治理越来越广大的疆土，商王朝采取了内外服分而治之的统治形式。

所谓“内服”，就是商王直辖之地，委派官员治理；而“外服”，则是王室贵族的领地，或吸纳臣服的方国，拱卫在王畿四周，既是防御和缓冲地带，也是对外扩张的前线。早商王朝的征途似乎一路奏凯，然而，历史的伏笔此刻再现，改变了时局。建城约200年之后，郑州商城逐渐走向衰落，宫室少有人迹，作坊不事生产。考古工作者在城门附近，发现了这一时期的三处大型青铜器窖藏，在深达4~6米的窖坑中，掺杂朱砂或铺垫木板，藏有不同数量和种类的青铜礼器。在那个时代，少有一座城能够长久地作为都邑，王朝需要寻找新的、能维续庞大人口生存的地方。或许在离开旧都前，商人举行了最后的庄严祭祀，将青铜重器瘗埋于地下，为的是祈求祖先和神灵的护佑。

当郑州商城衰落之际，距其20公里外的索须河畔，却生长起一个大型的中心都邑——小双桥，城中那大型的礼仪性建筑、大规模的祭祀活动、工艺精湛的青铜重器，都彰显着这并非一座普通的城址，而可能是商王朝的第二座都城——隞。在这座城内，除了商人的活动痕迹之外，竟还出现了一些具有山东海岱地区文化特征的器物，城中宏伟的高台祭坛，曾举行过比郑州商城中规模更大的祭祀活动；数个祭祀坑内，都填埋着大量非正常死亡的人骨。此时的商人，似已忘却了当年成汤以己身代人牲的苦心。仲丁统治期间，迁都至隞，此时，商人与东夷的关系逐渐紧张，仲丁征战告捷后，带回大量东夷俘虏，以他们为人牲举行了献祭仪式。不久后，一支东夷军队却突袭隞都，杀死了大量来不及逃走的商人。

几番鏖战后，商王朝虽然夺回都城，却也元气大伤。更不利的是，商人王位的传承次序向不固定，极易造成不稳定的局面。仲丁以后，王室贵族们为了王位同室操戈，由此引发的混乱波及九代商王。由于各种因素的叠加，商人不得不放弃旧都，离开中原，踏上了在黄河两岸间来回迁徙的漫长路途。直到距今约3300年前，商王盘庚迁都至殷墟，局势才自此安定。大邑商，这气势恢宏的晚商王都，继续着“四方之极”的构建。西北方的高冈上，商王陵与2800多座祭祀坑俯临俗世。王都内，宫殿和宗庙被环护在洹河和壕沟共同组成的森严屏翳中。西南角，是祭祀王朝土地之神的“社”，东南角是祭祀商人远祖的“宗”。社与宗之北，是为外朝，典礼和仪式活动都在此处举行。

再北则是内寝，也就是商王的内宫，始建于武丁时期，沿用到帝乙、帝辛时期。寝宫的长案之上，摆着几件青铜盥洗之器，侍女向铜盂里放入香茅，再将青铜鼎中的热水舀入盂中，一时满室馨香。铜盘盛满清水，可以调节盂中的水温。商王武丁面色凝重，进行着占卜前的准备，侍女将布浸在盂中，拧至半干，为武丁精心擦拭身体。室内火烛闪烁，映照着武丁健壮的身躯，那一道道伤痕是无数次征战的痕迹，侍女手握陶磢抚摩皮肤，以去除污垢。武丁的心事，是妻子妇好分娩过程的吉凶。妇好是他的配偶，是孩子的母亲，也是朝堂上的股肱。她主持祭祀，领兵征伐，于国于家，不可或缺。邦畿千里，日理万机，妇好的事总是让武丁牵挂：

她每有牙痛、鼻患、腹痛或是偶感风寒，武丁都会郑重地占卜吉凶，祭祀消灾。数年前，他为将要临盆的妇好卜问吉日，31天后的甲寅日，妇好分娩，却偏偏不是求得的吉日，这一次，他不愿再有任何闪失。卸下戎装的妇好仍是个爱美的女子，她最喜欢玉，玉牛、玉熊、玉蛙、玉羊首、玉蝉、玉虎、玉人，这些精雕细琢的美玉，凝聚着天地的灵气。玉镯、玉串饰和玛瑙串饰，还有精致的骨笄，摆满了她的首饰盒。她常常对镜梳妆，即使对王室来说，铜镜也是稀罕之物，身份尊贵如妇好，方能拥有。她用玉梳精心梳理并挽起长发，再从众多的发饰中，挑出一只夔首骨笄，紧紧固定住发髻。今夜皓月当空，商王与王后，一起来至早已备好的祭场中进行祷告。

这钺，是妇好征战时号令军旅的统帅标志，也是仪式中驱祟的法器，钺未装柄时便重达9公斤。钺身上饰有侧身对立、张口咆哮的双虎，虎口相对间嵌有一颗人首，虎身后各饰一条卷尾之龙，虎爪下是巨大的獠牙。商人崇信神祇，通过这样的仪式，能得神灵与祖先的护佑，顺利产子，延续商祚。此刻的密室中央，炭火熊熊，贞人宾拿起卜骨，在炭火上轻轻烘烤，身旁的矮几上放着整治、钻凿和刻画甲骨的工具。贞人，不但是占卜活动的具体操作者，还承担着巫的职责。少数掌管占卜机构的卜官，是商王之侧的高级贵族。武丁左手扶住卜骨，右手持炭棒灼烧着钻槽，通红的棒头神秘地闪烁，不久后，裂缝在钻槽边绽开，发出脆响。

武丁屏息细看卜骨上裂开的兆纹，脸上渐渐露出喜色，这次占卜所得是大吉之兆。“丁酉卜，宾贞：妇好有受生？”“王占曰：吉！”“其有受生！”

宾拿起青铜刻刀，开始在龟甲上契刻。传说仓颉造出文字之时，天雨粟、鬼夜哭，天地间同为这亘古未有的创造而震撼。文字之于古人，如劈入混沌世界的雷电，惊醒并照亮了蛰伏的文明。自此以后，人类拥有了能够世代累积和叠加的智慧与情怀，万物灵长之文明永难磨灭。从武丁时代开始，中国拥有了迄今可以确认的最早的成熟文字系统。象形文字，是商朝的祖先留给我们最宝贵的遗产。至今，这方正汉字的一笔一画间，仍存续着商人眼中那日升月落、风雨雷电的万象，古今攸同，仍维系着这三千多年来每一颗以汉字启蒙的心灵，无问西东。

武丁时期，开始实行“周祭”制度，也就是轮流祭祀自先祖上甲以来的历代先公、先王和王后，周而复始，一年完成一个祭祀周期。而升入天宫的历代先王们，此刻正陪伴在“帝”的身边，“帝”是至高之神，可赐予丰收，护佑凯旋，唯有在世商王才具备与祖先和“帝”沟通的能力。通过祭祀、占卜，获取“帝”的旨意和世间人事的预兆，祖先的遗骨垒出了信仰的神殿，而商王，正是唯一的大祭司。丰盛的食物来自于天赐与人功，也需以青铜盛装，才显出在人间的尊崇和对神明的敬重。巨大的方鼎，可烹煮大型祭牲，造型浑厚而庄重，如天地间最均衡和稳定的存在。簋，盛满煮熟的谷物；豆，盛放肉类或调料；甗，是复合器物，下为鬲，上为甑，可实现蒸煮一体的功能。

佳肴粢盛，盛以吉金重器，神祇与祖先想必能尽情享用。商人嗜酒，酒也是祭祀时不可或缺的祭品，它有助于让司祭者进入一种仿若通灵的状态。因此在青铜器中，酒器便成为最重要的一种器类。盛酒的瓿、彝、罍、尊、盉、卣、壶、觥；温酒的爵、角、鬲；饮酒的觚、觯，玉醴香醪，配以金铜炫曜，这来自人间的美酒芬芳想必能直达云霄。这日，武丁接到西边的周族被降服的消息，尽管当时生活在关中平原的这支小族实在势微力弱、不值一提，但恰逢大祭前夕，武丁觉得是吉兆，仍十分高兴。

每逢献祭，武丁便饮尽杯中酒，他注视着眼前那尊金色的大鼎，燎柴之烟氤氲缭绕，神光离合间，鼎上的蟠龙升腾，神鸟飞起，那双目炯炯的饕餮幻化为庞然巨物，凌驾于上空，恍惚间又只余缕缕青烟，与天上云气若混作一体，不知何处飞来的一只锦雉，竟立于鼎上，引吭清啼。武丁大感惶恐，生怕祭仪中这样一个偶发事件会成为不祥之兆，身旁的臣子祖己却劝解商王，只要勤修政事，便无可畏惧。

商代晚期的青铜器，各种器形已臻于成熟之境，与稳重简洁的造型形成对比的是纹饰的繁缛。铺满整器的地纹中，又以更粗重的线条与凹凸，勾勒出各种具有丰富意象的纹样。

狞厉的饕餮、蜷曲的夔龙、昂首的凤鸟、重生的蜕蝉，这些复杂纹饰熔铸于圆弧的容器表面，却无一处废笔或错漏，非高超至极的工艺和设计能力不能达成。商人通过精湛的青铜重器，营造了一种光怪陆离、神秘莫测的氛围，这也是商人浩荡而深窈的精神世界。神灵漫天、异兽满地，人之在其中，惟有商王可以沟通天地。而这些集稀缺资源、复杂工艺和神秘知识于一体的青铜礼器，其组织生产、使用方法和器用制度，也都一概掌握在商王的手中。独占神权，只实现了精神层面的统摄，这尚不足以使商的统治者高枕无忧。唯有配合以强大的军事实力，才能构成王位之畔的坚固藩篱。

商的军队组织周密而完备，装备有最精良的杀伐利器，钺是砍杀之器，也是军权和王权的象征；戈宜钩杀，矛宜刺杀，强弓劲矢可远程伏射，甲胄能护身御敌。狩猎，是日常练兵的一种方式，远处一头野牛出现，武丁令战车驱逐，不料，身后的小臣叶所驾之车却撞将上来，两驾车都毁坏，王子子央也堕车受伤。祭祀狩猎涂朱牛骨刻辞，记录了这起追尾事故，在写到“车”字时，利用象形文字的优势，极形象地展示了车祸现场，一“车”车轴断裂，一“车”倒翻。

距今约4000年前，乌拉尔山地区出现了世界上最早的马拉车，中国最早的成熟马车则出现于晚商的武丁时期，形制与欧亚草原的马车相似，匹配刀、弓形器、马策等工具组，表明其驾驭方式也与草原马车近似，暗示这些马车是草原地区传入，而非本土起源。

但以成套铜饰件装饰马车的做法则是商人的发明。马车的出现改变了战争形态，在此后的900年间，车战成为古中国大规模战争的主要形式，千乘之国是强国的象征，万乘之尊是帝王的称号。战争于胜者而言是荣耀，却也不可避免会带来伤亡。这位年轻的商军士兵被敌人箭镞射入头骨，死在了战场之上。即便在战场上侥幸苟活，一旦被俘，却可能面临着更加悲惨的命运，作为人牲献祭，是大多数战俘最终的结局。卜辞记载，武丁时期用于祭祀的人牲数量可达9021人，最大规模的一次足足献祭了500名人牲。殷墟曾出土过两片刻有文字的人头骨，其中一片头骨上刻写的内容是“方伯用”，意思是用“方伯”来献祭；另一篇则刻了“方伯，祖乙伐”，也就是说，用“方伯”来献祭商王祖乙。

方伯，是商人对方国首领的称呼。商人将这两位怀有不臣之心的一方之主制成了刻字的祭品，以此献给那长随天地之侧的祖先。“大邑商”，是商人对自己这巍峨国都的称呼，宫殿、宗庙与王陵，在世俗政治和宗教神权体系里，都有着至高无上的地位，统摄着整个王朝。以殷墟为中心，北起沙丘、南到朝歌、西至太行山、东抵古黄河，这便是商的王畿所在。水渠和道路由这中心生长出去，如虬须盘根般蔓延至商的每个角落，由王畿向外延伸数百公里。侯、甸、男、卫等各级贵族的领地，构成商的藩屏，方国势力交错其间。黄昏时分的豫北平原，北行的牛车队伍，正将四方奉献的贡品载向王都，其中不仅有谷物、象牙、龟甲，还有海贝、铜锭，甚至是鲸鱼骨骼。

另一辆双马轻车在路上奔驰，驭手是长族首领亚长的属下，亚长奉商王之命出征，已领军与东夷人激战月余，眼见胜利在望，亚长令驭手返回王都，向商王汇报战况。此刻天色近晚，不远处便是“羁”，也就是商人的驿站，按照原定的计划，在此休息一夜，明日一早起身赶路，午间便能回到王都。驭手念及将能见到分别许久的家人，心中欢喜，亚长献给商王的战利品，由驿卒代为看管，这包裹里却是他自己攒下来留给家人的礼物，给妻子的玉饰、给儿子的青铜刀、给女儿的海螺。他检视了一番，确定这些宝物并未因路途的颠簸而毁损，才放心地珍重藏好。驿站外忽然传来一阵急促的马蹄声，一位骑士推门进屋。

驭手认得，这是亚长的侍卫，他的突然出现让驭手心中十分不安。果然，侍卫带来了噩耗，在驭手出发后不久，本以为稳操胜券的亚长竟血染疆场。侍卫转告驭手不必再赶回王都，因为护送遗体的人马，将于三日后启程返回大邑商，驭手需在驿站中等候，而侍卫则要星夜兼程赶回王都报讯。驭手取出那个珍藏已久的包裹交给侍卫，嘱托他一定要带给自己的家人。驭手心中悲痛，他知道这样好的夜色，自己不能再看几回了。按照商人的习俗和礼仪，亚长一死，意味着作为亚长家臣的他，势必要殉葬于地下。迎回亚长遗体的家人，正沉浸于一片哀痛中。

亚长的遗体已经清洗干净，敌人集中袭击了他的身体中部，左大腿根部正面被锐斧劈砍，他同时受到来自左后方的利戈的袭击，动脉被割破，盆骨右侧被长矛刺穿，左臂留下至少3处锐器砍伤痕迹，左肋被锐器击伤。亚长伤痕累累的身体，将被花椒完全覆盖，这是亚长家乡特有的防腐方法，外面再裹以麻布和锦缎，面朝下放入棺中，这样的俯身葬也是家乡的习俗。深夜，乡村的茅草屋中，驭手的妻子仍等待着离人归来，这战火再起的时代，或许没有消息就是最好的消息。马车，本是身份地位的象征，乘坐马车者，通常为王公贵族，他们死后，马车将作为陪葬品埋于墓内或其附近，马匹和驭手常常被一并埋入。冬日阴霾的天空下，西风萧瑟，犬鸣马啸。

驭手跪坐在商王为亚长择定的墓旁，和他一起的，还有另外14位将一同殉葬的家臣，其中也有那夜报讯的侍卫，那个面容稚嫩的年轻人。有人牵来了要殉葬的狗；临行前，驭手饮了些酒，这是他人生的最后一杯酒。

殷墟王宫区的东南侧，花园庄东地的54号大墓中，墓主是35岁左右的长族首领亚长，与他一同下葬的还有15位殉人和15条殉狗，各类随葬品约1800余件，其中方鬲、方鼎、牛尊等，是高等级人群方能使用的器物。青铜礼器，尤其是觚和爵的数量，是社会等级的标志。妇好墓出土53觚40爵，亚长墓随葬9觚9爵，数量之别，正是等级之分。商人构建的这套以青铜器为载体的礼制系统，将祭祀礼器制度与社会等级制度相结合，并以特定器类的数量展现出来。

这样的礼制设计，直观而明确，覆盖面和接受度都很广，从而成为商代构建强大国家体制的重要内容。

对亚长牙齿的锶同位素检测结果显示，他的幼年并非在殷墟度过。他所属的长族，封地在今河南东南部的鹿邑，毗邻东夷和淮夷，是商的东南门户。为了争夺土地，控制海盐，商与东夷交战不休。济南的大辛庄遗址，是一处商王朝深楔入东夷腹地的军事据点，出土了与商王朝同样的青铜礼器组合，还发现了刻辞卜甲，这是晚商时期，殷墟以外罕有的发现了甲骨文的地方，无论是字形、文法，还是甲骨修正、钻凿形态，都与安阳殷墟卜辞属于同一系统。山东青州，则是另一处两军交锋的前沿。苏埠屯1号墓的墓主，是“亚醜”族的首领，采用了殉人、腰坑殉狗等典型的商人葬俗。

带“亚醜”徽记的铜钺，表明他拥有商王赋予的军权，为商镇守东土。

那么，遥远的南土，商人又该如何控制？青年是商的王族，受商王之命，巡视长江沿岸形势，筹划治理南土的良策。早先镇守南土的盘龙城早已凋敝，王朝疲于应付与东夷、西北的舌方（吾方是殷墟出土卜辞中见到的一个商代的方国部落）、鬼方、羌人之间的冲突，久已疏于南土的治理。青年望向江面，回想起在赣江中游所见的奇特祭器，那狰狞的猛虎躯体硕大，呈伏蹲欲纵之态，虎周身遍饰花纹，背上伏卧一只小鸟，宛若猛虎的驾驭者。头生双角的人面，神与人似融合一体，面容神秘诡异，充满威严。青年打开一幅帛画，所绘青铜礼器，正属于湘江中游的方国。四羊方尊、人面方鼎，虽与商器相类，但形象生动，别具特色。

商族青年逆流而上，进入三峡，两岸群山掩映，江水变得湍急起来。此刻，古蜀国神殿的密室中烛火通明，工匠正在打着一副黄金面具。方面大耳，直鼻阔口，三角形眼孔，幽暗深邃。经由宽阔的江河，抵达终古的密林，商族青年似被感召，径直踏入这古蜀国的神秘地界。硕大的太阳转动不休，威严的瞋目俯瞰尘世，商族青年恍惚之中，忽然心惊，商王独占了与帝沟通、祭祀先祖的神权。可是若这天下，各地有各地信仰之神呢？王又如何能控制四域？古蜀神殿的大门，正在徐徐敞开，硕大的青铜面具凸目阔口，造型意象上体现出人神合一的特点，或许是古蜀人的祖先神造像，双手举持法器的青铜立人，可能是其神、巫、王三者身份于一体的权威人物。

青铜神树，似从象征“神山”的底座上生出，枝叶蜷曲，奇花绽放，枝桠上立着昂首的神鸟，树侧有逶迤而下的铜龙。这是青铜鸟足神像，最下部为青铜方座罍，罍表可见少量红彩残留，器物中部为青铜神像，纵目、獠牙、额生牛角，以双手为支撑，倒立在青铜罍之上，他头顶一件青铜觚形尊，尊盖上还站着手持龙形器的立人，整器，似在表现鸟族人面的神祇从天而降。这件青铜神坛的镂空台基之上，共有13个小型青铜人像，其中有一人，以跪姿背罍，似为女性；抬架之上，立有一只如犬似马的青铜神兽，在神兽的头与尾之间，跪坐一人，手举山形座，神坛，由四根细体觚支撑，侧立青铜持鸟立人像，圆盘上方，有四个顶尊跪坐人像。

这件器物，构成了一个祭祀活动场景，可能反映了重大仪式活动中的人员分工以及宗教观念。牙璋，是自二里头继承来的玉礼器，传至古蜀国后，往往在牙璋顶端劈出尖嘴，使其酷似鱼形，有时又在鱼嘴中雕出一只振翅之鸟，仿佛鱼鸟化生。这件跪坐小人像，总高不足5厘米，下身着裙，上身赤裸，双手持握一璋，是古蜀人以璋行祭的实物例证。三星堆人，对商礼制系统中，尊、罍等器物的推崇，以及对中原独有的青铜范铸技术的采用，表现出自身与中原王朝的密切联系。四川广汉三星堆古城，可调控环境参数的考古舱内，精细的发掘正在进行。三千多年前，古蜀人可能在仪式中，亲手砸碎和埋葬了这些青铜器。

三千多年后，经考古工作者的双手，它们重见天日，并将在此后被一一修复、渐渐弥合，从而再现出神秘古蜀的宏大时空。三星堆的历史，是见证中国多元一体交融汇聚的缩影，距今4500年开始，西北地区的先民带着彩陶文化，经过川西山区进入成都平原的西北地区，同时抵达的还有经三峡地区而来的长江中游先民。此后漫长的岁月里，在中原和长江流域的青铜文明的共同作用下，才诞生了古蜀国这样一个区域文明。步入王国社会鼎盛时期的中原王朝，此刻已感知到潜在的危机，由他们树立的青铜礼制、文字体系、神权信仰和军事规范，有力维系和影响着辽阔疆域中的各种力量。但与此同时，也培养了自己的挑战者。

长江沿线，东有江西新干大洋洲文化，中有湖南宁乡青铜文化，西有四川广汉三星堆文化，商的南土，已经被这些新兴方国所控制。在西北地区更是强敌四起，陕北一带，卜辞中的“鬼方”，晋西山地，卜辞中的“舌方”，双双屹立在沟壑纵横的黄土高原。被后世称为纣王的帝辛，征服了强大的东夷后，继而又征服了蓝夷，将淮河和长江下游纳入实际控制区内，《史记·殷本纪》中的帝辛，见闻广博，勇力过人，但穷奢极欲，残暴不仁，其中有多少是史实，多少是周人的贬低，不得而知。殷墟考古中，各个时期的墓葬数量以帝辛时最多，各族邑遗址中发现的大型四合院建筑群、带墓道的大型墓葬及车马坑，也多建于帝辛时期，这表明王朝人口的蓬勃增长和国度的空前繁荣。

但帝辛终未能改变商的命运，靡不有初，鲜克有终，成汤时对子民的顾念，并未被后来的商王们继承。商人好问，他们问天、问神、问祖先，凡事他们都要问卜，直到有一天，神不再应答。不过，来自商的祖先，无意中还留下了他们的样子，这尊陶土塑成的、比心还小的脸庞。他是谁？他是为主人殉葬的驭手，他是向命运索取答案的贞人，他是顺着江河漂泊的青铜匠人……他望着我们，仿佛有很多话要说，那是甲骨文也无法言传的、一个人平凡的悲欢离合……而在西北的黄土高原上，出自戎狄之间的周人已经蓄势待发，他们，将带着我们的中国，走向闪烁人性光芒的家国天下。

\chapter{第七集 \quad 家国}

\section{正文}

陕西长武，泾水支流黑河的北岸上，碾子坡遗址笼罩在一片苍茫暮色中。这里，或许就是3000年前的豳，这里是黄土高原的南缘，千沟万壑分割出支离破碎的梁峁塬地，这里也是农耕区与半农半牧区的分界带。商以农为本，夷狄以牧为生，而周人的祖先曾摇摆于这两种生活方式之间，他们曾被戎狄侵扰，却又被商人当作戎狄一样征伐。商王武丁时期的甲骨上，历历可见与“伐周”相关的卜辞。古公亶父成为周族首领时，带着族人们从事农作，过上了家有余粮的安定生活。他常常站在豳原的高阜上南望，山势连绵，直接天际。他知道，翻过那一道道山梁，就是广袤的关中平原，那里，才是更适宜耕种的沃土，秋风萧瑟，熏育戎狄再次来犯，先掠才无，又意图占有土地和人口。

正当族人准备与戎狄决战时，古公亶父却制止了众人。他的理由是，对于普通人来说，在谁的治下生活并无不同，但战争却随时会令人丧生，因此他愿离开豳，将土地和人口都留给戎狄。

“古公亶父，来朝走马”，“率西水浒，至于岐下”。这是《诗经·大雅·緜》中周人对先祖自豳地迁往岐邑的记忆。古公亶父的退让，被豳人视为爱民的表现，纷纷离开家乡，扶老携幼，跟随他来到岐山之下。周边部族听闻后，也陆续前来归附，他们自此不再游牧，而在周原起城郭、筑屋舍、务农事，但这也意味着他们进入了“大邦商”的势力范围。商，实力强大，四处征伐，殷墟西北岗王陵中，小屯乙七建筑基址的祭祀坑，非王族的墓葬中，都发现了大量从葬的人骨，此时的周，不过是商邑西陲的蕞尔小邦，只能暂时的臣服于这个强大的王国。

与那个时代漠视生命的风气不同，姬姓周人的墓葬中，自始至终都极少有殉人，也几乎不见殉牲。多年来辗转求生的周族，深刻的认识到只有依靠群体的力量，才能存活于暑雨祁寒间，而仁德，就是凝聚起群体力量的黏合剂。古公亶父牵着小小的姬昌，在田间缓缓行走，耐心的教他辨识每一种农作物，娓娓讲述他们的习性，稚童眼神明亮，听得认真，努力记住祖父所说的一切，祖父常常讲起先祖后稷作农，并惠泽天下的故事。祖父还告诉他，人心，才是弱小的周族拥有的最大利器。夕阳一片，粟黍青青，祖孙俩脚下踏过的这片沃土，就是周原。周原遗址在今天的陕西扶风、岐山两县交界处，总面积约30多平方公里。诗经中说“周原膴膴，堇荼如饴”，意思是在周原这片丰饶的田野上，连生长出来的野菜也甘美如饴。

土地的膏腴不纯然来自天赐，也来自人力的开垦。从空中俯瞰，周原之上，周人曾以人工开挖的河道沟通了自然水系，以水库为调节枢纽，构成庞大的水利系统，灌溉着千里沃野。在周原的王家嘴，考古工作者们找到了灭商以前周人的世居之地。一号建筑基址正是那时的官室，它的总面积超过2200平方米，由前堂、后室、东西厢房、前后庭院和附属建筑组成。这种前堂后室的两进四合院式布局，明显以商的宫式建筑为模板。在王家嘴的东北方向，一水之隔处，发现了大小两重城垣，西北部的规整长方形小城，始建于商周之际至西周早期，大城城墙则是到西周晚期才扩建的，是万民所居的郭城。小城北部正中的凤雏建筑群可能是宫殿群，3号建筑基址是目前所知西周规模最大的单体建筑，在院落内发现了与土地祭祀有关的社石。

在这处建筑的南侧，还发现了一辆驷马驱驾的青铜轮牙马车，马匹及车辆装饰有大量镶嵌绿松石的青铜饰件，车厢外围装饰有薄壁青铜兽面、玉器、绿松石饰件。这辆华丽的西周马车，应该是高等级贵族的礼宾用车，以仁德和辑人心的周渐渐强大，在商王的授意下，周逐渐承担起为商治理西土的角色。姬昌的父亲季历，被商王帝乙封为西伯，率领周师四处征伐戎狄，并定期朝拜商王。仲春之月，渭水两岸桃花绽放，如云蒸霞蔚，河洲之上，雎鸠关关，众人簇拥着年轻的姬昌从奇异的宫殿中迤逦而出，来到渭水边，河对岸走来了他的新婚妻子太姒。周族传统的联姻对象是同样生活在商人西土的姜族，传说中正是姜族的女子姜嫄，诞育了周族的祖先后稷。姬昌的祖母也是姜族的太姜。

但自姬昌的父亲季历起，周族首领也开始与东方诸族联姻，试图融入中原社会。姬昌的母亲太任，是殷商的贵族挚任氏，姬昌的妻子太姒也从东来。凤雏建筑遗址出土了17000多片卜甲卜骨，其中有字卜甲有190多片，字体纤细如微雕，所书文字与商人文字无异，所记载的占卜程序也与商人相同。卜辞多属姬昌时期，内容包括向殷商王室的祖先成汤、太甲、文丁、帝乙等祭祀，其中还出现了周方伯的称呼，这些都足见周人对商的顺服。另一方面，姬昌继承了古公亶父以来的理念，笃仁、敬老、慈少、礼贤，前来归顺于周的人越来越多。岐山苍翠的林中，几声清啼打破清晨的寂静，一只锦雉掠过晨雾落在树梢上，而洹水之畔的大邑商，商王帝辛彻夜未眠，他摒弃了职业的贞人，亲自手持甲骨，细观纹理，推究上帝和祖先的启示。

突然，一阵鸟叫声打破了他的迷思，他惶惑地望向窗外，周的强大终于引起了商人的警惕，姬昌被帝辛囚禁于羑里七年，经周臣的多番奔走才终能得释。姬昌归去后，一面继续修德政，一面为商王征伐叛邦，赢得越来越多诸侯的归附，但很快姬昌在迁至丰京的次年死去。姬发即位后，修建了镐京。丰镐遗址在今西安市长安区，沣河纵贯两京之间，丰京在河西，镐京在河东，隔河连片，蔚然可观。镐京遗址内，已发现了十余座大型夯土建筑，还有道路和排水设施。郿坞岭高地14号建筑的规模宏伟，附近还有瘗埋动物骨骼的祭祀坑，或许就是周人的宗庙。周人以血缘确定贵族的亲疏、等级、分封和世袭关系。姬发居于宗周，也就是丰镐，他是姬姓宗族的族长，也是国君，掌握着族权和国家最高权力。

都邑之外，环绕着王族子弟的采邑。这里是国也是家，有国之架构，也有家之血缘。周都的东迁和规模扩大，无不宣示着这个曾经的西陲小邦如今的壮大。他们披仁德之坚，执宗法之锐，酝酿着对东边那个大邦的奋力一击。暮色渐深，马车仍在山路间向残阳落下的方向奔行，艰难的路途如同鸿沟，隔离了远方的山河。太师疵和少师彊两位商的重臣牢牢护着身旁的编铙，这是商王朝最神圣的祭乐之器，也是他们家族世代职责所司。他们生怕一路的颠簸令其稍有毁损。由于帝辛的残暴，他们不得不和许多殷商贵族一样，选择离开故土与故国，逃往以仁德闻名的周。在伐商的檄文《牧誓》中，武王痛责纣王的失德，他将克商之战称为“行天之罚”。

在那个极度崇拜神巫的时代，周人却从世俗层面，提出统治的合法性应该来自天命，而天命不在于祖先与天的特殊关系，不在于独占通天之技术和资源，主要在于统治者的德。有德之君，就会被天授予神圣的大命，得此大命者，可以革除失德王朝的旧命，建立新的王朝，抚育万民。陕西临潼零口镇出土的利簋，其底部铭文记录了武王伐商的事件，“珷征商，佳甲子朝，岁鼎，克昏夙有商”，意思就是甲子日的清晨，时逢岁星当空。武王的军队仅用一日便结束了战争。身为右使的利，由于跟随武王作战得胜后受到奖赏，因此铸造了这件青铜簋用以记功，并作为祭奠祖先的礼器。考古学家根据铭文中提及的特殊天象，再加上对干支纪年的推算，认为武王伐商的具体年代可能是在公元前1046年。

商，以灿烂的青铜和完善的文字体系，将我们的中国文明推送到新的发展高度。但是商的覆亡也说明，对神权的独占以及大规模的征战和杀殉，已经不足以威服四方。面对辽阔的国土和多元的族群，如何将之真正纳入到统一中央权力的管辖之下，周人需要新的政治方略。周人代商而立后，妥善安置了殷商遗民，帝辛之子武庚被封于殷商故都，仍可延续殷祀。但同时，武王又将自己的三位胞弟管叔、蔡叔、霍叔分封在殷都附近，起着监控的作用。他最信任的4位重臣则被分封在原商朝疆域的东、北、南境。姜尚在齐国，周公在鲁国，召公在燕国，南公在曾国，对商构成巨大威胁的东夷、淮夷曾活跃在这些地方，武王不敢掉以轻心。这一年，箕子从东边浮海而归，前往周都朝见周王。

归途中，他特地折往商人旧都，只见宫倾城圮的废墟上，已是禾黍青青，箕子是殷商王族，也是与微子比干齐名的重臣。曾经的朝堂之上，面对一意孤行的帝辛，他没有像比干一样死谏，也没有像微子一样离去，而选择了装疯，终致被囚为奴。武王之前，他也曾面陈《洪范》献上治国之法，但终不愿仕周为官而远走。眼前，太行西峙，洹水东流，故景依旧。那个曾令天下“莫敢不来享，莫敢不来王”的天邑商，却已无迹可寻。箕子忍不住悲从中来，在长天后土之间，怆然的放声唱道，“麦秀渐渐兮，禾黍油油，彼狡童兮不与我好兮”。诗中的狡童，正是商朝的末代之王，他的侄子帝辛，早已在武王兵临城下那日自焚于鹿台之上。殷商遗民听到这支歌，莫不垂泪。他们的王朝啊，终是沉入历史的长夜了。

月华如洗，周王的寝殿中依旧灯火通明，武王克商返回镐京后，却陷入比出征前更深的忧虑中，这天，他又再一次的熬到长夜欲曙，忍不住派人换来了周公。武王和周公的彻夜长谈，牵涉到周人疆域和身份观念的转型。周人本来自认为是西土之人，而商人所在则为东国，全部占有商王朝的原有领土后，周人的志向并不仅仅是取商人而代之，还希望能建立一个规模空前的王国。一番长谈后，二人商定，要在伊洛二水之北营建洛邑，以“定天保，依天室”，天室，就是太室山。太室山为天下之中的观念由来已久，夏商两代的都城都在这附近。武王返回镐京途中，曾在姜尚的陪同下，登上太室山山顶，俯瞰四方山川形势，那时他大概就已经存在营建洛邑的念头了。

这是何尊，出土于距周人旧都岐邑不远的陕西宝鸡贾村镇。尊，是酒器，同样也是礼器，这件青铜器是武王之子周成王在位第5年时，周的宗室贵族何，为纪念成王的训诫和封赏而作。何尊高38.8厘米，重达14.6公斤，内底铸有12行122字铭文，铭文中提及，成王将政治中心迁往成周时，追述了武王克商后的敬告上天之辞。“余其宅兹中国，自之乂民”，这是我们目前看到的最早的以文字书写的“中国”之名。成王即位后，在周公、召公的辅政下，完成了武王营建洛邑的遗愿，自此后，丰镐为宗周，洛邑为成周，也即是中国。在中国立都，周人便不再是西土之人，周也不再是西土方国，而是位居天下之中的宗主国。普天之下，莫非王土？这里是周人以天下视角确立的自然地理和政治地理核心。

周公旦在余生中仍常常想起他与哥哥姬发的那次彻夜长谈，除了营建洛邑之外，那时已知道自己病体难支的武王，还提出了另一个请求，若自己身死，便请周公继位。震惊之余，周公答道，周人的立族之本乃是宗法，族权自来是父死子继，国，亦当如此。王权的子继之法一旦实行，嫡庶之制、宗法之制、分封之制，便随即顺理成章，将密织成一整套坚固的上层建筑。周公所思，武王岂能不懂？身为国君，他惟愿自己死后仁德如周公者，能接过这未竟的事业，带领根基未稳的周，真正成为长久存在的大邦。周公的心术与规摹，并不在于眼前一时的权利，而在于周的万世治安大计。因此他选择了不登王位，不就封国，只留在王畿辅佐年少的成王，平武庚之乱东征淮夷，劳心竭力地辅政，敦促成王立德修身，既是为宗族，亦是为天下。

在周公的辅佐下，周王朝一面牢牢掌控着王畿之地，一面通过分封的诸侯国于四方之土，探索着不同的治理模式，使各地的不同部族，得以融入周人的政治秩序。鲁国的伯禽，面临着相当复杂的局面，当地遍布东方诸族，文化迥异。当看到熟悉的家门时，伯禽忍不住加快了步伐，走入堂中。他离开父母前往鲁地就封已经整整3年。因为父亲周公要留在成王身旁辅政，便只能由他代为就封。此番借着回来报政，他终于可以与父母相聚几日，堂上的周公白发暗生，面容因为常年的殚精竭力而略显疲惫，但双目依旧精光四射，仿若能洞悉世间的一切事理。他和蔼而专注的听伯禽汇报着三年来在鲁国的治理成果，频频颔首。末了，周公似乎想起了什么，问伯禽为何3年才回来报政？伯禽早有准备，从容答道，因为需要变其俗，革其礼。

周公告诉伯禽，太公姜尚就封鲁国之邻的齐国后，精简君臣之礼，依从当地习俗，仅过5月便回来报政。伯禽愣住了，周公告诫道，只有所施政策平易近民，才能得民心归附。周公所称许的齐国，是周初三公之一太公姜尚的封地。齐地虽有工、商、渔、盐之饶，却也是直接面对东夷诸部族的最前线。近年在山东临淄高青发现一处西周墓地，墓葬规格很高，也有祭祀遗迹和车马坑。大墓内出土的铜器，出现了齐公字样的铭文。若这里是齐公墓地，那么或许最早的齐国都城便在这一带。莒县西大庄春秋墓中出土的一件铜甗，可能是齐侯嫁女的陪嫁。莒是东夷势力之一，藉由此器，可以想见当时齐国通过各种手段融合东夷的努力。燕国所在之地，是周人克商后新辟的封地，为周王朝守卫着北境，当地同时还生活着迁徙而来的殷移民以及本地的土著族群。

这是北京房山琉璃河遗址出土的克盉，铭文“命克侯于燕”一句，意指周王将克封于燕，治理这里的疆土和六族人民。克，是太保召公奭的长子，因为召公留在王畿辅佐周王，故由长子代为就封。然而，1902号墓中出土的作册奂卣中太保墉燕的铭文，说明召公本人不仅亲自来过这里，还主持了燕都的营建。1902号墓的墓主人是一位殷遗民，他与周人同葬于一片墓地内。这位名叫奂的史官，生前颇受重用，随葬的青铜器上，铭刻着曾受周人封赏的荣光。他规矩的沿袭着商人殉狗的旧习俗，却也认同于周人建构的新秩序，由此可见。周人在族群面貌复杂之地采取的并存策略，收效应该比较显著。与燕地一样，汉水流域也是周王朝灭商后新扩的领土。

考古工作者在湖北随州附近，发现了一个失载于文献的曾国，曾之始祖南宫适也属姬姓，是辅佐文王和武王完成灭商大业的重臣，因为本人要留在王畿领职，只能由儿子代为就封。叶家山墓地中发现了曾侯犺、曾侯谏等数代西周曾侯的墓葬。这个从未出现在任何传世文献中的国家，现在人们普遍相信，他就是左传中汉阳诸姬之首的随国，一直存在了700余年，留下众多两周时期的珍贵文物和历史信息。这件方鼎是叔虞自作的铜器，叔虞是武王之子，成王之弟，亦在成王时期获封唐地，留下了桐叶封弟的传说。叔虞之子燮父继位后，迁都晋水之傍，将国号自唐改为晋。晋南是夏人故地，自然启以夏政，那里又遍布戎狄部族，亦须疆以戎索，融合各方，正是晋国肩负的重要责任。

燮父墓中出土了一件精美的鸟尊，器身造型生动无比，通体纹饰华美，堪称西周青铜艺术的杰作。鸟尊盖内铭文中的晋侯，就是燮父的自称，这个晋字见证着燮父迁唐于晋的历史。在山西曲沃曲村天马遗址的北赵晋侯墓地中，发现了9组19座排列规整的晋侯夫妇墓及附属的车马坑，视为国君夫妇的专属公墓。每代晋侯夫妇几乎都是一夫一妻异穴合葬，唯有64号墓中的晋穆侯，身旁并穴而葬了两位夫人。其中一位夫人墓内随葬的组玉佩，是迄今所见最复杂的。周代组玉佩是王室贵族们进行祭祀、朝聘、婚配等活动时佩戴的仪礼重器，不但规范了君臣尊卑，也将周人的道德观念赋予其上，是西周社会礼制日臻完善的重要体现。成年的姬诵头戴冠冕，身着朝服，肃容立于大朝会坛场上。周公在左，太公姜尚在右，皆衣冠济楚。

堂下为首的是尧舜禹商后裔，他们的出现象征着周的合法王统。诸侯们依次献上国内珍产，东方的海蛤、玄贝，南方的翠鸟、大象，西方的孔雀、朱砂，北方的麋鹿、大鲵，琳琅满目的异产，象征着各地对周的臣服。这场隆重的大朝会礼仪，从时间和空间两个维度上，宣示着周天子的权威。周初共分封71国，其中与周王族同姓的姬姓就有53国。分封，是周人总结商人治理得失，结合自身文化传统的创造。这项制度的核心，不是强力的军事控制，而是不同族群在周人统领下的融合，外封的诸侯都承担着以周文化融合当地族群的使命。周人将根植血缘关系的宗法制与国家治理紧密联系，使得家国一体。凭借稳定的社会体系，周天子的礼仪政令得以有效传递，周王朝的国家机器得以有效运行。

陕西扶风庄白一号窖藏，是新中国成立以来，出土青铜器数量最多的窖藏，在面积仅2平方米的窖穴内，就收纳了103件青铜器，其中74件都有铭文。在这件共王时铸造的史墙盘上，追述了文武之后迄至穆王的各代周王功绩。成王严明纲纪，封建邦国，康王开辟疆域，昭王广拓荆楚，穆王深谋大略。铭文在追述周人世系时，最大的变化，是直接以始祖后稷与上帝并列，称上帝后稷亢保。这意味着，周人将本族的祖先后稷，往前推至于舜禹同时，将族史与国史交织，逐渐融入正统华夏的叙事。这一古史系统的完成，应是在众多变革发生的穆王时期。夏将至，伯中受穆王之命，率兵在棫林追击淮夷，最终克敌制胜，穆王后王俎姜派人来到前线对他进行赏赐。凯旋回朝后，他便马上赶往故乡告祭父母。父亲在他幼年时便已辞世，是母亲将他抚育成人。

以往每次出征前，母亲总会端正他的衣冠，细细检查他的甲胄是否有破损，嘱咐他既要勇敢战斗，也要保全自己。淮夷之战的不久前，母亲刚刚去世，再无人那样絮絮地叮咛。此刻，他在心里默默告诉母亲，他平安归来了。伯中为母亲所做的簋上，铭刻着一个儿子对亡母的殷殷思念之情。“母亲的美好品行，在我心中涤荡，永远护佑我身，让我战无不胜，战毕，我身无伤”。同墓所出还有伯中为父母所做的两件方鼎，一件记述了穆王之命，一件记述了穆王后之赐，从周穆王、共王开始，伯中墓中这样鼎、簋配合的食器，成为了新的礼器体系主流商人创立的以觚、爵等酒器为主的青铜礼器系统被逐渐摒弃。周人认为，商人纵酒是致使其亡国的因由之一。

从武王的《牧誓》到周公的《酒诰》，从西周早期康王的大盂鼎，到西周末年宣王的毛公鼎的铭文中，都一再警示周人上下勿要酗酒，或许因此，周人才逐渐创立了自己的礼器体系。此时，青铜礼器的风格也发生了改变。商代晚期以来，华丽繁缛的青铜器装饰风格逐渐被简化，质朴简洁的鸟纹、几何纹逐渐成为主流。分封制构建的家国体系、以德配天的政治理念被寄寓在这些世俗色彩浓厚的青铜礼器中，物化成为礼制。周人之礼，如春雨润物无声，如阳光普照四方，衣食住行、婚丧嫁娶，家国朝野，无处不在。墓葬中的青铜礼器，便是这套礼制最直观的表现。根据《春秋公羊传》的记述，在礼仪活动中，天子使用九鼎。在他之下的诸侯、大夫、元士，则依次使用七鼎、五鼎、三鼎，这套礼器制度，在三门峡上村岭虢国墓地中展露无遗。

这些自西周末年延续至虢国灭亡的墓葬中，以食器、车马器、乐器的数量和有无分为五等。上世纪90年代，虢国墓地进行了大规模考古工作，发现了虢季墓和虢仲墓两座国君级大墓。虢仲墓随葬各类鼎20余件，其中包括一套大小成列的七鼎，与其王室卿士的身份相符。第二等级之下的墓葬皆无乐器，其余器类的数量逐级递减，五个等级之间，不止青铜礼器的种类和数量有别，食器所盛菜肴，车马所载出行，黄钟所奏雅乐更是不同。这百余年间，数百座墓葬都遵守如一的规则，正是雍雍穆穆的周礼。清道光末年出土于岐山的毛公鼎，是毛公为纪念周宣王的册封和赏赐而作。铭文中宣王先是追怀文王、武王时以仁德肇国，君臣一体治理天下的盛世，继而感怆时艰，为局势不宁、国家濒危而忧心忡忡。

由是恳切请求身兼重臣和宗亲的毛公，要为我邦我家而操心劳力，需勤理政事，勿要搪塞庶民，勿让官吏徇私，勿令鳏寡失养，望能约束僚属，勿使酗酒，你切莫失职，要夙夕谨记守业不易。从宣王的话语中，那个时代的风雨如晦扑面而来。西周王朝自始建后，在相当长的一段时间里，都将疆土经营的重心放在东土、北土和南土，而对宗周所在的西土，或许正因为是本族发祥之地，反而自信的认为能将其置于掌控下，这样的疏于防备留下的隐患。宣王时期的短暂中兴，终究无力回天。至幽王时，这种东西失衡的战略布局终于彻底暴露其弊端。申侯以女儿申后和外孙宜臼被废一事为由，联合缯国与犬戎，杀幽王于骊山下。此后，平王迁往成周，数百年来的祖居故地不得不被放弃。

在周原一带，发现了西周晚期的百余处青铜器窖藏，其中出土铜器数量较多，且有长篇铭文的重要发现就有十几处。当时这些铜器群应该分别被不同的贵族家族所拥有。然而他们为何几乎在同一时间段将这些重器藏于地下，最终又未再取出？有人联系西周末年的局势，推测可能是与犬戎入侵有关，贵族们仓促掩埋了难以带走的宝器，匆匆跟随周王东行。此后，由于各种原因，致使他们就此睽违天日，长埋地下数千年，成为后人了解一部煌煌西周史的锁钥。在周由盛转衰的同时，一支此前默默无闻的西陲小族，进入了周王室和东方诸侯的视线，那便是秦。秦人早期的历史，几乎重复着周人先祖的道路，从艰难求存于西戎之侧，到为周王朝守卫西土。幽王被杀时，秦襄公曾带兵相救，虽未能力挽狂澜，却得以崭露头角，后又以兵护送平王东徙成周，平王也因此正式锡封襄公为诸侯，并赐以岐西之地。

后来，春秋时秦武公所铸的秦公镈上便有追忆此事的铭文，“我先祖受天命，赏宅受国”。自此之后，西边的秦国开始崛起于汧渭之间，那时候，谁也不知道他们即将成为下一个时代的主角，并开启中国历史上一段浓墨重彩的篇章。孔子曰，郁郁乎文哉，吾从周。周人通过新的政治制度和政治理念的构建，形成了中华文明的人文基调，在随后的春秋战国时代焕发出更璀璨的光辉。华夏这一文明体的概念自此确立，并开启了一种内含天下结构的政治共同体形成模式。从此，不论是何种语言、文化或生活方式的族群，不论从何处发源，只要在孕育华夏文明的土地上生活，通过不断的调适认同基础，终将能融入这套无外的天下秩序里。那最初在伊洛之间的中国，势必如滚雪球般，生长成一个灿烂而盛大的政治共同体！

\chapter{第八集 \quad 天下}

\section{正文}

2400多年前的暮春，天下仍未太平，孔子却在梦中看见了久违的安宁。梦里，他和曾皙、颜回、子路这些弟子们，在春日的沂水边，身着新衣，沐浴春风，歌咏舞蹈。恍然间，对岸桃花林中有人肃立，隐约如周公。周公是孔子倾慕已久的圣人，以礼制终结了商代神权至上的治世理念，引领天下走向以人世为本的发展道路。“至哉周德”，这是目前所见，最早版本的“论语”类文献中，孔子对西周治世观的由衷推崇之语。但此刻，孔子梦中那位西周礼制开创者的身影，却已渐渐看不真切。杏林外，车马喧嚣，惊醒了孔子的梦，那是鲁哀公归城的车骑，人们纷纷传言他今日猎得麒麟。麒麟是传说中预兆圣王现世的瑞兽。

东周王室衰微，诸侯纷争不止。鲁是周公之后，理应是礼制的卫护者，如今却连鲁公也利用瑞兽来造僭越礼制的舆论之势。孔子惊觉，在有生之年里，他所怀抱的治世理想再不能如愿了。孔子理想中的天下，礼乐严明，雍容揖让。礼，是进退有止、井然有序的行为规范；乐，如舒缓优容、清和庄重的内心修养。这是曾侯乙墓编钟，由六位佩剑武士托起的巨大木架上，分三层依序悬挂65件铜钟，钟体上的铭文——标注了各钟的发音、律调、阶名，并精心标明，这些阶名与楚、周、齐、申等列国律调的对应关系。这样一套由钮钟、甬钟、镈钟组合而成的大型乐器，能演奏出由不同音色混响合成的典雅乐章，错落有致的金石之声，正暗喻着秩序整齐的礼制架构。

孔子一生起起落落，游历列国十数载，教导弟子三千，毕生正是致力于让这样雍雍穆穆的礼制思想流布天下。然而，七十高龄的他目见耳闻处，皆是礼崩乐坏。公元前548年春天的齐国，发生了一场激烈的践踏与卫护礼制之争。实权在握的齐国大夫崔杼，不顾君臣纲纪，公然杀死国君齐庄公。太史不为当权者讳，“崔杼弑其君”，短短五字，如实记录。崔杼自然不会容忍，下令杀太史、毁史册。然而太史还有两个弟弟，在那个以职为氏的年代，他们继承了史官的职责和风骨，执着地在简册上留下直笔。崔杼连杀太史一家两兄弟后，仍未能迫使史官们为其曲隐，只能作罢。不惜以性命为代价的秉笔直书，是史官们对日渐崩溃的礼制秩序的维护，也是他们在乱世中试图弥合离散的方法。

然而，在春秋三百年的动荡岁月中，旧秩序的约束力终究日渐微弱，政治格局经由血与火的洗礼，不断更迭，曾经由周天子号令的天下，转而被各怀野心的强国诸侯，搅动争霸风云。在这种情势下，宋、郑等中原小国遭兵燹之祸尤深，唯有殚精竭虑，以不同的方式辗转求生。宋国期冀通过外交斡旋来促成“弭兵”，以换取暂时的和平。宋国都城在河南商丘，今已深埋于黄河泛滥的洪积层下，经过考古勘探和发掘，我们基本了解了这段沿用千余年之久的南城墙，在东周时的模样，大部分城墙仍保留数米、甚至10米高，宽度约15米，墙体有多排纴木洞，这是修补城墙时留下的痕迹。公元前546年，这座都城曾见证了东周数百年风云中，规模最大的一场弭兵和谈。

宋都西门外，以晋、楚为首的14个诸侯国，都派出了执政的大夫参与弭战的会盟。清华简《系年》在记述此事时，虽仅有“弭天下甲兵”寥寥数字，却足以写尽沧海横流的时世里，那无数得以苟全的生灵。郑国上卿子产执政时，恰逢弭兵后的短暂和平，他从庙堂之高，俯身看到民间的芸芸众生。他深知，郑国地处中原腹心，强国环伺，朝夕有不虞之患。而子产希望邦国安宁，成为民众的依傍。因此，安内政、强军力乃当务之急。他下令“为田洫”，划田定界，防治田地不均；“作丘赋”，打破国野界限，给予乡人从军的义务和政治权利。他亦拒绝了僚属毁去乡校的建议，部分郑人对子产的咒骂传遍列国，子产自然也听说。

“苟利社稷，死生以之”。他对大夫子宽这样说道，心中的一掬赤血，终究无法剖之以示众。子产最石破天惊的一项措施是“铸刑书”，也就是把法律条文铸于铜鼎之上，公示于世人面前。列国普遍推崇礼治，仅将法制视为补充，且只有贵族才掌握律令的内容。子产却希望能让刑罚原则公开化，用成文法来规束各阶层的权利和义务，以此稳固政局。子产这一系列措施引发了部分郑国贵族的不满，也受到了列国贵族的诟责。车轮辘辘，行走在乡间的道路上，归程的子产目光悠远。夏日的木槿花娇艳而炽烈，女孩们欢声笑语，这烽火暂时消散的山水间，有人哼唱着动听的歌调，“青青子衿，悠悠我心”，“纵我不往，子宁不嗣音”。

郑地的音乐宛转悠扬，往往直抒胸臆，明快且富有生命力。虽然被正统的礼乐拥护者们，视为不入流的“淫声”，却迅速流行于各地，这不正是民间审美的抬头吗？返至家中，子产在回复叔向的责难之书中，写下了他一切作为的初心，“吾以救世也”，子产知道，战火还将再起。他肝胆俱焚，只求所有的政策能及时奏效，为郑国多争取一些时间，让那山水间平和的景象，存在的再久一些。子产被孔子称为“古之遗爱”，从子产到孔子，无不试图在这动荡的时世里，寻求重建社会秩序之法。终其一生，孔子未能见证治世理想的实践，子产也未能挽救颠坠的时局。世道艰难，但那样多的人仍在各自的位置上尽职尽责，为社稷家国死生以之，为天下大宁，匹夫怀责，这是他们的理想与信仰。

郑国国君大墓所出的一对莲鹤方壶，壶身以双龙为耳，双兽为足，腹部四角各铸一飞龙攀援欲上，通体满饰蟠螭纹，壶盖宛如一朵绽开的重瓣莲花，朵朵花瓣皆为镂空，花间站立着一只引吭振翅的鹤，昂扬向上的姿态，正如高悬于时世之上那蓬勃的朝气和理想。“鹤鸣于九皋，声闻于天”，这声声穿透混沌的清唳，在未来的历史中将会得到回应。春秋时代，小国设法求存，大国则竞逐争霸，在一众强国间，依靠实力成为天下认同的霸主。这是大国不同于小国的生存策略，黄河流域的齐国、晋国先后称霸。数十年后，齐国势头稍弱，长江流域的楚国则后来居上，与晋国形成均势。楚的崛起，与长江流域的矿产资源有关。

青铜，仍是春秋时期重要的战略资源，长江以南富集的铜、锡矿藏，原本由曾国，这一与周王室同宗的姬姓封国掌控，曾国从分封之初，除了抵御南方的淮夷之外，还承担着为周王室控制与管理铜矿资源及其运输通道的职责。曾伯漆簠的铭文中提到的“金道锡行”，指的就是连接黄河流域和长江流域的“随枣走廊”。经由崇山峻岭之间的这条狭窄山道，铜、锡等矿产资源，源源不断的北上。曾国，正是周王室这一战略生命线的安全保障。曾国国都的东南方向约200多公里处就是铜绿山，这片铜草花盛开的地方，正是东周时期规模最大的矿冶中心。已知的12个矿体中，有9个在当时便已开采。遗址中发现的矿井、炼炉、采掘和运载工具、铜锭等，展现了青铜制造业从冶炼、生产到流通的诸多环节。

长期以来，曾国都是长江与汉水之间一股强盛的力量，楚国起初完全无法与之抗衡，在随州叶家山墓地出土的斗子鼎上，郑重记载着西周初年曾国先主参与周成王会盟的经历，“丁巳，王大祓”，“己未，王赏多邦伯”，这场被曾国国君视为荣耀往事的会盟，却是楚国国君的心酸记忆。那场在岐阳举行的盛大会盟，楚国先主熊绎，也曾长途跋涉的奔赴，当他踏进仪式会场，毕恭毕敬的呈上楚地的青茅，却因实力卑弱无法正式参盟，只能与其他小国之君一起在院中看护火堆。那一夜，大殿内灯火通明，觥筹交错；冷清的庭院里，篝火熊熊，为争得列鼎而食的一席之地，南方楚人的目光里，从此留下了向北的执着。此后的楚国多代国君筚路蓝缕，方崛起于江汉西部。

春秋时期，楚国向汉水以东不断扩张，征服曾国之后，接管了铜绿山等重要矿产地，此外还掌握了丰富的锡矿资源。锡价昂贵，楚人却甚至能以锡礼器、锡箔饰随葬，获得这些重要战略资源的楚人，如虎添翼，强盛之势益发不可阻挡。数百年时移世易，楚国独霸于南方，但与姬姓的曾国联姻，对于楚国进入中原认可的秩序，仍然有着重要的意义。数代楚王都将同族的芈姓女子嫁入曾国为夫人，随州枣树林的曾国公墓中，葬有多代曾侯夫妇，191号墓中发现了夫人芈渔的器物，169号墓则出土了成套带有“随仲芈加”铭文的铜器，显示这是楚王为女儿芈加所铸的陪嫁媵器，出土的编钟上铭刻有“行相曾邦”、“余为妇为夫”等文字，记述了她曾行使国君权力的事实。

溠水汤汤，一身丧服的芈加摈退了侍从们，她是楚穆王长女曾侯宝之妻，如今丈夫早死，她与幼子面前是一个隐忧重重的曾国朝堂。芈加在楚国的支持下，决定临朝听政，从前的曾国紧紧追随周文王、周武王。随州文峰塔曾侯舆墓内出土的编钟铭文，在追忆先祖历史时，便使用了“左右文武”一语。而在附近的另一座曾侯墓出土的甬钟上，同样的语式，从属的对象却已成为了“楚王”。从“左右文武”到“左右楚王”，这一转变，正是发生在两代芈姓夫人的时代。当楚致力于对长江以南进行整合时，秦则在西部的高原绝岭间，艰难求取更大的生存空间。甘肃礼县，西汉水北岸的大堡子山上，葬着春秋早期的秦文公、静公父子。

国君大墓旁的乐器坑内，瘗埋有“秦子”镈钟、甬钟，应是文公为早死的儿子静公所铸。这套乐钟与秦公鼎、秦公簋等重器，显示出秦国对周人青铜礼器制度的忠实缵承。秦国先主原本就生活在东方，大约在西周穆王时，才迁至陇山的西边，负有为周人“保西陲”之责。由于护送周平王东迁有功，秦军得以跻身诸侯之列。继承宗周故地的秦国，被连绵山脉屏蔽于东方诸国之外，要逐鹿中原，只能通过结盟或战争的方式借道东邻的晋国。群峰寂静，大雨滂沱，刚刚自晋国凯旋的秦穆公，特地回到崤山中这处承载秦人惨痛记忆的战场。三年前，东出中原，灭滑而返的秦军在此遭遇晋军联同姜戎的伏击，数千兵士无一生还。

三年来，秦军上下无不希望一雪前恨，却屡败于晋军。直到今次，将士们渡过黄河后，将舟船尽焚，背水一战，终于大败晋军，这才敢前往崤山，为三年前阵亡的众兵士举丧。穆公于军前痛斥自己当初不听劝阻，执意发动崤之战的罪过。秦军在崤山中大哭三日，四面山谷回音阵阵，天地仿若同悲。这已是秦穆公执政的第36个年头，环顾天下，周室衰微，惟有秦与齐、楚、晋可称强国，多年来，穆公在西陲远望着东方霸主迭出，一种时不我待的急迫感，日夜催动着他。崤之战后，秦人终于认识到，东出之路已被晋国封锁，强盛之法惟有向西开拓。穆公用百里奚为左庶长，在执政的第37年，奋秦人累世经营西部之力，开地千里，称雄西戎，周天子派召公贺以金鼓。

自此，秦穆公亦被奉为春秋一霸。但崤之战也宣告着，秦晋之好已成过眼云烟，秦国东出之路暂时受阻，此后许多年，中原的政治秩序将与秦国无缘，秦国在一个相对静息的环境里，蓄力等待逐鹿中原的时机。这样的闭塞，也使得秦国的文化发展略显缓滞，列国早已抛弃的一些旧俗仍在这里延续。陕西凤翔发现的秦公一号大墓，墓主人为春秋晚期的秦景公，墓葬全长300米，是已知规模最大的先秦墓葬，墓内共发现殉葬186人，多为臣属、妃嫔以及奴婢和侍者。殉人的旧俗，铺张的墓葬，是秦国国君煊赫权力的彰示，却迥异于以周礼为纲的东方诸国。秦要参与到中原的政治格局里，还需要一场深彻的变革。

此时，以黄河岸边的函谷关为界，关城以西的秦国，蛰伏于变革的前夜之时，关城以东的局势，也正涌动起变革的浪潮，长江下游的吴国和越国，开始加入到逐鹿中原的行列。吴越壮大的基础与楚类似，也是得益于控辖的金属资源。除此外，越国更是掌握制作坚兵利器的技术。这把越王勾践剑代表了春秋时期短兵器制造的顶尖水平，该剑主要由锡青铜铸成，剑之正面，有“越王鸠浅自作用剑”八个鸟篆铭文，剑格镶有蓝色琉璃和绿松石，剑身的两面满饰黑色菱形花纹，刃薄而锋利。随着楚的强大，吴越的崛起，春秋后期的争霸主战场，逐渐从黄河流域转到长江流域，黄河流域的传统强国们则忙于应对旧有秩序瓦解的危机。其中对时局影响最深的是晋国之衰。

晋国的危机缘于公室衰微，大权旁落，六卿并立，诸卿之中，又以执政的赵氏最强。太原金胜村251号春秋晋墓，墓主赵鞅虽为正卿，却以七鼎级别随葬，俨然比肩诸侯国君。而彼时，即便是赵氏内部也人心浮动、纷争不断，卿大夫之间为维持内部团结，合力打击政敌，不得不反复的举行盟誓仪式。晋国宗庙的坛场中，司盟者手起刀落，将侍者抬来的羊杀死，殷红的牲血落入玉敦内，侍从捧着玉敦上前，参盟者逐一送出已完成书写的盟书，并以指蘸血抹于口旁。其后，侍者将杀死的祭牲埋入坎内，盟书亦被郑重地放置进去。考古工作者在盟誓遗址的400余个坎内，共发现了5000多件盟书，多以朱书写于玉圭或石圭片上，中心内容是，每个与盟人都要效忠盟主，诛讨敌对势力，并不准其重返晋地。

数量如此之多的盟书，正体现了晋国内部存在的裂隙，山盟海誓无法改变历史的走向。公元前453年，晋国被韩、赵、魏三家大夫瓜分，春秋之世历时290余年，单在《春秋》一书中所载的军事行动便有483次。春秋之初，九州天下析分为140多国，至战国之初，仅剩十余国。公室王侯、生民百姓，无不渴望结束这样的乱局。数百年乱局的一缕天光，来自一项新技术的普及。铁器，早在两周之际，炼铁技术便日趋成熟。春秋时期，秦国出土的块炼铁器数量，已远远超过其他各国。宝鸡益门村2号秦墓中，随葬有20多把异常精美的金柄剑，经鉴定，剑身的部分均为块炼铁，也就是俗称的“熟铁”。到战国时，铁器技术更是不断革新，渐成风行。相比青铜器，铁器技术具有成本低、产量大的优势，能够迅速惠及社会的各个阶层。

从考古发现的东周铁器来看，与民生密切相关的农具和手工工具最多；其次才是用于军事活动的兵器。走向普及化的铁农具和工具，极大地推动了生产力的变革，提高了人类对于自然的利用和改造能力。这样的利器，也成就了都江堰、郑国渠等各项关乎民生社稷的工程，战国时期因而迎来了社会经济的飞跃。国力与战力，固然是争夺天下的利器，但思想与文化才是深潜在世人血脉里的精髓与根本，这便是周人缔造的家国观念。现藏于中国国家博物馆的秦公簋上，刻有追述族史的铭文，“丕显朕皇祖受天命”，“宅禹蹟，十有二公在帝之坯”，自述其先祖原本生活在大禹敷布之地，凤翔秦公大墓出土石磬的铭文上，则明确提到“天子宴喜，共桓是嗣”，“高阳有灵，四方以鼏”，南方的楚人也像秦人一样，将黄帝之孙高阳追认为始祖，东方的齐国则追溯先祖为黄帝。

古史系统日益增繁，最终形成天下四方共同认可的谱系与传说，尽管干戈仍未止息，四海却已视同一家。在西方，秦与西戎之间时有摩擦，却也相互交融。位于甘肃张家川的马家塬战国墓地，正是秦文化与戎文化共生并存的写照。特殊的偏洞室墓，和随葬的陶器组合是当地戎人的传统。戎人墓主穿戴奢华，颈部戴金银半环项圈，或由各种材质珠子组成的项链，手臂上套有镶嵌繁复的金臂钏；头部饰以圆形金片和绿松石、肉红石髓珠子相间穿缀而成的发网，腰带以金银铜锡质地牌饰和珠饰组成，有的还在两侧悬挂复杂的装饰，用各类金管、肉红石髓珠、费昂斯、蜻蜓眼等串联起以高浮雕金饰或细金镶嵌玛瑙釉砂工艺的圆牌，显得斑斓绚丽。

最引人注目的，是贵族墓中随葬的大量车马，那完全按照中原样式打造的涂漆木车，车体表面，却装饰大量草原风格的金银铜铁饰件，雄鹿、猛虎和野山羊，更是属于整个欧亚草原游牧人群共有的动物形象，丰富的装饰元素充满了异族文化的张扬之美，规整连续的排列组合，又于重复中，体现华夏文化的律动与秩序。墓中随葬的铜器则体现出了秦对戎人的控制和影响力，秦与戎，这对曾经战场上的敌手，如今成了为彼此打开东西交流通途的向导。在北方，太行山东麓的中山国，由鲜虞人缔建，鲜虞，是白狄的一支，春秋时，仍被视为与中原有别的族群。在与邻近的三晋之地，发生频繁的交流、互动和摩擦之后，至战国时已渐习华风。中山王厝之冢，至今仍然高达9米，原来应是回廊环绕的三层楼阁式覆瓦建筑，顶部是享堂。

墓内随葬的青铜错金“兆域图”，绘制了陵园格局，这是迄今发现的最早的建筑设计图。该墓还出土了包括15件鼎在内的各类青铜礼器，其中一件铜方壶上錾刻着鎏金异彩的文字，言道“夫古之圣王，务在得贤，其即得民”。从墓葬建筑、随葬礼器到器物铭文，处处可见到对于周人礼制规范的遵循。而虎噬鹿器座、四龙四凤方案等青铜器，则在虎、鹿等北方式纹饰主题中，巧妙融入中原传统的龙、凤等形象，艺术与观念上的兼容并蓄，带来了生动无比的表现力。广袤大地上，处处皆是这样文化认同与文化碰撞交相辉映的动人景观。在南方，越灭吴后，成为长江下游的强国，此时，亦汇入了华夏秩序的大潮。无锡鸿山邱承墩的墓主，是地位仅次于越王的大夫，墓内随葬大量原始青瓷器，其中包括成套的编钟和编磬。

以越地的制瓷工艺，模仿青铜礼乐器型，不但象征财富与权力地位，也是越文化接触并融会华夏礼制的写照。在多元文化激荡的民族大融合时期，华夏一边坚守着文化本位，一边凝聚和吸收各族群文化中的精髓，以兼容并蓄的姿态，成为天下视野的轴心。这天下，便如一个无远弗届的同心圆，山河湖海间，处处闪耀着碰撞与交融带来的勃勃生机。

“中国”的观念，遂在这连年烽火里逐渐扩容，正如这金石乐悬是华夏礼制中最高端的“正声”，各国竞相效仿。随着时间的推移，四方与中原渐为一体，更为边远的族群，也逐渐进入时人的视野范围，“天下远近，小大若一”的理想将伴随着民族融合、国家统一的进程逐渐成为现实。

春秋战国时期，生产力的快速发展，促进了商业繁荣与城市兴盛。齐都临淄城内，曾经轻歌曼舞，百贾千工，一派繁华景象，而诸侯的争霸，又带来了对人才的广泛需求。临淄西门之外，是稷下学宫之所在，这里一度是战国时期的学术中心，曾容纳了当时诸子百家的各派学者，多达千人左右，这些学者纵横于各国之间，打破了地域的藩篱，交响出文化融会的新声。于是，新思想和新秩序的建设也随之而来。彼时的诸侯，无不反复追问，天下如何才能安定？孟子的回答是三个字：定于一。春秋战国扰攘五六百年，在中国历史上，常被当作乱世。然而，正是在这样一个时代里，中国经历了前所未有的碰撞与融合，孕育出天下大一统的基础。

进入战国时代以后，最有可能实践统一这幅蓝图的，是所谓的“七雄”，也就是函谷关以东的齐、楚、魏、燕、韩、赵六国加一个函谷关以西的秦国。此时的秦早已摈弃殉人等旧俗，但要与东边的强国抗衡，还需要更为积极的变革。秦孝公即位后，听取了魏国人商鞅的建议，在宫室盛大的栎阳城掀起了东周五百年间最风云激荡的变法。这次变法的核心，是以重典维护国家安全，进行战备积累，又以奖励军功提供上升途径，激活国力。公元前350年，秦人从栎阳迁都咸阳，商鞅在此进行了第二次变法。商鞅方升，是商鞅为统一度量衡所监制的标准量器，容积为202.15毫升，是当时的一升。这件体量不大的青铜器，却成为古中国经济秩序奠基之途上的“重器”。

商鞅两次变法，废井田，开阡陌，重农抑商，奖励军功，严革法度，推行郡县，编订户口，是一场覆盖政治、军事、经济、社会和文化制度的全方位改革，广泛调动了各种积极因素，推动秦国社会进步、经济发展、国力壮大、全面强盛。同时代的战国诸雄，虽时有一些改革作为，但却远不及秦深彻变法的决心和格局。在这风云激荡的时代里，崭新的秦国立于关中，眼望向广袤的万里河山，即将开启结束乱世，天下一统的征程。公元前328年，秦取上郡，拓地至陕北；公元前316年，秦灭巴蜀，囊括四川盆地；公元前272年，秦灭义渠，攻克了黄土高原最大的戎狄势力。秦国就这样，以关中盆地为营，分别向北、西、南三个方向扩展战略纵深，随后剑锋东指，匡合天下。

在商鞅变法的百年之后，秦的统一大业已势如破竹。湖北云梦，本为楚地，被秦人占据后，成为攻灭楚国的战略要地。埋葬在睡虎地的喜，是秦军中的一位下级军吏，亲身经历了秦完成统一的最后阶段。随葬的《编年记》载录着他的这段履历，嬴政在位的第十三年，喜从军，十五年从平阳军，十七年，攻韩，十八年，攻赵，十九年，南郡备警，廿一年，韩王死，廿二年，攻魏梁，廿三年，攻荆，楚昌文君死。甘肃礼县，海拔1867米的四角子山之巅，以中央方形土台，及其上半地穴房址为中心，四周对称分布着平行的双排建筑。这是秦始皇统一六国后，在秦人祖居之地“犬丘”修建的祭天场所。公元前221年，秦始皇熔化天下兵器，铸十二铜人，颁发诏版，宣告天下“廿六年，皇帝尽并兼天下诸侯，黔首大安”。三晋的强兵劲弩，楚的坚守抵抗，齐的合纵连横，燕的舍生忘死，都是向着天下一统驱驰的时代侧影。

秦，终结了数百年的纷争，延续和实现了周人的家国理想，将华夏九州推入统一多民族国家的时代。站在江河入海的时代，回顾奔涌的历史长流，是蹈海踏浪而来的文明记忆。曾有人在黑夜里，擎起天地间的第一簇火种，曾有人在漫天星斗下，隔着山川湖海遥望。人聚成邑，邑聚为国，曾有古国在大江大河间初铸辉煌，自九州之中，酿就四方众望、人心所向的风潮。经夏商周而立秦汉，生长在世界东方的古代中国文明，跋涉数千年岁月而来，穿梭在生死枯荣的古文明之林间，历风尘霜雪而初心未改。蹈蹚出一条不同于世界其他文明体的发展道路，创造出灿烂盛大的中华文明。《何以中国》从始至终，亦是站在考古学搭起的阶梯上，向千百代来成就今日之中国的祖先们致敬！自先人起，便未更改对于这片土地的深情，以及对于和平与统一的刻骨记忆，让我们在这片祖先们曾守望过的天下，携手同心，去往那光明盛放的前方！



\end{document}
